% preamble.tex -- layout, packages, notation and generated-input helpers shared
% by main.tex and MMHM_v1_erratum.tex. \input after \documentclass[10pt]{article}.

% ---------------------------------------------------------------------------
% Page layout and fonts: the preprint look of version 1 (its arxiv style),
% rebuilt here because arxiv.sty is not part of the TeX distribution.
% ---------------------------------------------------------------------------
\usepackage[letterpaper,textwidth=6.5in,textheight=9in,top=1in,
            headheight=14pt,headsep=25pt,footskip=30pt]{geometry}
\usepackage[T1]{fontenc}
\usepackage[utf8]{inputenc}
\renewcommand{\rmdefault}{ptm}            % Times for text; math stays Computer Modern
\usepackage[scaled=0.9]{helvet}           % sans serif (method names)
\usepackage[protrusion=true,expansion=true]{microtype}

\makeatletter
% Font sizes with the tight leading of version 1 (10/11pt body).
\renewcommand{\normalsize}{%
  \@setfontsize\normalsize\@xpt\@xipt
  \abovedisplayskip 7\p@ \@plus 2\p@ \@minus 5\p@
  \abovedisplayshortskip \z@ \@plus 3\p@
  \belowdisplayskip \abovedisplayskip
  \belowdisplayshortskip 4\p@ \@plus 3\p@ \@minus 3\p@}
\normalsize
\renewcommand{\small}{\@setfontsize\small\@ixpt\@xpt}
\renewcommand{\footnotesize}{\@setfontsize\footnotesize\@ixpt\@xpt}
\renewcommand{\scriptsize}{\@setfontsize\scriptsize\@viipt\@viiipt}
\renewcommand{\tiny}{\@setfontsize\tiny\@vipt\@viipt}
\renewcommand{\large}{\@setfontsize\large\@xiipt{14}}
\renewcommand{\Large}{\@setfontsize\Large\@xivpt{16}}
\renewcommand{\LARGE}{\@setfontsize\LARGE\@xviipt{20}}

% Paragraphs: no indentation, space between paragraphs.
\setlength{\parindent}{\z@}
\setlength{\parskip}{5.5\p@}
\widowpenalty=10000
\clubpenalty=10000

% Section headings.
\renewcommand{\section}{\@startsection{section}{1}{\z@}%
  {-2.0ex \@plus -0.5ex \@minus -0.2ex}{1.5ex \@plus 0.3ex \@minus 0.2ex}%
  {\large\bfseries\raggedright}}
\renewcommand{\subsection}{\@startsection{subsection}{2}{\z@}%
  {-1.8ex \@plus -0.5ex \@minus -0.2ex}{0.8ex \@plus 0.2ex}%
  {\normalsize\bfseries\raggedright}}
\renewcommand{\subsubsection}{\@startsection{subsubsection}{3}{\z@}%
  {-1.5ex \@plus -0.5ex \@minus -0.2ex}{0.5ex \@plus 0.2ex}%
  {\normalsize\bfseries\raggedright}}
\renewcommand{\paragraph}{\@startsection{paragraph}{4}{\z@}%
  {1.5ex \@plus 0.5ex \@minus 0.2ex}{-1em}{\normalsize\bfseries}}

% Floats.
\renewcommand{\topfraction}{0.85}
\renewcommand{\bottomfraction}{0.4}
\renewcommand{\textfraction}{0.1}
\renewcommand{\floatpagefraction}{0.7}
\setlength{\textfloatsep}{12\p@ \@plus 2\p@ \@minus 2\p@}
\setlength{\floatsep}{12\p@ \@plus 2\p@ \@minus 2\p@}
\setlength{\@fptop}{\z@}            % float-only pages start at the top
\renewcommand{\footnoterule}{\kern-3\p@ \hrule width 12pc \kern 2.6\p@}

% Title block: rules above and below a small-caps title, "A Preprint" under it,
% centred author tabular whose first row (the name) is bold.
\newcommand{\shorttitle}{Morse-based Modular Homology for Evolving Simplicial Complexes}
\renewcommand{\@maketitle}{%
  \vbox{%
    \hsize\textwidth
    \linewidth\hsize
    \vskip 0.1in
    \hrule height 2\p@
    \vskip 0.25in
    \vskip -\parskip
    \centering
    {\LARGE\scshape \@title\par}
    \vskip 0.29in
    \vskip -\parskip
    \hrule height 2\p@
    \vskip 0.09in
    {\scshape A Preprint}\par
    \vskip 0.1in
    \begin{tabular}[t]{c}\bfseries\rule{\z@}{24\p@}\ignorespaces\@author\end{tabular}\par
    \vskip 0.4in \@minus 0.1in
    {\@date\par}
    \vskip 0.2in
  }%
}
\renewcommand{\maketitle}{\par\begingroup\@maketitle\endgroup
  \thispagestyle{empty}\setcounter{footnote}{0}}

% Abstract and keywords as in version 1.
\renewenvironment{abstract}%
  {\centerline{\large\bfseries\scshape Abstract}\begin{quote}}%
  {\end{quote}}
\newcommand{\keywords}[1]{\par\addvspace\medskipamount
  {\rightskip=\z@ \@plus 1cm\noindent{\bfseries\itshape Keywords}\enspace
   \ignorespaces#1\par}}
\makeatother

% Running head: title centred, "A Preprint" right, rule below.
\usepackage{fancyhdr}
\fancyhf{}
\pagestyle{fancy}
\renewcommand{\headrulewidth}{0.4pt}
\fancyheadoffset{0pt}
\fancyhead[C]{\shorttitle}
\fancyhead[R]{\scshape A Preprint}
\fancyfoot[C]{\thepage}

% ---------------------------------------------------------------------------
% Packages
% ---------------------------------------------------------------------------
\usepackage{amsmath,amssymb,amsthm}
\usepackage{graphicx}
\graphicspath{{../figures/}{figures/}}
\usepackage{booktabs}
\usepackage{array}
\usepackage{tabularx}
\usepackage[table]{xcolor}
\usepackage{tikz}
\usetikzlibrary{arrows.meta,positioning,calc,fit,backgrounds,decorations.pathreplacing}
\usepackage[font=small,labelfont=bf,skip=6pt]{caption}
\usepackage{enumitem}
\setlist{topsep=2pt,itemsep=1pt,parsep=1pt}
\usepackage[numbers,square,sort&compress]{natbib}
\usepackage{url}
\usepackage[hyperfootnotes=false]{hyperref}
% Link colours echo the link boxes of version 1 (red: internal, green:
% citations, cyan: URLs), darkened so that they print well in grey.
\definecolor{linkred}{rgb}{0.55,0.05,0.05}
\definecolor{citegreen}{rgb}{0.05,0.42,0.10}
\definecolor{urlcyan}{rgb}{0.0,0.30,0.50}
\hypersetup{
  colorlinks=true,
  linkcolor=linkred,
  citecolor=citegreen,
  urlcolor=urlcyan,
  pdftitle={Morse-based Modular Homology for Evolving Simplicial Complexes},
  pdfauthor={Anqiao Ouyang},
  pdfkeywords={simplicial homology, Betti numbers, dynamic algorithms,
    discrete Morse theory, matrix reduction},
}
\usepackage{doi}
\usepackage[capitalise,noabbrev,nameinlink]{cleveref}

% Shared notation and theorem environments (after cleveref, so that theorem,
% lemma and proposition, which share a counter, keep their own names).
% macros.tex -- shared notation and theorem environments for the MMHM v2 paper.
% \input this file in the preamble of main.tex (after \documentclass).
% Requires amsmath, amssymb, amsthm (loaded here; loading them again is harmless).
\usepackage{amsmath,amssymb,amsthm}

% ---------------------------------------------------------------------------
% Theorem-like environments.
% Theorem, lemma and proposition share one counter so that the main results are
% Theorem 1, Lemma 2, Proposition 3, Proposition 4 (as cited in SPEC.md);
% corollaries are numbered after their parent result (Corollary 2.1, 2.2).
% Observations have their own counter; definitions, examples and remarks share
% a further counter.
% ---------------------------------------------------------------------------
\makeatletter
\theoremstyle{plain}
\@ifundefined{theorem}{\newtheorem{theorem}{Theorem}}{}
\@ifundefined{lemma}{\newtheorem{lemma}[theorem]{Lemma}}{}
\@ifundefined{proposition}{\newtheorem{proposition}[theorem]{Proposition}}{}
\@ifundefined{conjecture}{\newtheorem{conjecture}[theorem]{Conjecture}}{}
\@ifundefined{corollary}{\newtheorem{corollary}{Corollary}[theorem]}{}
\@ifundefined{claim}{\newtheorem*{claim}{Claim}}{}
\@ifundefined{observation}{\newtheorem{observation}{Observation}}{}
\theoremstyle{definition}
\@ifundefined{definition}{\newtheorem{definition}{Definition}}{}
\@ifundefined{example}{\newtheorem{example}[definition]{Example}}{}
\theoremstyle{remark}
\@ifundefined{remark}{\newtheorem{remark}[definition]{Remark}}{}
\makeatother

% ---------------------------------------------------------------------------
% Notation (SPEC.md sections 1--4)
% ---------------------------------------------------------------------------
\providecommand{\Z}{\mathbb{Z}_2}                 % the coefficient field
\providecommand{\N}{\mathbb{N}}
\providecommand{\Kc}[1]{K_{#1}}                   % k-simplices of K
\providecommand{\Cc}[1]{C_{#1}}                   % chain group C_k(K; Z2)
\providecommand{\Bm}[1]{B_{#1}}                   % boundary matrix of d_k
\providecommand{\Zc}[1]{Z_{#1}}                   % cycle space ker d_k
\providecommand{\bd}{\partial}
\providecommand{\dM}{\partial^{M}}                % Morse boundary
\providecommand{\ord}{\operatorname{ord}}
\providecommand{\low}{\operatorname{low}}
\providecommand{\occ}{\operatorname{occ}}
\providecommand{\pivot}{\operatorname{pivot}}
\providecommand{\rank}{\operatorname{rank}}
\providecommand{\nnz}{\operatorname{nnz}}
\providecommand{\supp}{\operatorname{supp}}
\providecommand{\rmt}{\operatorname{rm}}           % removal time in the collapse
\providecommand{\Crit}{\mathrm{Crit}}             % critical cells
\providecommand{\Lo}{\mathrm{Low}}                % lower cells (matched with a coface)
\providecommand{\Up}{\mathrm{Up}}                 % upper cells (matched with a facet)
\providecommand{\flow}{\Pi}                       % gradient-flow projection
\providecommand{\symdiff}{\mathbin{\triangle}}     % symmetric difference = Z2 sum
\providecommand{\facets}{\operatorname{fac}}
\providecommand{\method}[1]{\textsf{#1}}          % method names: \method{mmhm}
\providecommand{\Ins}{\textup{\textsc{Insert}}}
\providecommand{\Del}{\textup{\textsc{Delete}}}
\providecommand{\Collapse}{\textup{\textsc{Collapse}}}
\providecommand{\Ot}{\widetilde{O}}

\crefname{observation}{Observation}{Observations}
\crefname{conjecture}{Conjecture}{Conjectures}
\crefname{claim}{Claim}{Claims}
\crefname{remark}{Remark}{Remarks}
\crefname{example}{Example}{Examples}
\crefname{definition}{Definition}{Definitions}
\crefname{corollary}{Corollary}{Corollaries}

% ---------------------------------------------------------------------------
% Generated inputs. A missing file renders as a visible box and leaves a
% "Package mmhm Warning" in the log (grep 'Package mmhm Warning' main.log).
% ---------------------------------------------------------------------------
\makeatletter
% a missing figure file
\newcommand{\figpending}[2]{% #1 file stem, #2 height
  \PackageWarning{mmhm}{Figure file ../figures/#1.pdf not found}%
  \fcolorbox{black!55}{black!5}{\parbox[c][#2][c]{\dimexpr\linewidth-2\fboxsep-2\fboxrule\relax}{%
    \centering\small\sffamily Figure missing:
    \texttt{figures/\detokenize{#1}.pdf}}}}
% include a result figure if it exists
\newcommand{\resultfigure}[3][width=\linewidth]{% [graphicx options] {file stem} {placeholder height}
  \IfFileExists{../figures/#2.pdf}%
    {\includegraphics[#1]{#2.pdf}}%
    {\figpending{#2}{#3}}}
% include a generated table body (a complete tabular) if it exists
\newcommand{\generatedtable}[1]{% #1 file stem in generated/
  \IfFileExists{generated/#1.tex}%
    {\input{generated/#1}}%
    {\PackageWarning{mmhm}{Generated table generated/#1.tex not found}%
     \fcolorbox{black!55}{black!5}{\parbox{0.8\linewidth}{\centering\small\sffamily
       Table missing: \texttt{generated/\detokenize{#1}.tex}}}}}

% Measured numbers: generated/numbers.tex defines \Num... macros. Every macro
% used in the text gets a fallback that prints a bold "??" and warns.
% Contract: the text uses every \Num macro in text mode, never inside $...$,
% so a generated value may be plain text (2000, 1.2{,}345) or contain math.
\IfFileExists{generated/numbers.tex}{% Generated by experiments/aggregate.py -- do not edit.
% Each macro: \newcommand{\NumX}{value} preceded by a comment saying
% what it is and which config/script produced it. See README.md.
% CPU-hours of all runs (wall time of each single-threaded process, summed) [run.py manifests]
\newcommand{\NumComputeHours}{5.2}
% E8 w3_N20_k2 drb-id (3 seeds pooled, non-rebuild updates): median additions per update [aggregate.py]
\newcommand{\NumEEightWThreeNTwentyDrbIdAddsMed}{5{,}950}
% E8 w3_N20_k2 drb-id (3 seeds pooled, non-rebuild updates): p99 additions per update [aggregate.py]
\newcommand{\NumEEightWThreeNTwentyDrbIdAddsPNinetyNine}{34{,}000}
% E8 w3_N20_k2 drb-id (3 seeds pooled, non-rebuild updates): median footprint_sum per update [aggregate.py]
\newcommand{\NumEEightWThreeNTwentyDrbIdFootprintMed}{607}
% E8 w3_N20_k2 drb-id (3 seeds pooled, non-rebuild updates): p99 footprint_sum [aggregate.py]
\newcommand{\NumEEightWThreeNTwentyDrbIdFootprintPNinetyNine}{34{,}700}
% E8 w3_N20_k2 drb-id-rb (3 seeds pooled, non-rebuild updates): median additions per update [aggregate.py]
\newcommand{\NumEEightWThreeNTwentyDrbIdRbAddsMed}{4{,}960}
% E8 w3_N20_k2 drb-id-rb (3 seeds pooled, non-rebuild updates): p99 additions per update [aggregate.py]
\newcommand{\NumEEightWThreeNTwentyDrbIdRbAddsPNinetyNine}{19{,}100}
% E8 w3_N20_k2 drb-id-rb (3 seeds pooled, non-rebuild updates): median footprint_sum per update [aggregate.py]
\newcommand{\NumEEightWThreeNTwentyDrbIdRbFootprintMed}{671}
% E8 w3_N20_k2 drb-id-rb (3 seeds pooled, non-rebuild updates): p99 footprint_sum [aggregate.py]
\newcommand{\NumEEightWThreeNTwentyDrbIdRbFootprintPNinetyNine}{13{,}300}
% E8 w3_N20_k2 drb-id-rb (3 seeds pooled, non-rebuild updates): % of updates with zero additions [aggregate.py]
\newcommand{\NumEEightWThreeNTwentyDrbIdRbZeroAddsPct}{0.3}
% E8 w3_N20_k2 drb-id (3 seeds pooled, non-rebuild updates): % of updates with zero additions [aggregate.py]
\newcommand{\NumEEightWThreeNTwentyDrbIdZeroAddsPct}{0.3}
% E8 w3_N20_k2 mmhm (3 seeds pooled, non-rebuild updates): median additions per update [aggregate.py]
\newcommand{\NumEEightWThreeNTwentyMmhmAddsMed}{4{,}170}
% E8 w3_N20_k2 mmhm (3 seeds pooled, non-rebuild updates): p99 additions per update [aggregate.py]
\newcommand{\NumEEightWThreeNTwentyMmhmAddsPNinetyNine}{21{,}300}
% E8 w3_N20_k2 mmhm (3 seeds pooled, non-rebuild updates): median footprint_sum per update [aggregate.py]
\newcommand{\NumEEightWThreeNTwentyMmhmFootprintMed}{594}
% E8 w3_N20_k2 mmhm (3 seeds pooled, non-rebuild updates): p99 footprint_sum [aggregate.py]
\newcommand{\NumEEightWThreeNTwentyMmhmFootprintPNinetyNine}{17{,}000}
% E8 w3_N20_k2 mmhm-norb (3 seeds pooled, non-rebuild updates): median additions per update [aggregate.py]
\newcommand{\NumEEightWThreeNTwentyMmhmNorbAddsMed}{4{,}570}
% E8 w3_N20_k2 mmhm-norb (3 seeds pooled, non-rebuild updates): p99 additions per update [aggregate.py]
\newcommand{\NumEEightWThreeNTwentyMmhmNorbAddsPNinetyNine}{31{,}200}
% E8 w3_N20_k2 mmhm-norb (3 seeds pooled, non-rebuild updates): median footprint_sum per update [aggregate.py]
\newcommand{\NumEEightWThreeNTwentyMmhmNorbFootprintMed}{580}
% E8 w3_N20_k2 mmhm-norb (3 seeds pooled, non-rebuild updates): p99 footprint_sum [aggregate.py]
\newcommand{\NumEEightWThreeNTwentyMmhmNorbFootprintPNinetyNine}{31{,}900}
% E8 w3_N20_k2 mmhm-norb (3 seeds pooled, non-rebuild updates): % of updates with zero additions [aggregate.py]
\newcommand{\NumEEightWThreeNTwentyMmhmNorbZeroAddsPct}{0.3}
% E8 w3_N20_k2 mmhm (3 seeds pooled, non-rebuild updates): % of updates with zero additions [aggregate.py]
\newcommand{\NumEEightWThreeNTwentyMmhmZeroAddsPct}{0.3}
% E7 w3_N20_k2 drb-id-rb alpha=4 period=off: mean additions per non-rebuild update [configs/E7.json, aggregate.py]
\newcommand{\NumESevenWThreeNTwentyDrbIdRbAFourPOffAdds}{7{,}180}
% E7 w3_N20_k2 drb-id-rb alpha=4 period=off: amortized latency (us) [configs/E7.json, aggregate.py]
\newcommand{\NumESevenWThreeNTwentyDrbIdRbAFourPOffAmort}{4{,}390}
% E7 w3_N20_k2 drb-id-rb alpha=4 period=off: median latency (us) [configs/E7.json, aggregate.py]
\newcommand{\NumESevenWThreeNTwentyDrbIdRbAFourPOffMed}{2{,}630}
% E7 w3_N20_k2 drb-id-rb alpha=4 period=off: p99 latency (us) [configs/E7.json, aggregate.py]
\newcommand{\NumESevenWThreeNTwentyDrbIdRbAFourPOffPNinetyNine}{26{,}500}
% E7 w3_N20_k2 drb-id-rb alpha=4 period=off: rebuilds per run (mean of seeds) [configs/E7.json, aggregate.py]
\newcommand{\NumESevenWThreeNTwentyDrbIdRbAFourPOffRebuilds}{1.0}
% E7 w3_N20_k2 drb-id-rb alpha=4 period=64: mean additions per non-rebuild update [configs/E7.json, aggregate.py]
\newcommand{\NumESevenWThreeNTwentyDrbIdRbAFourPSixtyFourAdds}{5{,}380}
% E7 w3_N20_k2 drb-id-rb alpha=4 period=64: amortized latency (us) [configs/E7.json, aggregate.py]
\newcommand{\NumESevenWThreeNTwentyDrbIdRbAFourPSixtyFourAmort}{2{,}130}
% E7 w3_N20_k2 drb-id-rb alpha=4 period=64: median latency (us) [configs/E7.json, aggregate.py]
\newcommand{\NumESevenWThreeNTwentyDrbIdRbAFourPSixtyFourMed}{1{,}330}
% E7 w3_N20_k2 drb-id-rb alpha=4 period=64: p99 latency (us) [configs/E7.json, aggregate.py]
\newcommand{\NumESevenWThreeNTwentyDrbIdRbAFourPSixtyFourPNinetyNine}{31{,}100}
% E7 w3_N20_k2 drb-id-rb alpha=4 period=64: rebuilds per run (mean of seeds) [configs/E7.json, aggregate.py]
\newcommand{\NumESevenWThreeNTwentyDrbIdRbAFourPSixtyFourRebuilds}{31}
% E7 w3_N20_k2 drb-id-rb alpha=4 period=256: mean additions per non-rebuild update [configs/E7.json, aggregate.py]
\newcommand{\NumESevenWThreeNTwentyDrbIdRbAFourPTwoHundredFiftySixAdds}{6{,}210}
% E7 w3_N20_k2 drb-id-rb alpha=4 period=256: amortized latency (us) [configs/E7.json, aggregate.py]
\newcommand{\NumESevenWThreeNTwentyDrbIdRbAFourPTwoHundredFiftySixAmort}{2{,}670}
% E7 w3_N20_k2 drb-id-rb alpha=4 period=256: median latency (us) [configs/E7.json, aggregate.py]
\newcommand{\NumESevenWThreeNTwentyDrbIdRbAFourPTwoHundredFiftySixMed}{1{,}680}
% E7 w3_N20_k2 drb-id-rb alpha=4 period=256: p99 latency (us) [configs/E7.json, aggregate.py]
\newcommand{\NumESevenWThreeNTwentyDrbIdRbAFourPTwoHundredFiftySixPNinetyNine}{14{,}700}
% E7 w3_N20_k2 drb-id-rb alpha=4 period=256: rebuilds per run (mean of seeds) [configs/E7.json, aggregate.py]
\newcommand{\NumESevenWThreeNTwentyDrbIdRbAFourPTwoHundredFiftySixRebuilds}{7.0}
% E7 w3_N20_k2 drb-id-rb alpha=0.5 period=off: mean additions per non-rebuild update [configs/E7.json, aggregate.py]
\newcommand{\NumESevenWThreeNTwentyDrbIdRbAHalfPOffAdds}{5{,}380}
% E7 w3_N20_k2 drb-id-rb alpha=0.5 period=off: amortized latency (us) [configs/E7.json, aggregate.py]
\newcommand{\NumESevenWThreeNTwentyDrbIdRbAHalfPOffAmort}{2{,}090}
% E7 w3_N20_k2 drb-id-rb alpha=0.5 period=off: median latency (us) [configs/E7.json, aggregate.py]
\newcommand{\NumESevenWThreeNTwentyDrbIdRbAHalfPOffMed}{1{,}340}
% E7 w3_N20_k2 drb-id-rb alpha=0.5 period=off: p99 latency (us) [configs/E7.json, aggregate.py]
\newcommand{\NumESevenWThreeNTwentyDrbIdRbAHalfPOffPNinetyNine}{23{,}300}
% E7 w3_N20_k2 drb-id-rb alpha=0.5 period=off: rebuilds per run (mean of seeds) [configs/E7.json, aggregate.py]
\newcommand{\NumESevenWThreeNTwentyDrbIdRbAHalfPOffRebuilds}{22}
% E7 w3_N20_k2 drb-id-rb alpha=0.5 period=64: mean additions per non-rebuild update [configs/E7.json, aggregate.py]
\newcommand{\NumESevenWThreeNTwentyDrbIdRbAHalfPSixtyFourAdds}{5{,}350}
% E7 w3_N20_k2 drb-id-rb alpha=0.5 period=64: amortized latency (us) [configs/E7.json, aggregate.py]
\newcommand{\NumESevenWThreeNTwentyDrbIdRbAHalfPSixtyFourAmort}{2{,}120}
% E7 w3_N20_k2 drb-id-rb alpha=0.5 period=64: median latency (us) [configs/E7.json, aggregate.py]
\newcommand{\NumESevenWThreeNTwentyDrbIdRbAHalfPSixtyFourMed}{1{,}290}
% E7 w3_N20_k2 drb-id-rb alpha=0.5 period=64: p99 latency (us) [configs/E7.json, aggregate.py]
\newcommand{\NumESevenWThreeNTwentyDrbIdRbAHalfPSixtyFourPNinetyNine}{31{,}900}
% E7 w3_N20_k2 drb-id-rb alpha=0.5 period=64: rebuilds per run (mean of seeds) [configs/E7.json, aggregate.py]
\newcommand{\NumESevenWThreeNTwentyDrbIdRbAHalfPSixtyFourRebuilds}{31}
% E7 w3_N20_k2 drb-id-rb alpha=0.5 period=256: mean additions per non-rebuild update [configs/E7.json, aggregate.py]
\newcommand{\NumESevenWThreeNTwentyDrbIdRbAHalfPTwoHundredFiftySixAdds}{5{,}380}
% E7 w3_N20_k2 drb-id-rb alpha=0.5 period=256: amortized latency (us) [configs/E7.json, aggregate.py]
\newcommand{\NumESevenWThreeNTwentyDrbIdRbAHalfPTwoHundredFiftySixAmort}{2{,}060}
% E7 w3_N20_k2 drb-id-rb alpha=0.5 period=256: median latency (us) [configs/E7.json, aggregate.py]
\newcommand{\NumESevenWThreeNTwentyDrbIdRbAHalfPTwoHundredFiftySixMed}{1{,}310}
% E7 w3_N20_k2 drb-id-rb alpha=0.5 period=256: p99 latency (us) [configs/E7.json, aggregate.py]
\newcommand{\NumESevenWThreeNTwentyDrbIdRbAHalfPTwoHundredFiftySixPNinetyNine}{22{,}900}
% E7 w3_N20_k2 drb-id-rb alpha=0.5 period=256: rebuilds per run (mean of seeds) [configs/E7.json, aggregate.py]
\newcommand{\NumESevenWThreeNTwentyDrbIdRbAHalfPTwoHundredFiftySixRebuilds}{22}
% E7 w3_N20_k2 drb-id-rb alpha=off period=off: mean additions per non-rebuild update [configs/E7.json, aggregate.py]
\newcommand{\NumESevenWThreeNTwentyDrbIdRbAOffPOffAdds}{7{,}400}
% E7 w3_N20_k2 drb-id-rb alpha=off period=off: amortized latency (us) [configs/E7.json, aggregate.py]
\newcommand{\NumESevenWThreeNTwentyDrbIdRbAOffPOffAmort}{6{,}910}
% E7 w3_N20_k2 drb-id-rb alpha=off period=off: median latency (us) [configs/E7.json, aggregate.py]
\newcommand{\NumESevenWThreeNTwentyDrbIdRbAOffPOffMed}{3{,}520}
% E7 w3_N20_k2 drb-id-rb alpha=off period=off: p99 latency (us) [configs/E7.json, aggregate.py]
\newcommand{\NumESevenWThreeNTwentyDrbIdRbAOffPOffPNinetyNine}{49{,}900}
% E7 w3_N20_k2 drb-id-rb alpha=off period=off: rebuilds per run (mean of seeds) [configs/E7.json, aggregate.py]
\newcommand{\NumESevenWThreeNTwentyDrbIdRbAOffPOffRebuilds}{0}
% E7 w3_N20_k2 drb-id-rb alpha=off period=64: mean additions per non-rebuild update [configs/E7.json, aggregate.py]
\newcommand{\NumESevenWThreeNTwentyDrbIdRbAOffPSixtyFourAdds}{5{,}380}
% E7 w3_N20_k2 drb-id-rb alpha=off period=64: amortized latency (us) [configs/E7.json, aggregate.py]
\newcommand{\NumESevenWThreeNTwentyDrbIdRbAOffPSixtyFourAmort}{2{,}070}
% E7 w3_N20_k2 drb-id-rb alpha=off period=64: median latency (us) [configs/E7.json, aggregate.py]
\newcommand{\NumESevenWThreeNTwentyDrbIdRbAOffPSixtyFourMed}{1{,}260}
% E7 w3_N20_k2 drb-id-rb alpha=off period=64: p99 latency (us) [configs/E7.json, aggregate.py]
\newcommand{\NumESevenWThreeNTwentyDrbIdRbAOffPSixtyFourPNinetyNine}{29{,}000}
% E7 w3_N20_k2 drb-id-rb alpha=off period=64: rebuilds per run (mean of seeds) [configs/E7.json, aggregate.py]
\newcommand{\NumESevenWThreeNTwentyDrbIdRbAOffPSixtyFourRebuilds}{31}
% E7 w3_N20_k2 drb-id-rb alpha=off period=256: mean additions per non-rebuild update [configs/E7.json, aggregate.py]
\newcommand{\NumESevenWThreeNTwentyDrbIdRbAOffPTwoHundredFiftySixAdds}{6{,}210}
% E7 w3_N20_k2 drb-id-rb alpha=off period=256: amortized latency (us) [configs/E7.json, aggregate.py]
\newcommand{\NumESevenWThreeNTwentyDrbIdRbAOffPTwoHundredFiftySixAmort}{2{,}330}
% E7 w3_N20_k2 drb-id-rb alpha=off period=256: median latency (us) [configs/E7.json, aggregate.py]
\newcommand{\NumESevenWThreeNTwentyDrbIdRbAOffPTwoHundredFiftySixMed}{1{,}410}
% E7 w3_N20_k2 drb-id-rb alpha=off period=256: p99 latency (us) [configs/E7.json, aggregate.py]
\newcommand{\NumESevenWThreeNTwentyDrbIdRbAOffPTwoHundredFiftySixPNinetyNine}{12{,}800}
% E7 w3_N20_k2 drb-id-rb alpha=off period=256: rebuilds per run (mean of seeds) [configs/E7.json, aggregate.py]
\newcommand{\NumESevenWThreeNTwentyDrbIdRbAOffPTwoHundredFiftySixRebuilds}{7.0}
% E7 w3_N20_k2 drb-id-rb alpha=1 period=off: mean additions per non-rebuild update [configs/E7.json, aggregate.py]
\newcommand{\NumESevenWThreeNTwentyDrbIdRbAOnePOffAdds}{5{,}780}
% E7 w3_N20_k2 drb-id-rb alpha=1 period=off: amortized latency (us) [configs/E7.json, aggregate.py]
\newcommand{\NumESevenWThreeNTwentyDrbIdRbAOnePOffAmort}{2{,}230}
% E7 w3_N20_k2 drb-id-rb alpha=1 period=off: median latency (us) [configs/E7.json, aggregate.py]
\newcommand{\NumESevenWThreeNTwentyDrbIdRbAOnePOffMed}{1{,}490}
% E7 w3_N20_k2 drb-id-rb alpha=1 period=off: p99 latency (us) [configs/E7.json, aggregate.py]
\newcommand{\NumESevenWThreeNTwentyDrbIdRbAOnePOffPNinetyNine}{10{,}000}
% E7 w3_N20_k2 drb-id-rb alpha=1 period=off: rebuilds per run (mean of seeds) [configs/E7.json, aggregate.py]
\newcommand{\NumESevenWThreeNTwentyDrbIdRbAOnePOffRebuilds}{10}
% E7 w3_N20_k2 drb-id-rb alpha=1 period=64: mean additions per non-rebuild update [configs/E7.json, aggregate.py]
\newcommand{\NumESevenWThreeNTwentyDrbIdRbAOnePSixtyFourAdds}{5{,}350}
% E7 w3_N20_k2 drb-id-rb alpha=1 period=64: amortized latency (us) [configs/E7.json, aggregate.py]
\newcommand{\NumESevenWThreeNTwentyDrbIdRbAOnePSixtyFourAmort}{2{,}130}
% E7 w3_N20_k2 drb-id-rb alpha=1 period=64: median latency (us) [configs/E7.json, aggregate.py]
\newcommand{\NumESevenWThreeNTwentyDrbIdRbAOnePSixtyFourMed}{1{,}290}
% E7 w3_N20_k2 drb-id-rb alpha=1 period=64: p99 latency (us) [configs/E7.json, aggregate.py]
\newcommand{\NumESevenWThreeNTwentyDrbIdRbAOnePSixtyFourPNinetyNine}{32{,}100}
% E7 w3_N20_k2 drb-id-rb alpha=1 period=64: rebuilds per run (mean of seeds) [configs/E7.json, aggregate.py]
\newcommand{\NumESevenWThreeNTwentyDrbIdRbAOnePSixtyFourRebuilds}{31}
% E7 w3_N20_k2 drb-id-rb alpha=1 period=256: mean additions per non-rebuild update [configs/E7.json, aggregate.py]
\newcommand{\NumESevenWThreeNTwentyDrbIdRbAOnePTwoHundredFiftySixAdds}{5{,}780}
% E7 w3_N20_k2 drb-id-rb alpha=1 period=256: amortized latency (us) [configs/E7.json, aggregate.py]
\newcommand{\NumESevenWThreeNTwentyDrbIdRbAOnePTwoHundredFiftySixAmort}{2{,}250}
% E7 w3_N20_k2 drb-id-rb alpha=1 period=256: median latency (us) [configs/E7.json, aggregate.py]
\newcommand{\NumESevenWThreeNTwentyDrbIdRbAOnePTwoHundredFiftySixMed}{1{,}490}
% E7 w3_N20_k2 drb-id-rb alpha=1 period=256: p99 latency (us) [configs/E7.json, aggregate.py]
\newcommand{\NumESevenWThreeNTwentyDrbIdRbAOnePTwoHundredFiftySixPNinetyNine}{9{,}470}
% E7 w3_N20_k2 drb-id-rb alpha=1 period=256: rebuilds per run (mean of seeds) [configs/E7.json, aggregate.py]
\newcommand{\NumESevenWThreeNTwentyDrbIdRbAOnePTwoHundredFiftySixRebuilds}{10}
% E7 w3_N20_k2 drb-id-rb alpha=0.25 period=off: mean additions per non-rebuild update [configs/E7.json, aggregate.py]
\newcommand{\NumESevenWThreeNTwentyDrbIdRbAQuarterPOffAdds}{5{,}150}
% E7 w3_N20_k2 drb-id-rb alpha=0.25 period=off: amortized latency (us) [configs/E7.json, aggregate.py]
\newcommand{\NumESevenWThreeNTwentyDrbIdRbAQuarterPOffAmort}{2{,}130}
% E7 w3_N20_k2 drb-id-rb alpha=0.25 period=off: median latency (us) [configs/E7.json, aggregate.py]
\newcommand{\NumESevenWThreeNTwentyDrbIdRbAQuarterPOffMed}{1{,}200}
% E7 w3_N20_k2 drb-id-rb alpha=0.25 period=off: p99 latency (us) [configs/E7.json, aggregate.py]
\newcommand{\NumESevenWThreeNTwentyDrbIdRbAQuarterPOffPNinetyNine}{32{,}300}
% E7 w3_N20_k2 drb-id-rb alpha=0.25 period=off: rebuilds per run (mean of seeds) [configs/E7.json, aggregate.py]
\newcommand{\NumESevenWThreeNTwentyDrbIdRbAQuarterPOffRebuilds}{44}
% E7 w3_N20_k2 drb-id-rb alpha=0.25 period=64: mean additions per non-rebuild update [configs/E7.json, aggregate.py]
\newcommand{\NumESevenWThreeNTwentyDrbIdRbAQuarterPSixtyFourAdds}{5{,}150}
% E7 w3_N20_k2 drb-id-rb alpha=0.25 period=64: amortized latency (us) [configs/E7.json, aggregate.py]
\newcommand{\NumESevenWThreeNTwentyDrbIdRbAQuarterPSixtyFourAmort}{2{,}130}
% E7 w3_N20_k2 drb-id-rb alpha=0.25 period=64: median latency (us) [configs/E7.json, aggregate.py]
\newcommand{\NumESevenWThreeNTwentyDrbIdRbAQuarterPSixtyFourMed}{1{,}180}
% E7 w3_N20_k2 drb-id-rb alpha=0.25 period=64: p99 latency (us) [configs/E7.json, aggregate.py]
\newcommand{\NumESevenWThreeNTwentyDrbIdRbAQuarterPSixtyFourPNinetyNine}{31{,}900}
% E7 w3_N20_k2 drb-id-rb alpha=0.25 period=64: rebuilds per run (mean of seeds) [configs/E7.json, aggregate.py]
\newcommand{\NumESevenWThreeNTwentyDrbIdRbAQuarterPSixtyFourRebuilds}{44}
% E7 w3_N20_k2 drb-id-rb alpha=0.25 period=256: mean additions per non-rebuild update [configs/E7.json, aggregate.py]
\newcommand{\NumESevenWThreeNTwentyDrbIdRbAQuarterPTwoHundredFiftySixAdds}{5{,}150}
% E7 w3_N20_k2 drb-id-rb alpha=0.25 period=256: amortized latency (us) [configs/E7.json, aggregate.py]
\newcommand{\NumESevenWThreeNTwentyDrbIdRbAQuarterPTwoHundredFiftySixAmort}{2{,}150}
% E7 w3_N20_k2 drb-id-rb alpha=0.25 period=256: median latency (us) [configs/E7.json, aggregate.py]
\newcommand{\NumESevenWThreeNTwentyDrbIdRbAQuarterPTwoHundredFiftySixMed}{1{,}210}
% E7 w3_N20_k2 drb-id-rb alpha=0.25 period=256: p99 latency (us) [configs/E7.json, aggregate.py]
\newcommand{\NumESevenWThreeNTwentyDrbIdRbAQuarterPTwoHundredFiftySixPNinetyNine}{32{,}300}
% E7 w3_N20_k2 drb-id-rb alpha=0.25 period=256: rebuilds per run (mean of seeds) [configs/E7.json, aggregate.py]
\newcommand{\NumESevenWThreeNTwentyDrbIdRbAQuarterPTwoHundredFiftySixRebuilds}{44}
% E7 w3_N20_k2 drb-id-rb alpha=2 period=off: mean additions per non-rebuild update [configs/E7.json, aggregate.py]
\newcommand{\NumESevenWThreeNTwentyDrbIdRbATwoPOffAdds}{6{,}420}
% E7 w3_N20_k2 drb-id-rb alpha=2 period=off: amortized latency (us) [configs/E7.json, aggregate.py]
\newcommand{\NumESevenWThreeNTwentyDrbIdRbATwoPOffAmort}{2{,}880}
% E7 w3_N20_k2 drb-id-rb alpha=2 period=off: median latency (us) [configs/E7.json, aggregate.py]
\newcommand{\NumESevenWThreeNTwentyDrbIdRbATwoPOffMed}{1{,}860}
% E7 w3_N20_k2 drb-id-rb alpha=2 period=off: p99 latency (us) [configs/E7.json, aggregate.py]
\newcommand{\NumESevenWThreeNTwentyDrbIdRbATwoPOffPNinetyNine}{14{,}500}
% E7 w3_N20_k2 drb-id-rb alpha=2 period=off: rebuilds per run (mean of seeds) [configs/E7.json, aggregate.py]
\newcommand{\NumESevenWThreeNTwentyDrbIdRbATwoPOffRebuilds}{4.5}
% E7 w3_N20_k2 drb-id-rb alpha=2 period=64: mean additions per non-rebuild update [configs/E7.json, aggregate.py]
\newcommand{\NumESevenWThreeNTwentyDrbIdRbATwoPSixtyFourAdds}{5{,}350}
% E7 w3_N20_k2 drb-id-rb alpha=2 period=64: amortized latency (us) [configs/E7.json, aggregate.py]
\newcommand{\NumESevenWThreeNTwentyDrbIdRbATwoPSixtyFourAmort}{2{,}060}
% E7 w3_N20_k2 drb-id-rb alpha=2 period=64: median latency (us) [configs/E7.json, aggregate.py]
\newcommand{\NumESevenWThreeNTwentyDrbIdRbATwoPSixtyFourMed}{1{,}250}
% E7 w3_N20_k2 drb-id-rb alpha=2 period=64: p99 latency (us) [configs/E7.json, aggregate.py]
\newcommand{\NumESevenWThreeNTwentyDrbIdRbATwoPSixtyFourPNinetyNine}{31{,}600}
% E7 w3_N20_k2 drb-id-rb alpha=2 period=64: rebuilds per run (mean of seeds) [configs/E7.json, aggregate.py]
\newcommand{\NumESevenWThreeNTwentyDrbIdRbATwoPSixtyFourRebuilds}{31}
% E7 w3_N20_k2 drb-id-rb alpha=2 period=256: mean additions per non-rebuild update [configs/E7.json, aggregate.py]
\newcommand{\NumESevenWThreeNTwentyDrbIdRbATwoPTwoHundredFiftySixAdds}{6{,}210}
% E7 w3_N20_k2 drb-id-rb alpha=2 period=256: amortized latency (us) [configs/E7.json, aggregate.py]
\newcommand{\NumESevenWThreeNTwentyDrbIdRbATwoPTwoHundredFiftySixAmort}{2{,}490}
% E7 w3_N20_k2 drb-id-rb alpha=2 period=256: median latency (us) [configs/E7.json, aggregate.py]
\newcommand{\NumESevenWThreeNTwentyDrbIdRbATwoPTwoHundredFiftySixMed}{1{,}560}
% E7 w3_N20_k2 drb-id-rb alpha=2 period=256: p99 latency (us) [configs/E7.json, aggregate.py]
\newcommand{\NumESevenWThreeNTwentyDrbIdRbATwoPTwoHundredFiftySixPNinetyNine}{13{,}000}
% E7 w3_N20_k2 drb-id-rb alpha=2 period=256: rebuilds per run (mean of seeds) [configs/E7.json, aggregate.py]
\newcommand{\NumESevenWThreeNTwentyDrbIdRbATwoPTwoHundredFiftySixRebuilds}{7.0}
% E7 w3_N20_k2 mmhm alpha=4 period=off: mean additions per non-rebuild update [configs/E7.json, aggregate.py]
\newcommand{\NumESevenWThreeNTwentyMmhmAFourPOffAdds}{5{,}420}
% E7 w3_N20_k2 mmhm alpha=4 period=off: amortized latency (us) [configs/E7.json, aggregate.py]
\newcommand{\NumESevenWThreeNTwentyMmhmAFourPOffAmort}{3{,}750}
% E7 w3_N20_k2 mmhm alpha=4 period=off: median latency (us) [configs/E7.json, aggregate.py]
\newcommand{\NumESevenWThreeNTwentyMmhmAFourPOffMed}{2{,}180}
% E7 w3_N20_k2 mmhm alpha=4 period=off: p99 latency (us) [configs/E7.json, aggregate.py]
\newcommand{\NumESevenWThreeNTwentyMmhmAFourPOffPNinetyNine}{24{,}800}
% E7 w3_N20_k2 mmhm alpha=4 period=off: rebuilds per run (mean of seeds) [configs/E7.json, aggregate.py]
\newcommand{\NumESevenWThreeNTwentyMmhmAFourPOffRebuilds}{0.50}
% E7 w3_N20_k2 mmhm alpha=4 period=64: mean additions per non-rebuild update [configs/E7.json, aggregate.py]
\newcommand{\NumESevenWThreeNTwentyMmhmAFourPSixtyFourAdds}{4{,}860}
% E7 w3_N20_k2 mmhm alpha=4 period=64: amortized latency (us) [configs/E7.json, aggregate.py]
\newcommand{\NumESevenWThreeNTwentyMmhmAFourPSixtyFourAmort}{1{,}950}
% E7 w3_N20_k2 mmhm alpha=4 period=64: median latency (us) [configs/E7.json, aggregate.py]
\newcommand{\NumESevenWThreeNTwentyMmhmAFourPSixtyFourMed}{957}
% E7 w3_N20_k2 mmhm alpha=4 period=64: p99 latency (us) [configs/E7.json, aggregate.py]
\newcommand{\NumESevenWThreeNTwentyMmhmAFourPSixtyFourPNinetyNine}{39{,}800}
% E7 w3_N20_k2 mmhm alpha=4 period=64: rebuilds per run (mean of seeds) [configs/E7.json, aggregate.py]
\newcommand{\NumESevenWThreeNTwentyMmhmAFourPSixtyFourRebuilds}{31}
% E7 w3_N20_k2 mmhm alpha=4 period=256: mean additions per non-rebuild update [configs/E7.json, aggregate.py]
\newcommand{\NumESevenWThreeNTwentyMmhmAFourPTwoHundredFiftySixAdds}{4{,}830}
% E7 w3_N20_k2 mmhm alpha=4 period=256: amortized latency (us) [configs/E7.json, aggregate.py]
\newcommand{\NumESevenWThreeNTwentyMmhmAFourPTwoHundredFiftySixAmort}{1{,}720}
% E7 w3_N20_k2 mmhm alpha=4 period=256: median latency (us) [configs/E7.json, aggregate.py]
\newcommand{\NumESevenWThreeNTwentyMmhmAFourPTwoHundredFiftySixMed}{1{,}100}
% E7 w3_N20_k2 mmhm alpha=4 period=256: p99 latency (us) [configs/E7.json, aggregate.py]
\newcommand{\NumESevenWThreeNTwentyMmhmAFourPTwoHundredFiftySixPNinetyNine}{7{,}400}
% E7 w3_N20_k2 mmhm alpha=4 period=256: rebuilds per run (mean of seeds) [configs/E7.json, aggregate.py]
\newcommand{\NumESevenWThreeNTwentyMmhmAFourPTwoHundredFiftySixRebuilds}{7.0}
% E7 w3_N20_k2 mmhm alpha=0.5 period=off: mean additions per non-rebuild update [configs/E7.json, aggregate.py]
\newcommand{\NumESevenWThreeNTwentyMmhmAHalfPOffAdds}{4{,}540}
% E7 w3_N20_k2 mmhm alpha=0.5 period=off: amortized latency (us) [configs/E7.json, aggregate.py]
\newcommand{\NumESevenWThreeNTwentyMmhmAHalfPOffAmort}{1{,}620}
% E7 w3_N20_k2 mmhm alpha=0.5 period=off: median latency (us) [configs/E7.json, aggregate.py]
\newcommand{\NumESevenWThreeNTwentyMmhmAHalfPOffMed}{961}
% E7 w3_N20_k2 mmhm alpha=0.5 period=off: p99 latency (us) [configs/E7.json, aggregate.py]
\newcommand{\NumESevenWThreeNTwentyMmhmAHalfPOffPNinetyNine}{6{,}790}
% E7 w3_N20_k2 mmhm alpha=0.5 period=off: rebuilds per run (mean of seeds) [configs/E7.json, aggregate.py]
\newcommand{\NumESevenWThreeNTwentyMmhmAHalfPOffRebuilds}{15}
% E7 w3_N20_k2 mmhm alpha=0.5 period=64: mean additions per non-rebuild update [configs/E7.json, aggregate.py]
\newcommand{\NumESevenWThreeNTwentyMmhmAHalfPSixtyFourAdds}{4{,}860}
% E7 w3_N20_k2 mmhm alpha=0.5 period=64: amortized latency (us) [configs/E7.json, aggregate.py]
\newcommand{\NumESevenWThreeNTwentyMmhmAHalfPSixtyFourAmort}{1{,}810}
% E7 w3_N20_k2 mmhm alpha=0.5 period=64: median latency (us) [configs/E7.json, aggregate.py]
\newcommand{\NumESevenWThreeNTwentyMmhmAHalfPSixtyFourMed}{892}
% E7 w3_N20_k2 mmhm alpha=0.5 period=64: p99 latency (us) [configs/E7.json, aggregate.py]
\newcommand{\NumESevenWThreeNTwentyMmhmAHalfPSixtyFourPNinetyNine}{38{,}500}
% E7 w3_N20_k2 mmhm alpha=0.5 period=64: rebuilds per run (mean of seeds) [configs/E7.json, aggregate.py]
\newcommand{\NumESevenWThreeNTwentyMmhmAHalfPSixtyFourRebuilds}{31}
% E7 w3_N20_k2 mmhm alpha=0.5 period=256: mean additions per non-rebuild update [configs/E7.json, aggregate.py]
\newcommand{\NumESevenWThreeNTwentyMmhmAHalfPTwoHundredFiftySixAdds}{4{,}540}
% E7 w3_N20_k2 mmhm alpha=0.5 period=256: amortized latency (us) [configs/E7.json, aggregate.py]
\newcommand{\NumESevenWThreeNTwentyMmhmAHalfPTwoHundredFiftySixAmort}{1{,}570}
% E7 w3_N20_k2 mmhm alpha=0.5 period=256: median latency (us) [configs/E7.json, aggregate.py]
\newcommand{\NumESevenWThreeNTwentyMmhmAHalfPTwoHundredFiftySixMed}{940}
% E7 w3_N20_k2 mmhm alpha=0.5 period=256: p99 latency (us) [configs/E7.json, aggregate.py]
\newcommand{\NumESevenWThreeNTwentyMmhmAHalfPTwoHundredFiftySixPNinetyNine}{6{,}200}
% E7 w3_N20_k2 mmhm alpha=0.5 period=256: rebuilds per run (mean of seeds) [configs/E7.json, aggregate.py]
\newcommand{\NumESevenWThreeNTwentyMmhmAHalfPTwoHundredFiftySixRebuilds}{15}
% E7 w3_N20_k2 mmhm alpha=off period=off: mean additions per non-rebuild update [configs/E7.json, aggregate.py]
\newcommand{\NumESevenWThreeNTwentyMmhmAOffPOffAdds}{5{,}440}
% E7 w3_N20_k2 mmhm alpha=off period=off: amortized latency (us) [configs/E7.json, aggregate.py]
\newcommand{\NumESevenWThreeNTwentyMmhmAOffPOffAmort}{3{,}820}
% E7 w3_N20_k2 mmhm alpha=off period=off: median latency (us) [configs/E7.json, aggregate.py]
\newcommand{\NumESevenWThreeNTwentyMmhmAOffPOffMed}{2{,}230}
% E7 w3_N20_k2 mmhm alpha=off period=off: p99 latency (us) [configs/E7.json, aggregate.py]
\newcommand{\NumESevenWThreeNTwentyMmhmAOffPOffPNinetyNine}{25{,}400}
% E7 w3_N20_k2 mmhm alpha=off period=off: rebuilds per run (mean of seeds) [configs/E7.json, aggregate.py]
\newcommand{\NumESevenWThreeNTwentyMmhmAOffPOffRebuilds}{0}
% E7 w3_N20_k2 mmhm alpha=off period=64: mean additions per non-rebuild update [configs/E7.json, aggregate.py]
\newcommand{\NumESevenWThreeNTwentyMmhmAOffPSixtyFourAdds}{4{,}860}
% E7 w3_N20_k2 mmhm alpha=off period=64: amortized latency (us) [configs/E7.json, aggregate.py]
\newcommand{\NumESevenWThreeNTwentyMmhmAOffPSixtyFourAmort}{1{,}940}
% E7 w3_N20_k2 mmhm alpha=off period=64: median latency (us) [configs/E7.json, aggregate.py]
\newcommand{\NumESevenWThreeNTwentyMmhmAOffPSixtyFourMed}{960}
% E7 w3_N20_k2 mmhm alpha=off period=64: p99 latency (us) [configs/E7.json, aggregate.py]
\newcommand{\NumESevenWThreeNTwentyMmhmAOffPSixtyFourPNinetyNine}{40{,}700}
% E7 w3_N20_k2 mmhm alpha=off period=64: rebuilds per run (mean of seeds) [configs/E7.json, aggregate.py]
\newcommand{\NumESevenWThreeNTwentyMmhmAOffPSixtyFourRebuilds}{31}
% E7 w3_N20_k2 mmhm alpha=off period=256: mean additions per non-rebuild update [configs/E7.json, aggregate.py]
\newcommand{\NumESevenWThreeNTwentyMmhmAOffPTwoHundredFiftySixAdds}{4{,}830}
% E7 w3_N20_k2 mmhm alpha=off period=256: amortized latency (us) [configs/E7.json, aggregate.py]
\newcommand{\NumESevenWThreeNTwentyMmhmAOffPTwoHundredFiftySixAmort}{1{,}640}
% E7 w3_N20_k2 mmhm alpha=off period=256: median latency (us) [configs/E7.json, aggregate.py]
\newcommand{\NumESevenWThreeNTwentyMmhmAOffPTwoHundredFiftySixMed}{1{,}060}
% E7 w3_N20_k2 mmhm alpha=off period=256: p99 latency (us) [configs/E7.json, aggregate.py]
\newcommand{\NumESevenWThreeNTwentyMmhmAOffPTwoHundredFiftySixPNinetyNine}{6{,}810}
% E7 w3_N20_k2 mmhm alpha=off period=256: rebuilds per run (mean of seeds) [configs/E7.json, aggregate.py]
\newcommand{\NumESevenWThreeNTwentyMmhmAOffPTwoHundredFiftySixRebuilds}{7.0}
% E7 w3_N20_k2 mmhm alpha=1 period=off: mean additions per non-rebuild update [configs/E7.json, aggregate.py]
\newcommand{\NumESevenWThreeNTwentyMmhmAOnePOffAdds}{4{,}730}
% E7 w3_N20_k2 mmhm alpha=1 period=off: amortized latency (us) [configs/E7.json, aggregate.py]
\newcommand{\NumESevenWThreeNTwentyMmhmAOnePOffAmort}{1{,}510}
% E7 w3_N20_k2 mmhm alpha=1 period=off: median latency (us) [configs/E7.json, aggregate.py]
\newcommand{\NumESevenWThreeNTwentyMmhmAOnePOffMed}{1{,}000}
% E7 w3_N20_k2 mmhm alpha=1 period=off: p99 latency (us) [configs/E7.json, aggregate.py]
\newcommand{\NumESevenWThreeNTwentyMmhmAOnePOffPNinetyNine}{6{,}290}
% E7 w3_N20_k2 mmhm alpha=1 period=off: rebuilds per run (mean of seeds) [configs/E7.json, aggregate.py]
\newcommand{\NumESevenWThreeNTwentyMmhmAOnePOffRebuilds}{6.0}
% E7 w3_N20_k2 mmhm alpha=1 period=64: mean additions per non-rebuild update [configs/E7.json, aggregate.py]
\newcommand{\NumESevenWThreeNTwentyMmhmAOnePSixtyFourAdds}{4{,}860}
% E7 w3_N20_k2 mmhm alpha=1 period=64: amortized latency (us) [configs/E7.json, aggregate.py]
\newcommand{\NumESevenWThreeNTwentyMmhmAOnePSixtyFourAmort}{1{,}940}
% E7 w3_N20_k2 mmhm alpha=1 period=64: median latency (us) [configs/E7.json, aggregate.py]
\newcommand{\NumESevenWThreeNTwentyMmhmAOnePSixtyFourMed}{945}
% E7 w3_N20_k2 mmhm alpha=1 period=64: p99 latency (us) [configs/E7.json, aggregate.py]
\newcommand{\NumESevenWThreeNTwentyMmhmAOnePSixtyFourPNinetyNine}{40{,}100}
% E7 w3_N20_k2 mmhm alpha=1 period=64: rebuilds per run (mean of seeds) [configs/E7.json, aggregate.py]
\newcommand{\NumESevenWThreeNTwentyMmhmAOnePSixtyFourRebuilds}{31}
% E7 w3_N20_k2 mmhm alpha=1 period=256: mean additions per non-rebuild update [configs/E7.json, aggregate.py]
\newcommand{\NumESevenWThreeNTwentyMmhmAOnePTwoHundredFiftySixAdds}{4{,}820}
% E7 w3_N20_k2 mmhm alpha=1 period=256: amortized latency (us) [configs/E7.json, aggregate.py]
\newcommand{\NumESevenWThreeNTwentyMmhmAOnePTwoHundredFiftySixAmort}{1{,}620}
% E7 w3_N20_k2 mmhm alpha=1 period=256: median latency (us) [configs/E7.json, aggregate.py]
\newcommand{\NumESevenWThreeNTwentyMmhmAOnePTwoHundredFiftySixMed}{1{,}040}
% E7 w3_N20_k2 mmhm alpha=1 period=256: p99 latency (us) [configs/E7.json, aggregate.py]
\newcommand{\NumESevenWThreeNTwentyMmhmAOnePTwoHundredFiftySixPNinetyNine}{6{,}680}
% E7 w3_N20_k2 mmhm alpha=1 period=256: rebuilds per run (mean of seeds) [configs/E7.json, aggregate.py]
\newcommand{\NumESevenWThreeNTwentyMmhmAOnePTwoHundredFiftySixRebuilds}{7.0}
% E7 w3_N20_k2 mmhm alpha=0.25 period=off: mean additions per non-rebuild update [configs/E7.json, aggregate.py]
\newcommand{\NumESevenWThreeNTwentyMmhmAQuarterPOffAdds}{4{,}690}
% E7 w3_N20_k2 mmhm alpha=0.25 period=off: amortized latency (us) [configs/E7.json, aggregate.py]
\newcommand{\NumESevenWThreeNTwentyMmhmAQuarterPOffAmort}{1{,}890}
% E7 w3_N20_k2 mmhm alpha=0.25 period=off: median latency (us) [configs/E7.json, aggregate.py]
\newcommand{\NumESevenWThreeNTwentyMmhmAQuarterPOffMed}{913}
% E7 w3_N20_k2 mmhm alpha=0.25 period=off: p99 latency (us) [configs/E7.json, aggregate.py]
\newcommand{\NumESevenWThreeNTwentyMmhmAQuarterPOffPNinetyNine}{40{,}300}
% E7 w3_N20_k2 mmhm alpha=0.25 period=off: rebuilds per run (mean of seeds) [configs/E7.json, aggregate.py]
\newcommand{\NumESevenWThreeNTwentyMmhmAQuarterPOffRebuilds}{35}
% E7 w3_N20_k2 mmhm alpha=0.25 period=64: mean additions per non-rebuild update [configs/E7.json, aggregate.py]
\newcommand{\NumESevenWThreeNTwentyMmhmAQuarterPSixtyFourAdds}{4{,}800}
% E7 w3_N20_k2 mmhm alpha=0.25 period=64: amortized latency (us) [configs/E7.json, aggregate.py]
\newcommand{\NumESevenWThreeNTwentyMmhmAQuarterPSixtyFourAmort}{1{,}960}
% E7 w3_N20_k2 mmhm alpha=0.25 period=64: median latency (us) [configs/E7.json, aggregate.py]
\newcommand{\NumESevenWThreeNTwentyMmhmAQuarterPSixtyFourMed}{927}
% E7 w3_N20_k2 mmhm alpha=0.25 period=64: p99 latency (us) [configs/E7.json, aggregate.py]
\newcommand{\NumESevenWThreeNTwentyMmhmAQuarterPSixtyFourPNinetyNine}{40{,}300}
% E7 w3_N20_k2 mmhm alpha=0.25 period=64: rebuilds per run (mean of seeds) [configs/E7.json, aggregate.py]
\newcommand{\NumESevenWThreeNTwentyMmhmAQuarterPSixtyFourRebuilds}{37}
% E7 w3_N20_k2 mmhm alpha=0.25 period=256: mean additions per non-rebuild update [configs/E7.json, aggregate.py]
\newcommand{\NumESevenWThreeNTwentyMmhmAQuarterPTwoHundredFiftySixAdds}{4{,}690}
% E7 w3_N20_k2 mmhm alpha=0.25 period=256: amortized latency (us) [configs/E7.json, aggregate.py]
\newcommand{\NumESevenWThreeNTwentyMmhmAQuarterPTwoHundredFiftySixAmort}{1{,}920}
% E7 w3_N20_k2 mmhm alpha=0.25 period=256: median latency (us) [configs/E7.json, aggregate.py]
\newcommand{\NumESevenWThreeNTwentyMmhmAQuarterPTwoHundredFiftySixMed}{919}
% E7 w3_N20_k2 mmhm alpha=0.25 period=256: p99 latency (us) [configs/E7.json, aggregate.py]
\newcommand{\NumESevenWThreeNTwentyMmhmAQuarterPTwoHundredFiftySixPNinetyNine}{40{,}200}
% E7 w3_N20_k2 mmhm alpha=0.25 period=256: rebuilds per run (mean of seeds) [configs/E7.json, aggregate.py]
\newcommand{\NumESevenWThreeNTwentyMmhmAQuarterPTwoHundredFiftySixRebuilds}{35}
% E7 w3_N20_k2 mmhm alpha=2 period=off: mean additions per non-rebuild update [configs/E7.json, aggregate.py]
\newcommand{\NumESevenWThreeNTwentyMmhmATwoPOffAdds}{5{,}280}
% E7 w3_N20_k2 mmhm alpha=2 period=off: amortized latency (us) [configs/E7.json, aggregate.py]
\newcommand{\NumESevenWThreeNTwentyMmhmATwoPOffAmort}{2{,}090}
% E7 w3_N20_k2 mmhm alpha=2 period=off: median latency (us) [configs/E7.json, aggregate.py]
\newcommand{\NumESevenWThreeNTwentyMmhmATwoPOffMed}{1{,}360}
% E7 w3_N20_k2 mmhm alpha=2 period=off: p99 latency (us) [configs/E7.json, aggregate.py]
\newcommand{\NumESevenWThreeNTwentyMmhmATwoPOffPNinetyNine}{10{,}800}
% E7 w3_N20_k2 mmhm alpha=2 period=off: rebuilds per run (mean of seeds) [configs/E7.json, aggregate.py]
\newcommand{\NumESevenWThreeNTwentyMmhmATwoPOffRebuilds}{2.5}
% E7 w3_N20_k2 mmhm alpha=2 period=64: mean additions per non-rebuild update [configs/E7.json, aggregate.py]
\newcommand{\NumESevenWThreeNTwentyMmhmATwoPSixtyFourAdds}{4{,}860}
% E7 w3_N20_k2 mmhm alpha=2 period=64: amortized latency (us) [configs/E7.json, aggregate.py]
\newcommand{\NumESevenWThreeNTwentyMmhmATwoPSixtyFourAmort}{1{,}920}
% E7 w3_N20_k2 mmhm alpha=2 period=64: median latency (us) [configs/E7.json, aggregate.py]
\newcommand{\NumESevenWThreeNTwentyMmhmATwoPSixtyFourMed}{951}
% E7 w3_N20_k2 mmhm alpha=2 period=64: p99 latency (us) [configs/E7.json, aggregate.py]
\newcommand{\NumESevenWThreeNTwentyMmhmATwoPSixtyFourPNinetyNine}{40{,}100}
% E7 w3_N20_k2 mmhm alpha=2 period=64: rebuilds per run (mean of seeds) [configs/E7.json, aggregate.py]
\newcommand{\NumESevenWThreeNTwentyMmhmATwoPSixtyFourRebuilds}{31}
% E7 w3_N20_k2 mmhm alpha=2 period=256: mean additions per non-rebuild update [configs/E7.json, aggregate.py]
\newcommand{\NumESevenWThreeNTwentyMmhmATwoPTwoHundredFiftySixAdds}{4{,}830}
% E7 w3_N20_k2 mmhm alpha=2 period=256: amortized latency (us) [configs/E7.json, aggregate.py]
\newcommand{\NumESevenWThreeNTwentyMmhmATwoPTwoHundredFiftySixAmort}{1{,}780}
% E7 w3_N20_k2 mmhm alpha=2 period=256: median latency (us) [configs/E7.json, aggregate.py]
\newcommand{\NumESevenWThreeNTwentyMmhmATwoPTwoHundredFiftySixMed}{1{,}150}
% E7 w3_N20_k2 mmhm alpha=2 period=256: p99 latency (us) [configs/E7.json, aggregate.py]
\newcommand{\NumESevenWThreeNTwentyMmhmATwoPTwoHundredFiftySixPNinetyNine}{7{,}460}
% E7 w3_N20_k2 mmhm alpha=2 period=256: rebuilds per run (mean of seeds) [configs/E7.json, aggregate.py]
\newcommand{\NumESevenWThreeNTwentyMmhmATwoPTwoHundredFiftySixRebuilds}{7.0}
% mismatching states of exact methods, all oracles [verify.py]
\newcommand{\NumExactMismatches}{0}
% runs of exact methods checked (E1-E7) [verify.py]
\newcommand{\NumExactRuns}{348}
% gudhi runs all T updates iff |K_0| <= this [configs/make_configs.py]
\newcommand{\NumGudhiFullLimit}{300{,}000}
% largest |K_0| at which gudhi is timed on all T updates [aggregate.py]
\newcommand{\NumGudhiFullSizeMax}{286{,}589}
% largest |K_0| at which gudhi is timed (on the prefix T') [aggregate.py]
\newcommand{\NumGudhiSizeMax}{3{,}932{,}162}
% internal full-id verifications (--verify-every 50 + last update) in exact runs [verify.py]
\newcommand{\NumInternalChecks}{9{,}717}
% lcr: % of wrong updates over all E0 lcr runs [verify.py]
\newcommand{\NumLcrErrAll}{83.1}
% lcr: largest per-workload % of wrong updates [verify.py]
\newcommand{\NumLcrErrMax}{100.0}
% lcr: smallest per-workload % of wrong updates [verify.py]
\newcommand{\NumLcrErrMin}{0.0}
% lcr on w5a_L10000: % of updates with a wrong Betti vector (3 seeds, --verify-every 1) [configs/E0.json, verify.py]
\newcommand{\NumLcrWFiveCycleLTenThousandErrPct}{0.0}
% lcr on w5b_L5: % of updates with a wrong Betti vector (3 seeds, --verify-every 1) [configs/E0.json, verify.py]
\newcommand{\NumLcrWFiveSphereLFiveErrPct}{50.0}
% lcr on w5b_L5: first wrong update (min over seeds) [configs/E0.json, verify.py]
\newcommand{\NumLcrWFiveSphereLFiveFirstWrong}{1}
% lcr on w4_2d_n4000_D3: % of updates with a wrong Betti vector (3 seeds, --verify-every 1) [configs/E0.json, verify.py]
\newcommand{\NumLcrWFourPlaneNFourThousandDThreeErrPct}{99.6}
% lcr on w4_2d_n4000_D3: first wrong update (min over seeds) [configs/E0.json, verify.py]
\newcommand{\NumLcrWFourPlaneNFourThousandDThreeFirstWrong}{5}
% lcr on w4_2d_n4000_D2: % of updates with a wrong Betti vector (3 seeds, --verify-every 1) [configs/E0.json, verify.py]
\newcommand{\NumLcrWFourPlaneNFourThousandDTwoErrPct}{99.6}
% lcr on w4_2d_n4000_D2: first wrong update (min over seeds) [configs/E0.json, verify.py]
\newcommand{\NumLcrWFourPlaneNFourThousandDTwoFirstWrong}{5}
% lcr on w4_3d_n4000_D3: % of updates with a wrong Betti vector (3 seeds, --verify-every 1) [configs/E0.json, verify.py]
\newcommand{\NumLcrWFourSpaceNFourThousandDThreeErrPct}{99.9}
% lcr on w4_3d_n4000_D3: first wrong update (min over seeds) [configs/E0.json, verify.py]
\newcommand{\NumLcrWFourSpaceNFourThousandDThreeFirstWrong}{2}
% lcr on w1_sphere_L5: % of updates with a wrong Betti vector (3 seeds, --verify-every 1) [configs/E0.json, verify.py]
\newcommand{\NumLcrWOneSphereLFiveErrPct}{99.5}
% lcr on w1_sphere_L5: first wrong update (min over seeds) [configs/E0.json, verify.py]
\newcommand{\NumLcrWOneSphereLFiveFirstWrong}{1}
% lcr on w1_torus_128: % of updates with a wrong Betti vector (3 seeds, --verify-every 1) [configs/E0.json, verify.py]
\newcommand{\NumLcrWOneTorusOneHundredTwentyEightErrPct}{99.3}
% lcr on w1_torus_128: first wrong update (min over seeds) [configs/E0.json, verify.py]
\newcommand{\NumLcrWOneTorusOneHundredTwentyEightFirstWrong}{1}
% lcr on w3_N14_k2: % of updates with a wrong Betti vector (3 seeds, --verify-every 1) [configs/E0.json, verify.py]
\newcommand{\NumLcrWThreeNFourteenErrPct}{100.0}
% lcr on w3_N14_k2: first wrong update (min over seeds) [configs/E0.json, verify.py]
\newcommand{\NumLcrWThreeNFourteenFirstWrong}{1}
% lcr on w2_N15_s8: % of updates with a wrong Betti vector (3 seeds, --verify-every 1) [configs/E0.json, verify.py]
\newcommand{\NumLcrWTwoNFifteenSEightErrPct}{100.0}
% lcr on w2_N15_s8: first wrong update (min over seeds) [configs/E0.json, verify.py]
\newcommand{\NumLcrWTwoNFifteenSEightFirstWrong}{1}
% number of E0 lcr instances [verify.py]
\newcommand{\NumLcrWorkloads}{9}
% number of E0 lcr instances with at least one wrong update [verify.py]
\newcommand{\NumLcrWorkloadsWrong}{8}
% noise: CV (%) of the amortized latency over 5 identical runs of full-id [noise.py]
\newcommand{\NumNoiseCVFullIdAmort}{4.5}
% noise: CV (%) of the median per-update latency over 5 identical runs of full-id on w2_N22_s8_T2000_seed0 [noise.py]
\newcommand{\NumNoiseCVFullIdMed}{4.1}
% noise: CV (%) of the amortized latency over 5 identical runs of mmhm [noise.py]
\newcommand{\NumNoiseCVMmhmAmort}{8.4}
% noise: CV (%) of the median per-update latency over 5 identical runs of mmhm on w2_N22_s8_T2000_seed0 [noise.py]
\newcommand{\NumNoiseCVMmhmMed}{7.1}
% noise: change (%) of the median latency of full-id with one timed + two GUDHI runs running concurrently, vs alone [noise.py]
\newcommand{\NumNoiseFullIdFourBusyPct}{+1.5}
% noise: change (%) of the median latency of full-id with one other timed run running concurrently, vs alone [noise.py]
\newcommand{\NumNoiseFullIdTwoTimedPct}{+0.3}
% noise: change (%) of the median latency of mmhm with one timed + two GUDHI runs running concurrently, vs alone [noise.py]
\newcommand{\NumNoiseMmhmFourBusyPct}{-3.6}
% noise: change (%) of the median latency of mmhm with one other timed run running concurrently, vs alone [noise.py]
\newcommand{\NumNoiseMmhmTwoTimedPct}{-2.3}
% noise: identical repetitions per method [noise.py]
\newcommand{\NumNoiseReps}{5}
% T': prefix of the gudhi baseline above 3e5 simplices [configs/make_configs.py]
\newcommand{\NumPrefix}{200}
% Betti vectors reported by exact runs (INIT + updates) [verify.py]
\newcommand{\NumReportedStates}{647{,}748}
% sample cross-check: Betti vectors of exact configurations compared with GUDHI (all samples) [crosscheck_samples.py]
\newcommand{\NumSampleChecks}{9{,}218}
% samples on which lcr is wrong at least once [crosscheck_samples.py]
\newcommand{\NumSampleLcrWrong}{7}
% sample cross-check: mismatching states of exact configurations [crosscheck_samples.py]
\newcommand{\NumSampleMismatches}{0}
% committed workload samples [crosscheck_samples.py]
\newcommand{\NumSamples}{8}
% seeds per configuration [configs/make_configs.py]
\newcommand{\NumSeeds}{3}
% largest |K_0| in a timed run [aggregate.py]
\newcommand{\NumSizeMax}{3{,}932{,}162}
% smallest |K_0| in a timed run [aggregate.py]
\newcommand{\NumSizeMin}{10{,}739}
% E5 WOneSphere full-morse: least-squares slope of log10(amortized latency) vs log10(|K_0|), 2 points (w1_sphere_L6 w1_sphere_L7) [aggregate.py]
\newcommand{\NumSlopeEFiveWOneSphereFullMorseAmort}{1.37}
% number of points of \NumSlopeEFiveWOneSphereFullMorseAmort [aggregate.py]
\newcommand{\NumSlopeEFiveWOneSphereFullMorseAmortN}{2}
% R^2 of \NumSlopeEFiveWOneSphereFullMorseAmort [aggregate.py]
\newcommand{\NumSlopeEFiveWOneSphereFullMorseAmortRSq}{1.000}
% E5 WOneSphere full-morse: least-squares slope of log10(median latency) vs log10(|K_0|), 2 points (w1_sphere_L6 w1_sphere_L7) [aggregate.py]
\newcommand{\NumSlopeEFiveWOneSphereFullMorseMed}{1.39}
% number of points of \NumSlopeEFiveWOneSphereFullMorseMed [aggregate.py]
\newcommand{\NumSlopeEFiveWOneSphereFullMorseMedN}{2}
% R^2 of \NumSlopeEFiveWOneSphereFullMorseMed [aggregate.py]
\newcommand{\NumSlopeEFiveWOneSphereFullMorseMedRSq}{1.000}
% E5 WOneSphere gudhi: least-squares slope of log10(amortized latency) vs log10(|K_0|), 3 points (w1_sphere_L5 w1_sphere_L6 w1_sphere_L7) [aggregate.py]
\newcommand{\NumSlopeEFiveWOneSphereGudhiAmort}{1.03}
% number of points of \NumSlopeEFiveWOneSphereGudhiAmort [aggregate.py]
\newcommand{\NumSlopeEFiveWOneSphereGudhiAmortN}{3}
% R^2 of \NumSlopeEFiveWOneSphereGudhiAmort [aggregate.py]
\newcommand{\NumSlopeEFiveWOneSphereGudhiAmortRSq}{1.000}
% E5 WOneSphere gudhi: least-squares slope of log10(median latency) vs log10(|K_0|), 3 points (w1_sphere_L5 w1_sphere_L6 w1_sphere_L7) [aggregate.py]
\newcommand{\NumSlopeEFiveWOneSphereGudhiMed}{1.03}
% number of points of \NumSlopeEFiveWOneSphereGudhiMed [aggregate.py]
\newcommand{\NumSlopeEFiveWOneSphereGudhiMedN}{3}
% R^2 of \NumSlopeEFiveWOneSphereGudhiMed [aggregate.py]
\newcommand{\NumSlopeEFiveWOneSphereGudhiMedRSq}{1.000}
% E5 WOneTorus full-morse: least-squares slope of log10(amortized latency) vs log10(|K_0|), 2 points (w1_torus_256 w1_torus_500) [aggregate.py]
\newcommand{\NumSlopeEFiveWOneTorusFullMorseAmort}{1.34}
% number of points of \NumSlopeEFiveWOneTorusFullMorseAmort [aggregate.py]
\newcommand{\NumSlopeEFiveWOneTorusFullMorseAmortN}{2}
% R^2 of \NumSlopeEFiveWOneTorusFullMorseAmort [aggregate.py]
\newcommand{\NumSlopeEFiveWOneTorusFullMorseAmortRSq}{1.000}
% E5 WOneTorus full-morse: least-squares slope of log10(median latency) vs log10(|K_0|), 2 points (w1_torus_256 w1_torus_500) [aggregate.py]
\newcommand{\NumSlopeEFiveWOneTorusFullMorseMed}{1.36}
% number of points of \NumSlopeEFiveWOneTorusFullMorseMed [aggregate.py]
\newcommand{\NumSlopeEFiveWOneTorusFullMorseMedN}{2}
% R^2 of \NumSlopeEFiveWOneTorusFullMorseMed [aggregate.py]
\newcommand{\NumSlopeEFiveWOneTorusFullMorseMedRSq}{1.000}
% E5 WOneTorus gudhi: least-squares slope of log10(amortized latency) vs log10(|K_0|), 3 points (w1_torus_128 w1_torus_256 w1_torus_500) [aggregate.py]
\newcommand{\NumSlopeEFiveWOneTorusGudhiAmort}{1.05}
% number of points of \NumSlopeEFiveWOneTorusGudhiAmort [aggregate.py]
\newcommand{\NumSlopeEFiveWOneTorusGudhiAmortN}{3}
% R^2 of \NumSlopeEFiveWOneTorusGudhiAmort [aggregate.py]
\newcommand{\NumSlopeEFiveWOneTorusGudhiAmortRSq}{0.997}
% E5 WOneTorus gudhi: least-squares slope of log10(median latency) vs log10(|K_0|), 3 points (w1_torus_128 w1_torus_256 w1_torus_500) [aggregate.py]
\newcommand{\NumSlopeEFiveWOneTorusGudhiMed}{1.05}
% number of points of \NumSlopeEFiveWOneTorusGudhiMed [aggregate.py]
\newcommand{\NumSlopeEFiveWOneTorusGudhiMedN}{3}
% R^2 of \NumSlopeEFiveWOneTorusGudhiMed [aggregate.py]
\newcommand{\NumSlopeEFiveWOneTorusGudhiMedRSq}{0.998}
% E4 WFourPlaneDTwo full-morse: least-squares slope of log10(amortized latency) vs log10(|K_0|), 2 points (w4_2d_n4000_D2 w4_2d_n16000_D2) [aggregate.py]
\newcommand{\NumSlopeEFourWFourPlaneDTwoFullMorseAmort}{1.19}
% number of points of \NumSlopeEFourWFourPlaneDTwoFullMorseAmort [aggregate.py]
\newcommand{\NumSlopeEFourWFourPlaneDTwoFullMorseAmortN}{2}
% R^2 of \NumSlopeEFourWFourPlaneDTwoFullMorseAmort [aggregate.py]
\newcommand{\NumSlopeEFourWFourPlaneDTwoFullMorseAmortRSq}{1.000}
% E4 WFourPlaneDTwo full-morse: least-squares slope of log10(median latency) vs log10(|K_0|), 2 points (w4_2d_n4000_D2 w4_2d_n16000_D2) [aggregate.py]
\newcommand{\NumSlopeEFourWFourPlaneDTwoFullMorseMed}{1.19}
% number of points of \NumSlopeEFourWFourPlaneDTwoFullMorseMed [aggregate.py]
\newcommand{\NumSlopeEFourWFourPlaneDTwoFullMorseMedN}{2}
% R^2 of \NumSlopeEFourWFourPlaneDTwoFullMorseMed [aggregate.py]
\newcommand{\NumSlopeEFourWFourPlaneDTwoFullMorseMedRSq}{1.000}
% E4 WFourPlaneDTwo gudhi: least-squares slope of log10(amortized latency) vs log10(|K_0|), 2 points (w4_2d_n4000_D2 w4_2d_n16000_D2) [aggregate.py]
\newcommand{\NumSlopeEFourWFourPlaneDTwoGudhiAmort}{1.26}
% number of points of \NumSlopeEFourWFourPlaneDTwoGudhiAmort [aggregate.py]
\newcommand{\NumSlopeEFourWFourPlaneDTwoGudhiAmortN}{2}
% R^2 of \NumSlopeEFourWFourPlaneDTwoGudhiAmort [aggregate.py]
\newcommand{\NumSlopeEFourWFourPlaneDTwoGudhiAmortRSq}{1.000}
% E4 WFourPlaneDTwo gudhi: least-squares slope of log10(median latency) vs log10(|K_0|), 2 points (w4_2d_n4000_D2 w4_2d_n16000_D2) [aggregate.py]
\newcommand{\NumSlopeEFourWFourPlaneDTwoGudhiMed}{1.26}
% number of points of \NumSlopeEFourWFourPlaneDTwoGudhiMed [aggregate.py]
\newcommand{\NumSlopeEFourWFourPlaneDTwoGudhiMedN}{2}
% R^2 of \NumSlopeEFourWFourPlaneDTwoGudhiMed [aggregate.py]
\newcommand{\NumSlopeEFourWFourPlaneDTwoGudhiMedRSq}{1.000}
% E1 WTwoSEight full-id: least-squares slope of log10(amortized latency) vs log10(|K_0|), 2 points (w2_N22_s8 w2_N33_s8) [aggregate.py]
\newcommand{\NumSlopeEOneWTwoSEightFullIdAmort}{1.16}
% number of points of \NumSlopeEOneWTwoSEightFullIdAmort [aggregate.py]
\newcommand{\NumSlopeEOneWTwoSEightFullIdAmortN}{2}
% R^2 of \NumSlopeEOneWTwoSEightFullIdAmort [aggregate.py]
\newcommand{\NumSlopeEOneWTwoSEightFullIdAmortRSq}{1.000}
% E1 WTwoSEight full-id: least-squares slope of log10(median latency) vs log10(|K_0|), 2 points (w2_N22_s8 w2_N33_s8) [aggregate.py]
\newcommand{\NumSlopeEOneWTwoSEightFullIdMed}{1.17}
% number of points of \NumSlopeEOneWTwoSEightFullIdMed [aggregate.py]
\newcommand{\NumSlopeEOneWTwoSEightFullIdMedN}{2}
% R^2 of \NumSlopeEOneWTwoSEightFullIdMed [aggregate.py]
\newcommand{\NumSlopeEOneWTwoSEightFullIdMedRSq}{1.000}
% E1 WTwoSEight full-morse: least-squares slope of log10(amortized latency) vs log10(|K_0|), 5 points (w2_N7_s8 w2_N10_s8 w2_N15_s8 w2_N22_s8 w2_N33_s8) [aggregate.py]
\newcommand{\NumSlopeEOneWTwoSEightFullMorseAmort}{1.14}
% number of points of \NumSlopeEOneWTwoSEightFullMorseAmort [aggregate.py]
\newcommand{\NumSlopeEOneWTwoSEightFullMorseAmortN}{5}
% R^2 of \NumSlopeEOneWTwoSEightFullMorseAmort [aggregate.py]
\newcommand{\NumSlopeEOneWTwoSEightFullMorseAmortRSq}{0.994}
% E1 WTwoSEight full-morse: least-squares slope of log10(median latency) vs log10(|K_0|), 5 points (w2_N7_s8 w2_N10_s8 w2_N15_s8 w2_N22_s8 w2_N33_s8) [aggregate.py]
\newcommand{\NumSlopeEOneWTwoSEightFullMorseMed}{1.14}
% number of points of \NumSlopeEOneWTwoSEightFullMorseMed [aggregate.py]
\newcommand{\NumSlopeEOneWTwoSEightFullMorseMedN}{5}
% R^2 of \NumSlopeEOneWTwoSEightFullMorseMed [aggregate.py]
\newcommand{\NumSlopeEOneWTwoSEightFullMorseMedRSq}{0.994}
% E1 WTwoSEight gudhi: least-squares slope of log10(amortized latency) vs log10(|K_0|), 5 points (w2_N7_s8 w2_N10_s8 w2_N15_s8 w2_N22_s8 w2_N33_s8) [aggregate.py]
\newcommand{\NumSlopeEOneWTwoSEightGudhiAmort}{1.05}
% number of points of \NumSlopeEOneWTwoSEightGudhiAmort [aggregate.py]
\newcommand{\NumSlopeEOneWTwoSEightGudhiAmortN}{5}
% R^2 of \NumSlopeEOneWTwoSEightGudhiAmort [aggregate.py]
\newcommand{\NumSlopeEOneWTwoSEightGudhiAmortRSq}{0.999}
% E1 WTwoSEight gudhi: least-squares slope of log10(median latency) vs log10(|K_0|), 5 points (w2_N7_s8 w2_N10_s8 w2_N15_s8 w2_N22_s8 w2_N33_s8) [aggregate.py]
\newcommand{\NumSlopeEOneWTwoSEightGudhiMed}{1.06}
% number of points of \NumSlopeEOneWTwoSEightGudhiMed [aggregate.py]
\newcommand{\NumSlopeEOneWTwoSEightGudhiMedN}{5}
% R^2 of \NumSlopeEOneWTwoSEightGudhiMed [aggregate.py]
\newcommand{\NumSlopeEOneWTwoSEightGudhiMedRSq}{0.999}
% E6 WFiveCycle gudhi: least-squares slope of log10(amortized latency) vs log10(L), 4 points (w5a_L30000 w5a_L100000 w5a_L300000 w5a_L1000000) [aggregate.py]
\newcommand{\NumSlopeESixWFiveCycleGudhiAmort}{1.13}
% number of points of \NumSlopeESixWFiveCycleGudhiAmort [aggregate.py]
\newcommand{\NumSlopeESixWFiveCycleGudhiAmortN}{4}
% R^2 of \NumSlopeESixWFiveCycleGudhiAmort [aggregate.py]
\newcommand{\NumSlopeESixWFiveCycleGudhiAmortRSq}{1.000}
% E6 WFiveCycle gudhi: least-squares slope of log10(median latency) vs log10(L), 4 points (w5a_L30000 w5a_L100000 w5a_L300000 w5a_L1000000) [aggregate.py]
\newcommand{\NumSlopeESixWFiveCycleGudhiMed}{1.13}
% number of points of \NumSlopeESixWFiveCycleGudhiMed [aggregate.py]
\newcommand{\NumSlopeESixWFiveCycleGudhiMedN}{4}
% R^2 of \NumSlopeESixWFiveCycleGudhiMed [aggregate.py]
\newcommand{\NumSlopeESixWFiveCycleGudhiMedRSq}{1.000}
% E6 WFiveSphere full-id: least-squares slope of log10(amortized latency) vs log10(n_2), 2 points (w5b_L7 w5b_L8) [aggregate.py]
\newcommand{\NumSlopeESixWFiveSphereFullIdAmort}{1.01}
% number of points of \NumSlopeESixWFiveSphereFullIdAmort [aggregate.py]
\newcommand{\NumSlopeESixWFiveSphereFullIdAmortN}{2}
% R^2 of \NumSlopeESixWFiveSphereFullIdAmort [aggregate.py]
\newcommand{\NumSlopeESixWFiveSphereFullIdAmortRSq}{1.000}
% E6 WFiveSphere full-id: least-squares slope of log10(median latency) vs log10(n_2), 2 points (w5b_L7 w5b_L8) [aggregate.py]
\newcommand{\NumSlopeESixWFiveSphereFullIdMed}{1.02}
% number of points of \NumSlopeESixWFiveSphereFullIdMed [aggregate.py]
\newcommand{\NumSlopeESixWFiveSphereFullIdMedN}{2}
% R^2 of \NumSlopeESixWFiveSphereFullIdMed [aggregate.py]
\newcommand{\NumSlopeESixWFiveSphereFullIdMedRSq}{1.000}
% E6 WFiveSphere gudhi: least-squares slope of log10(amortized latency) vs log10(n_2), 4 points (w5b_L5 w5b_L6 w5b_L7 w5b_L8) [aggregate.py]
\newcommand{\NumSlopeESixWFiveSphereGudhiAmort}{1.13}
% number of points of \NumSlopeESixWFiveSphereGudhiAmort [aggregate.py]
\newcommand{\NumSlopeESixWFiveSphereGudhiAmortN}{4}
% R^2 of \NumSlopeESixWFiveSphereGudhiAmort [aggregate.py]
\newcommand{\NumSlopeESixWFiveSphereGudhiAmortRSq}{0.999}
% E6 WFiveSphere gudhi: least-squares slope of log10(median latency) vs log10(n_2), 4 points (w5b_L5 w5b_L6 w5b_L7 w5b_L8) [aggregate.py]
\newcommand{\NumSlopeESixWFiveSphereGudhiMed}{1.13}
% number of points of \NumSlopeESixWFiveSphereGudhiMed [aggregate.py]
\newcommand{\NumSlopeESixWFiveSphereGudhiMedN}{4}
% R^2 of \NumSlopeESixWFiveSphereGudhiMed [aggregate.py]
\newcommand{\NumSlopeESixWFiveSphereGudhiMedRSq}{0.999}
% E3 WThree drb-id: least-squares slope of log10(amortized latency) vs log10(|K_0|), 2 points (w3_N20_k2 w3_N27_k2) [aggregate.py]
\newcommand{\NumSlopeEThreeWThreeDrbIdAmort}{-0.39}
% number of points of \NumSlopeEThreeWThreeDrbIdAmort [aggregate.py]
\newcommand{\NumSlopeEThreeWThreeDrbIdAmortN}{2}
% R^2 of \NumSlopeEThreeWThreeDrbIdAmort [aggregate.py]
\newcommand{\NumSlopeEThreeWThreeDrbIdAmortRSq}{1.000}
% E3 WThree drb-id: least-squares slope of log10(median latency) vs log10(|K_0|), 2 points (w3_N20_k2 w3_N27_k2) [aggregate.py]
\newcommand{\NumSlopeEThreeWThreeDrbIdMed}{-0.28}
% number of points of \NumSlopeEThreeWThreeDrbIdMed [aggregate.py]
\newcommand{\NumSlopeEThreeWThreeDrbIdMedN}{2}
% R^2 of \NumSlopeEThreeWThreeDrbIdMed [aggregate.py]
\newcommand{\NumSlopeEThreeWThreeDrbIdMedRSq}{1.000}
% E3 WThree drb-id-rb: least-squares slope of log10(amortized latency) vs log10(|K_0|), 2 points (w3_N20_k2 w3_N27_k2) [aggregate.py]
\newcommand{\NumSlopeEThreeWThreeDrbIdRbAmort}{0.17}
% number of points of \NumSlopeEThreeWThreeDrbIdRbAmort [aggregate.py]
\newcommand{\NumSlopeEThreeWThreeDrbIdRbAmortN}{2}
% R^2 of \NumSlopeEThreeWThreeDrbIdRbAmort [aggregate.py]
\newcommand{\NumSlopeEThreeWThreeDrbIdRbAmortRSq}{1.000}
% E3 WThree drb-id-rb: least-squares slope of log10(median latency) vs log10(|K_0|), 2 points (w3_N20_k2 w3_N27_k2) [aggregate.py]
\newcommand{\NumSlopeEThreeWThreeDrbIdRbMed}{0.12}
% number of points of \NumSlopeEThreeWThreeDrbIdRbMed [aggregate.py]
\newcommand{\NumSlopeEThreeWThreeDrbIdRbMedN}{2}
% R^2 of \NumSlopeEThreeWThreeDrbIdRbMed [aggregate.py]
\newcommand{\NumSlopeEThreeWThreeDrbIdRbMedRSq}{1.000}
% E3 WThree full-id: least-squares slope of log10(amortized latency) vs log10(|K_0|), 3 points (w3_N14_k2 w3_N20_k2 w3_N27_k2) [aggregate.py]
\newcommand{\NumSlopeEThreeWThreeFullIdAmort}{1.03}
% number of points of \NumSlopeEThreeWThreeFullIdAmort [aggregate.py]
\newcommand{\NumSlopeEThreeWThreeFullIdAmortN}{3}
% R^2 of \NumSlopeEThreeWThreeFullIdAmort [aggregate.py]
\newcommand{\NumSlopeEThreeWThreeFullIdAmortRSq}{0.999}
% E3 WThree full-id: least-squares slope of log10(median latency) vs log10(|K_0|), 3 points (w3_N14_k2 w3_N20_k2 w3_N27_k2) [aggregate.py]
\newcommand{\NumSlopeEThreeWThreeFullIdMed}{1.05}
% number of points of \NumSlopeEThreeWThreeFullIdMed [aggregate.py]
\newcommand{\NumSlopeEThreeWThreeFullIdMedN}{3}
% R^2 of \NumSlopeEThreeWThreeFullIdMed [aggregate.py]
\newcommand{\NumSlopeEThreeWThreeFullIdMedRSq}{0.999}
% E3 WThree full-morse: least-squares slope of log10(amortized latency) vs log10(|K_0|), 3 points (w3_N14_k2 w3_N20_k2 w3_N27_k2) [aggregate.py]
\newcommand{\NumSlopeEThreeWThreeFullMorseAmort}{1.14}
% number of points of \NumSlopeEThreeWThreeFullMorseAmort [aggregate.py]
\newcommand{\NumSlopeEThreeWThreeFullMorseAmortN}{3}
% R^2 of \NumSlopeEThreeWThreeFullMorseAmort [aggregate.py]
\newcommand{\NumSlopeEThreeWThreeFullMorseAmortRSq}{0.998}
% E3 WThree full-morse: least-squares slope of log10(median latency) vs log10(|K_0|), 3 points (w3_N14_k2 w3_N20_k2 w3_N27_k2) [aggregate.py]
\newcommand{\NumSlopeEThreeWThreeFullMorseMed}{1.14}
% number of points of \NumSlopeEThreeWThreeFullMorseMed [aggregate.py]
\newcommand{\NumSlopeEThreeWThreeFullMorseMedN}{3}
% R^2 of \NumSlopeEThreeWThreeFullMorseMed [aggregate.py]
\newcommand{\NumSlopeEThreeWThreeFullMorseMedRSq}{0.998}
% E3 WThree gudhi: least-squares slope of log10(amortized latency) vs log10(|K_0|), 3 points (w3_N14_k2 w3_N20_k2 w3_N27_k2) [aggregate.py]
\newcommand{\NumSlopeEThreeWThreeGudhiAmort}{0.92}
% number of points of \NumSlopeEThreeWThreeGudhiAmort [aggregate.py]
\newcommand{\NumSlopeEThreeWThreeGudhiAmortN}{3}
% R^2 of \NumSlopeEThreeWThreeGudhiAmort [aggregate.py]
\newcommand{\NumSlopeEThreeWThreeGudhiAmortRSq}{0.983}
% E3 WThree gudhi: least-squares slope of log10(median latency) vs log10(|K_0|), 3 points (w3_N14_k2 w3_N20_k2 w3_N27_k2) [aggregate.py]
\newcommand{\NumSlopeEThreeWThreeGudhiMed}{0.93}
% number of points of \NumSlopeEThreeWThreeGudhiMed [aggregate.py]
\newcommand{\NumSlopeEThreeWThreeGudhiMedN}{3}
% R^2 of \NumSlopeEThreeWThreeGudhiMed [aggregate.py]
\newcommand{\NumSlopeEThreeWThreeGudhiMedRSq}{0.981}
% E3 WThree mmhm: least-squares slope of log10(amortized latency) vs log10(|K_0|), 2 points (w3_N20_k2 w3_N27_k2) [aggregate.py]
\newcommand{\NumSlopeEThreeWThreeMmhmAmort}{0.17}
% number of points of \NumSlopeEThreeWThreeMmhmAmort [aggregate.py]
\newcommand{\NumSlopeEThreeWThreeMmhmAmortN}{2}
% R^2 of \NumSlopeEThreeWThreeMmhmAmort [aggregate.py]
\newcommand{\NumSlopeEThreeWThreeMmhmAmortRSq}{1.000}
% E3 WThree mmhm: least-squares slope of log10(median latency) vs log10(|K_0|), 2 points (w3_N20_k2 w3_N27_k2) [aggregate.py]
\newcommand{\NumSlopeEThreeWThreeMmhmMed}{0.14}
% number of points of \NumSlopeEThreeWThreeMmhmMed [aggregate.py]
\newcommand{\NumSlopeEThreeWThreeMmhmMedN}{2}
% R^2 of \NumSlopeEThreeWThreeMmhmMed [aggregate.py]
\newcommand{\NumSlopeEThreeWThreeMmhmMedRSq}{1.000}
% E3 WThree mmhm-norb: least-squares slope of log10(amortized latency) vs log10(|K_0|), 2 points (w3_N20_k2 w3_N27_k2) [aggregate.py]
\newcommand{\NumSlopeEThreeWThreeMmhmNorbAmort}{-0.30}
% number of points of \NumSlopeEThreeWThreeMmhmNorbAmort [aggregate.py]
\newcommand{\NumSlopeEThreeWThreeMmhmNorbAmortN}{2}
% R^2 of \NumSlopeEThreeWThreeMmhmNorbAmort [aggregate.py]
\newcommand{\NumSlopeEThreeWThreeMmhmNorbAmortRSq}{1.000}
% E3 WThree mmhm-norb: least-squares slope of log10(median latency) vs log10(|K_0|), 2 points (w3_N20_k2 w3_N27_k2) [aggregate.py]
\newcommand{\NumSlopeEThreeWThreeMmhmNorbMed}{-0.23}
% number of points of \NumSlopeEThreeWThreeMmhmNorbMed [aggregate.py]
\newcommand{\NumSlopeEThreeWThreeMmhmNorbMedN}{2}
% R^2 of \NumSlopeEThreeWThreeMmhmNorbMed [aggregate.py]
\newcommand{\NumSlopeEThreeWThreeMmhmNorbMedRSq}{1.000}
% E2 WTwoNTwentyTwo full-id: least-squares slope of log10(amortized latency) vs log10(s), 4 points (w2_N22_s8 w2_N22_s16 w2_N22_s32 w2_N22_s64) [aggregate.py]
\newcommand{\NumSlopeETwoWTwoNTwentyTwoFullIdAmort}{0.12}
% number of points of \NumSlopeETwoWTwoNTwentyTwoFullIdAmort [aggregate.py]
\newcommand{\NumSlopeETwoWTwoNTwentyTwoFullIdAmortN}{4}
% R^2 of \NumSlopeETwoWTwoNTwentyTwoFullIdAmort [aggregate.py]
\newcommand{\NumSlopeETwoWTwoNTwentyTwoFullIdAmortRSq}{0.893}
% E2 WTwoNTwentyTwo full-id: least-squares slope of log10(median latency) vs log10(s), 4 points (w2_N22_s8 w2_N22_s16 w2_N22_s32 w2_N22_s64) [aggregate.py]
\newcommand{\NumSlopeETwoWTwoNTwentyTwoFullIdMed}{0.12}
% number of points of \NumSlopeETwoWTwoNTwentyTwoFullIdMed [aggregate.py]
\newcommand{\NumSlopeETwoWTwoNTwentyTwoFullIdMedN}{4}
% R^2 of \NumSlopeETwoWTwoNTwentyTwoFullIdMed [aggregate.py]
\newcommand{\NumSlopeETwoWTwoNTwentyTwoFullIdMedRSq}{0.920}
% E2 WTwoNTwentyTwo full-morse: least-squares slope of log10(amortized latency) vs log10(s), 7 points (w2_N22_s1 w2_N22_s2 w2_N22_s4 w2_N22_s8 w2_N22_s16 w2_N22_s32 w2_N22_s64) [aggregate.py]
\newcommand{\NumSlopeETwoWTwoNTwentyTwoFullMorseAmort}{0.07}
% number of points of \NumSlopeETwoWTwoNTwentyTwoFullMorseAmort [aggregate.py]
\newcommand{\NumSlopeETwoWTwoNTwentyTwoFullMorseAmortN}{7}
% R^2 of \NumSlopeETwoWTwoNTwentyTwoFullMorseAmort [aggregate.py]
\newcommand{\NumSlopeETwoWTwoNTwentyTwoFullMorseAmortRSq}{0.788}
% E2 WTwoNTwentyTwo full-morse: least-squares slope of log10(median latency) vs log10(s), 7 points (w2_N22_s1 w2_N22_s2 w2_N22_s4 w2_N22_s8 w2_N22_s16 w2_N22_s32 w2_N22_s64) [aggregate.py]
\newcommand{\NumSlopeETwoWTwoNTwentyTwoFullMorseMed}{0.07}
% number of points of \NumSlopeETwoWTwoNTwentyTwoFullMorseMed [aggregate.py]
\newcommand{\NumSlopeETwoWTwoNTwentyTwoFullMorseMedN}{7}
% R^2 of \NumSlopeETwoWTwoNTwentyTwoFullMorseMed [aggregate.py]
\newcommand{\NumSlopeETwoWTwoNTwentyTwoFullMorseMedRSq}{0.782}
% timed runs (E1-E7) [aggregate.py]
\newcommand{\NumTimedRuns}{332}
% T: updates per timed run [configs/make_configs.py]
\newcommand{\NumUpdates}{2{,}000}
% (run, state) pairs of exact methods compared with the engine full-id Betti file of the same instance [verify.py]
\newcommand{\NumVerifiedStates}{354{,}396}
% (run, state) pairs compared with the analytic .expected.betti [verify.py]
\newcommand{\NumVerifiedStatesExpected}{567{,}708}
% (run, state) pairs compared with the GUDHI oracle [verify.py]
\newcommand{\NumVerifiedStatesGudhi}{333{,}837}
% w5a_L100000 / gudhi: amortized latency = total time incl. rebuilds / updates (us) [configs/E6.json, aggregate.py]
\newcommand{\NumWFiveCycleLOneHundredThousandGudhiAmort}{26{,}800}
% w5a_L100000 / gudhi: INIT build time (ms), median of seeds [configs/E6.json, aggregate.py]
\newcommand{\NumWFiveCycleLOneHundredThousandGudhiInitMs}{77.7}
% w5a_L100000 / gudhi: max latency (us) [configs/E6.json, aggregate.py]
\newcommand{\NumWFiveCycleLOneHundredThousandGudhiMax}{77{,}400}
% w5a_L100000 / gudhi: median per-update latency (us), 3 seeds pooled [configs/E6.json, aggregate.py]
\newcommand{\NumWFiveCycleLOneHundredThousandGudhiMed}{26{,}500}
% w5a_L100000 / gudhi: p95 latency (us) [configs/E6.json, aggregate.py]
\newcommand{\NumWFiveCycleLOneHundredThousandGudhiPNinetyFive}{33{,}000}
% w5a_L100000 / gudhi: p99 latency (us) [configs/E6.json, aggregate.py]
\newcommand{\NumWFiveCycleLOneHundredThousandGudhiPNinetyNine}{38{,}200}
% w5a_L100000 / gudhi: 25th percentile latency (us) [configs/E6.json, aggregate.py]
\newcommand{\NumWFiveCycleLOneHundredThousandGudhiQOne}{25{,}000}
% w5a_L100000 / gudhi: 75th percentile latency (us) [configs/E6.json, aggregate.py]
\newcommand{\NumWFiveCycleLOneHundredThousandGudhiQThree}{27{,}600}
% w5a_L100000 / gudhi: updates timed per run (T or prefix T') [configs/E6.json, aggregate.py]
\newcommand{\NumWFiveCycleLOneHundredThousandGudhiT}{2{,}000}
% w5a_L1000000 / full-id: amortized latency = total time incl. rebuilds / updates (us) [configs/E6.json, aggregate.py]
\newcommand{\NumWFiveCycleLOneMillionFullIdAmort}{36{,}300}
% w5a_L1000000 / full-id: INIT build time (ms), median of seeds [configs/E6.json, aggregate.py]
\newcommand{\NumWFiveCycleLOneMillionFullIdInitMs}{772}
% w5a_L1000000 / full-id: max latency (us) [configs/E6.json, aggregate.py]
\newcommand{\NumWFiveCycleLOneMillionFullIdMax}{139{,}000}
% w5a_L1000000 / full-id: median per-update latency (us), 3 seeds pooled [configs/E6.json, aggregate.py]
\newcommand{\NumWFiveCycleLOneMillionFullIdMed}{35{,}900}
% w5a_L1000000 / full-id: p95 latency (us) [configs/E6.json, aggregate.py]
\newcommand{\NumWFiveCycleLOneMillionFullIdPNinetyFive}{42{,}200}
% w5a_L1000000 / full-id: p99 latency (us) [configs/E6.json, aggregate.py]
\newcommand{\NumWFiveCycleLOneMillionFullIdPNinetyNine}{48{,}100}
% w5a_L1000000 / full-id: 25th percentile latency (us) [configs/E6.json, aggregate.py]
\newcommand{\NumWFiveCycleLOneMillionFullIdQOne}{34{,}100}
% w5a_L1000000 / full-id: 75th percentile latency (us) [configs/E6.json, aggregate.py]
\newcommand{\NumWFiveCycleLOneMillionFullIdQThree}{37{,}700}
% w5a_L1000000 / full-id: amortized latency of gudhi (on its T or T' prefix) / this method [configs/E6.json, aggregate.py]
\newcommand{\NumWFiveCycleLOneMillionFullIdSpeedupGudhi}{10.4}
% w5a_L1000000 / full-id: updates timed per run (T or prefix T') [configs/E6.json, aggregate.py]
\newcommand{\NumWFiveCycleLOneMillionFullIdT}{2{,}000}
% w5a_L1000000 / gudhi: amortized latency = total time incl. rebuilds / updates (us) [configs/E6.json, aggregate.py]
\newcommand{\NumWFiveCycleLOneMillionGudhiAmort}{376{,}000}
% w5a_L1000000 / gudhi: INIT build time (ms), median of seeds [configs/E6.json, aggregate.py]
\newcommand{\NumWFiveCycleLOneMillionGudhiInitMs}{910}
% w5a_L1000000 / gudhi: max latency (us) [configs/E6.json, aggregate.py]
\newcommand{\NumWFiveCycleLOneMillionGudhiMax}{498{,}000}
% w5a_L1000000 / gudhi: median per-update latency (us), 3 seeds pooled [configs/E6.json, aggregate.py]
\newcommand{\NumWFiveCycleLOneMillionGudhiMed}{372{,}000}
% w5a_L1000000 / gudhi: p95 latency (us) [configs/E6.json, aggregate.py]
\newcommand{\NumWFiveCycleLOneMillionGudhiPNinetyFive}{429{,}000}
% w5a_L1000000 / gudhi: p99 latency (us) [configs/E6.json, aggregate.py]
\newcommand{\NumWFiveCycleLOneMillionGudhiPNinetyNine}{457{,}000}
% w5a_L1000000 / gudhi: 25th percentile latency (us) [configs/E6.json, aggregate.py]
\newcommand{\NumWFiveCycleLOneMillionGudhiQOne}{353{,}000}
% w5a_L1000000 / gudhi: 75th percentile latency (us) [configs/E6.json, aggregate.py]
\newcommand{\NumWFiveCycleLOneMillionGudhiQThree}{395{,}000}
% w5a_L1000000 / gudhi: amortized latency of full-id / this method [configs/E6.json, aggregate.py]
\newcommand{\NumWFiveCycleLOneMillionGudhiSpeedupFullId}{0.0964}
% w5a_L1000000 / gudhi: median latency of full-id / this method [configs/E6.json, aggregate.py]
\newcommand{\NumWFiveCycleLOneMillionGudhiSpeedupFullIdMed}{0.0964}
% w5a_L1000000 / gudhi: updates timed per run (T or prefix T') [configs/E6.json, aggregate.py]
\newcommand{\NumWFiveCycleLOneMillionGudhiT}{200}
% w5a_L1000000: mean atomic edits per update [configs/E6.json, aggregate.py]
\newcommand{\NumWFiveCycleLOneMillionOps}{1.00}
% w5a_L1000000: |K_0| (mean of seeds) [configs/E6.json, aggregate.py]
\newcommand{\NumWFiveCycleLOneMillionSize}{2{,}000{,}000}
% w5a_L30000 / gudhi: amortized latency = total time incl. rebuilds / updates (us) [configs/E6.json, aggregate.py]
\newcommand{\NumWFiveCycleLThirtyThousandGudhiAmort}{7{,}180}
% w5a_L30000 / gudhi: INIT build time (ms), median of seeds [configs/E6.json, aggregate.py]
\newcommand{\NumWFiveCycleLThirtyThousandGudhiInitMs}{20.2}
% w5a_L30000 / gudhi: max latency (us) [configs/E6.json, aggregate.py]
\newcommand{\NumWFiveCycleLThirtyThousandGudhiMax}{37{,}000}
% w5a_L30000 / gudhi: median per-update latency (us), 3 seeds pooled [configs/E6.json, aggregate.py]
\newcommand{\NumWFiveCycleLThirtyThousandGudhiMed}{7{,}140}
% w5a_L30000 / gudhi: p95 latency (us) [configs/E6.json, aggregate.py]
\newcommand{\NumWFiveCycleLThirtyThousandGudhiPNinetyFive}{8{,}950}
% w5a_L30000 / gudhi: p99 latency (us) [configs/E6.json, aggregate.py]
\newcommand{\NumWFiveCycleLThirtyThousandGudhiPNinetyNine}{10{,}600}
% w5a_L30000 / gudhi: 25th percentile latency (us) [configs/E6.json, aggregate.py]
\newcommand{\NumWFiveCycleLThirtyThousandGudhiQOne}{6{,}760}
% w5a_L30000 / gudhi: 75th percentile latency (us) [configs/E6.json, aggregate.py]
\newcommand{\NumWFiveCycleLThirtyThousandGudhiQThree}{7{,}400}
% w5a_L30000 / gudhi: updates timed per run (T or prefix T') [configs/E6.json, aggregate.py]
\newcommand{\NumWFiveCycleLThirtyThousandGudhiT}{2{,}000}
% w5a_L300000 / gudhi: amortized latency = total time incl. rebuilds / updates (us) [configs/E6.json, aggregate.py]
\newcommand{\NumWFiveCycleLThreeHundredThousandGudhiAmort}{91{,}200}
% w5a_L300000 / gudhi: INIT build time (ms), median of seeds [configs/E6.json, aggregate.py]
\newcommand{\NumWFiveCycleLThreeHundredThousandGudhiInitMs}{233}
% w5a_L300000 / gudhi: max latency (us) [configs/E6.json, aggregate.py]
\newcommand{\NumWFiveCycleLThreeHundredThousandGudhiMax}{129{,}000}
% w5a_L300000 / gudhi: median per-update latency (us), 3 seeds pooled [configs/E6.json, aggregate.py]
\newcommand{\NumWFiveCycleLThreeHundredThousandGudhiMed}{90{,}800}
% w5a_L300000 / gudhi: p95 latency (us) [configs/E6.json, aggregate.py]
\newcommand{\NumWFiveCycleLThreeHundredThousandGudhiPNinetyFive}{108{,}000}
% w5a_L300000 / gudhi: p99 latency (us) [configs/E6.json, aggregate.py]
\newcommand{\NumWFiveCycleLThreeHundredThousandGudhiPNinetyNine}{121{,}000}
% w5a_L300000 / gudhi: 25th percentile latency (us) [configs/E6.json, aggregate.py]
\newcommand{\NumWFiveCycleLThreeHundredThousandGudhiQOne}{85{,}400}
% w5a_L300000 / gudhi: 75th percentile latency (us) [configs/E6.json, aggregate.py]
\newcommand{\NumWFiveCycleLThreeHundredThousandGudhiQThree}{94{,}900}
% w5a_L300000 / gudhi: updates timed per run (T or prefix T') [configs/E6.json, aggregate.py]
\newcommand{\NumWFiveCycleLThreeHundredThousandGudhiT}{200}
% w5b_L8 / drb-id: mean column additions per update (non-rebuild updates) [configs/E6.json, aggregate.py]
\newcommand{\NumWFiveSphereLEightDrbIdAdds}{0.501}
% w5b_L8 / drb-id: amortized latency = total time incl. rebuilds / updates (us) [configs/E6.json, aggregate.py]
\newcommand{\NumWFiveSphereLEightDrbIdAmort}{21{,}500}
% w5b_L8 / drb-id: mean footprint_sum per update (sum of |occ| over deletions) [configs/E6.json, aggregate.py]
\newcommand{\NumWFiveSphereLEightDrbIdFootprint}{0.501}
% w5b_L8 / drb-id: INIT build time (ms), median of seeds [configs/E6.json, aggregate.py]
\newcommand{\NumWFiveSphereLEightDrbIdInitMs}{2{,}920}
% w5b_L8 / drb-id: max latency (us) [configs/E6.json, aggregate.py]
\newcommand{\NumWFiveSphereLEightDrbIdMax}{299{,}000}
% w5b_L8 / drb-id: median per-update latency (us), 3 seeds pooled [configs/E6.json, aggregate.py]
\newcommand{\NumWFiveSphereLEightDrbIdMed}{15{,}600}
% w5b_L8 / drb-id: nnz(R)+nnz(V) after the last update (mean of seeds) [configs/E6.json, aggregate.py]
\newcommand{\NumWFiveSphereLEightDrbIdNnzEnd}{14{,}586{,}951}
% w5b_L8 / drb-id: nnz at the end / nnz after INIT [configs/E6.json, aggregate.py]
\newcommand{\NumWFiveSphereLEightDrbIdNnzGrowth}{1.13}
% w5b_L8 / drb-id: max nnz(R)+nnz(V) over the run (mean of seeds) [configs/E6.json, aggregate.py]
\newcommand{\NumWFiveSphereLEightDrbIdNnzMax}{14{,}586{,}951}
% w5b_L8 / drb-id: nnz(R)+nnz(V) after INIT (mean of seeds) [configs/E6.json, aggregate.py]
\newcommand{\NumWFiveSphereLEightDrbIdNnzStart}{12{,}926{,}027}
% w5b_L8 / drb-id: p95 latency (us) [configs/E6.json, aggregate.py]
\newcommand{\NumWFiveSphereLEightDrbIdPNinetyFive}{110{,}000}
% w5b_L8 / drb-id: p99 latency (us) [configs/E6.json, aggregate.py]
\newcommand{\NumWFiveSphereLEightDrbIdPNinetyNine}{203{,}000}
% w5b_L8 / drb-id: 25th percentile latency (us) [configs/E6.json, aggregate.py]
\newcommand{\NumWFiveSphereLEightDrbIdQOne}{1{,}860}
% w5b_L8 / drb-id: 75th percentile latency (us) [configs/E6.json, aggregate.py]
\newcommand{\NumWFiveSphereLEightDrbIdQThree}{23{,}800}
% w5b_L8 / drb-id-rb: mean column additions per update (non-rebuild updates) [configs/E6.json, aggregate.py]
\newcommand{\NumWFiveSphereLEightDrbIdRbAdds}{0.501}
% w5b_L8 / drb-id-rb: amortized latency = total time incl. rebuilds / updates (us) [configs/E6.json, aggregate.py]
\newcommand{\NumWFiveSphereLEightDrbIdRbAmort}{21{,}700}
% w5b_L8 / drb-id-rb: mean footprint_sum per update (sum of |occ| over deletions) [configs/E6.json, aggregate.py]
\newcommand{\NumWFiveSphereLEightDrbIdRbFootprint}{0.501}
% w5b_L8 / drb-id-rb: INIT build time (ms), median of seeds [configs/E6.json, aggregate.py]
\newcommand{\NumWFiveSphereLEightDrbIdRbInitMs}{2{,}900}
% w5b_L8 / drb-id-rb: max latency (us) [configs/E6.json, aggregate.py]
\newcommand{\NumWFiveSphereLEightDrbIdRbMax}{343{,}000}
% w5b_L8 / drb-id-rb: median per-update latency (us), 3 seeds pooled [configs/E6.json, aggregate.py]
\newcommand{\NumWFiveSphereLEightDrbIdRbMed}{16{,}600}
% w5b_L8 / drb-id-rb: nnz(R)+nnz(V) after the last update (mean of seeds) [configs/E6.json, aggregate.py]
\newcommand{\NumWFiveSphereLEightDrbIdRbNnzEnd}{14{,}586{,}951}
% w5b_L8 / drb-id-rb: nnz at the end / nnz after INIT [configs/E6.json, aggregate.py]
\newcommand{\NumWFiveSphereLEightDrbIdRbNnzGrowth}{1.13}
% w5b_L8 / drb-id-rb: max nnz(R)+nnz(V) over the run (mean of seeds) [configs/E6.json, aggregate.py]
\newcommand{\NumWFiveSphereLEightDrbIdRbNnzMax}{14{,}586{,}951}
% w5b_L8 / drb-id-rb: nnz(R)+nnz(V) after INIT (mean of seeds) [configs/E6.json, aggregate.py]
\newcommand{\NumWFiveSphereLEightDrbIdRbNnzStart}{12{,}926{,}027}
% w5b_L8 / drb-id-rb: p95 latency (us) [configs/E6.json, aggregate.py]
\newcommand{\NumWFiveSphereLEightDrbIdRbPNinetyFive}{112{,}000}
% w5b_L8 / drb-id-rb: p99 latency (us) [configs/E6.json, aggregate.py]
\newcommand{\NumWFiveSphereLEightDrbIdRbPNinetyNine}{206{,}000}
% w5b_L8 / drb-id-rb: 25th percentile latency (us) [configs/E6.json, aggregate.py]
\newcommand{\NumWFiveSphereLEightDrbIdRbQOne}{1{,}840}
% w5b_L8 / drb-id-rb: 75th percentile latency (us) [configs/E6.json, aggregate.py]
\newcommand{\NumWFiveSphereLEightDrbIdRbQThree}{23{,}800}
% w5b_L8 / drb-id-rb: rebuilds per run (mean of seeds), default policy alpha=1, period off [configs/E6.json, aggregate.py]
\newcommand{\NumWFiveSphereLEightDrbIdRbRebuilds}{0}
% w5b_L8 / drb-id-rb: amortized latency of full-id / this method [configs/E6.json, aggregate.py]
\newcommand{\NumWFiveSphereLEightDrbIdRbSpeedupFullId}{8.18}
% w5b_L8 / drb-id-rb: median latency of full-id / this method [configs/E6.json, aggregate.py]
\newcommand{\NumWFiveSphereLEightDrbIdRbSpeedupFullIdMed}{11.4}
% w5b_L8 / drb-id-rb: amortized latency of gudhi (on its T or T' prefix) / this method [configs/E6.json, aggregate.py]
\newcommand{\NumWFiveSphereLEightDrbIdRbSpeedupGudhi}{66.9}
% w5b_L8 / drb-id-rb: updates timed per run (T or prefix T') [configs/E6.json, aggregate.py]
\newcommand{\NumWFiveSphereLEightDrbIdRbT}{2{,}000}
% w5b_L8 / drb-id: amortized latency of full-id / this method [configs/E6.json, aggregate.py]
\newcommand{\NumWFiveSphereLEightDrbIdSpeedupFullId}{8.27}
% w5b_L8 / drb-id: median latency of full-id / this method [configs/E6.json, aggregate.py]
\newcommand{\NumWFiveSphereLEightDrbIdSpeedupFullIdMed}{12.1}
% w5b_L8 / drb-id: amortized latency of gudhi (on its T or T' prefix) / this method [configs/E6.json, aggregate.py]
\newcommand{\NumWFiveSphereLEightDrbIdSpeedupGudhi}{67.7}
% w5b_L8 / drb-id: updates timed per run (T or prefix T') [configs/E6.json, aggregate.py]
\newcommand{\NumWFiveSphereLEightDrbIdT}{2{,}000}
% w5b_L8 / euler: amortized latency = total time incl. rebuilds / updates (us) [configs/E6.json, aggregate.py]
\newcommand{\NumWFiveSphereLEightEulerAmort}{0.267}
% w5b_L8 / euler: INIT build time (ms), median of seeds [configs/E6.json, aggregate.py]
\newcommand{\NumWFiveSphereLEightEulerInitMs}{2{,}440}
% w5b_L8 / euler: max latency (us) [configs/E6.json, aggregate.py]
\newcommand{\NumWFiveSphereLEightEulerMax}{36.3}
% w5b_L8 / euler: median per-update latency (us), 3 seeds pooled [configs/E6.json, aggregate.py]
\newcommand{\NumWFiveSphereLEightEulerMed}{0.108}
% w5b_L8 / euler: p95 latency (us) [configs/E6.json, aggregate.py]
\newcommand{\NumWFiveSphereLEightEulerPNinetyFive}{0.263}
% w5b_L8 / euler: p99 latency (us) [configs/E6.json, aggregate.py]
\newcommand{\NumWFiveSphereLEightEulerPNinetyNine}{4.07}
% w5b_L8 / euler: 25th percentile latency (us) [configs/E6.json, aggregate.py]
\newcommand{\NumWFiveSphereLEightEulerQOne}{0.0910}
% w5b_L8 / euler: 75th percentile latency (us) [configs/E6.json, aggregate.py]
\newcommand{\NumWFiveSphereLEightEulerQThree}{0.124}
% w5b_L8 / euler: amortized latency of full-id / this method [configs/E6.json, aggregate.py]
\newcommand{\NumWFiveSphereLEightEulerSpeedupFullId}{665{,}000}
% w5b_L8 / euler: median latency of full-id / this method [configs/E6.json, aggregate.py]
\newcommand{\NumWFiveSphereLEightEulerSpeedupFullIdMed}{1{,}740{,}000}
% w5b_L8 / euler: amortized latency of gudhi (on its T or T' prefix) / this method [configs/E6.json, aggregate.py]
\newcommand{\NumWFiveSphereLEightEulerSpeedupGudhi}{5{,}440{,}000}
% w5b_L8 / euler: updates timed per run (T or prefix T') [configs/E6.json, aggregate.py]
\newcommand{\NumWFiveSphereLEightEulerT}{2{,}000}
% w5b_L8 / full-id: amortized latency = total time incl. rebuilds / updates (us) [configs/E6.json, aggregate.py]
\newcommand{\NumWFiveSphereLEightFullIdAmort}{178{,}000}
% w5b_L8 / full-id: INIT build time (ms), median of seeds [configs/E6.json, aggregate.py]
\newcommand{\NumWFiveSphereLEightFullIdInitMs}{2{,}650}
% w5b_L8 / full-id: max latency (us) [configs/E6.json, aggregate.py]
\newcommand{\NumWFiveSphereLEightFullIdMax}{336{,}000}
% w5b_L8 / full-id: median per-update latency (us), 3 seeds pooled [configs/E6.json, aggregate.py]
\newcommand{\NumWFiveSphereLEightFullIdMed}{188{,}000}
% w5b_L8 / full-id: p95 latency (us) [configs/E6.json, aggregate.py]
\newcommand{\NumWFiveSphereLEightFullIdPNinetyFive}{218{,}000}
% w5b_L8 / full-id: p99 latency (us) [configs/E6.json, aggregate.py]
\newcommand{\NumWFiveSphereLEightFullIdPNinetyNine}{240{,}000}
% w5b_L8 / full-id: 25th percentile latency (us) [configs/E6.json, aggregate.py]
\newcommand{\NumWFiveSphereLEightFullIdQOne}{163{,}000}
% w5b_L8 / full-id: 75th percentile latency (us) [configs/E6.json, aggregate.py]
\newcommand{\NumWFiveSphereLEightFullIdQThree}{199{,}000}
% w5b_L8 / full-id: amortized latency of gudhi (on its T or T' prefix) / this method [configs/E6.json, aggregate.py]
\newcommand{\NumWFiveSphereLEightFullIdSpeedupGudhi}{8.18}
% w5b_L8 / full-id: updates timed per run (T or prefix T') [configs/E6.json, aggregate.py]
\newcommand{\NumWFiveSphereLEightFullIdT}{2{,}000}
% w5b_L8 / gudhi: amortized latency = total time incl. rebuilds / updates (us) [configs/E6.json, aggregate.py]
\newcommand{\NumWFiveSphereLEightGudhiAmort}{1{,}460{,}000}
% w5b_L8 / gudhi: INIT build time (ms), median of seeds [configs/E6.json, aggregate.py]
\newcommand{\NumWFiveSphereLEightGudhiInitMs}{2{,}990}
% w5b_L8 / gudhi: max latency (us) [configs/E6.json, aggregate.py]
\newcommand{\NumWFiveSphereLEightGudhiMax}{1{,}710{,}000}
% w5b_L8 / gudhi: median per-update latency (us), 3 seeds pooled [configs/E6.json, aggregate.py]
\newcommand{\NumWFiveSphereLEightGudhiMed}{1{,}460{,}000}
% w5b_L8 / gudhi: p95 latency (us) [configs/E6.json, aggregate.py]
\newcommand{\NumWFiveSphereLEightGudhiPNinetyFive}{1{,}580{,}000}
% w5b_L8 / gudhi: p99 latency (us) [configs/E6.json, aggregate.py]
\newcommand{\NumWFiveSphereLEightGudhiPNinetyNine}{1{,}650{,}000}
% w5b_L8 / gudhi: 25th percentile latency (us) [configs/E6.json, aggregate.py]
\newcommand{\NumWFiveSphereLEightGudhiQOne}{1{,}400{,}000}
% w5b_L8 / gudhi: 75th percentile latency (us) [configs/E6.json, aggregate.py]
\newcommand{\NumWFiveSphereLEightGudhiQThree}{1{,}510{,}000}
% w5b_L8 / gudhi: amortized latency of full-id / this method [configs/E6.json, aggregate.py]
\newcommand{\NumWFiveSphereLEightGudhiSpeedupFullId}{0.122}
% w5b_L8 / gudhi: median latency of full-id / this method [configs/E6.json, aggregate.py]
\newcommand{\NumWFiveSphereLEightGudhiSpeedupFullIdMed}{0.129}
% w5b_L8 / gudhi: updates timed per run (T or prefix T') [configs/E6.json, aggregate.py]
\newcommand{\NumWFiveSphereLEightGudhiT}{200}
% w5b_L8 / mmhm: mean column additions per update (non-rebuild updates) [configs/E6.json, aggregate.py]
\newcommand{\NumWFiveSphereLEightMmhmAdds}{0.500}
% w5b_L8 / mmhm: amortized latency = total time incl. rebuilds / updates (us) [configs/E6.json, aggregate.py]
\newcommand{\NumWFiveSphereLEightMmhmAmort}{28{,}700}
% w5b_L8 / mmhm: mean footprint_sum per update (sum of |occ| over deletions) [configs/E6.json, aggregate.py]
\newcommand{\NumWFiveSphereLEightMmhmFootprint}{0.500}
% w5b_L8 / mmhm: INIT build time (ms), median of seeds [configs/E6.json, aggregate.py]
\newcommand{\NumWFiveSphereLEightMmhmInitMs}{3{,}820}
% w5b_L8 / mmhm: max latency (us) [configs/E6.json, aggregate.py]
\newcommand{\NumWFiveSphereLEightMmhmMax}{383{,}000}
% w5b_L8 / mmhm: median per-update latency (us), 3 seeds pooled [configs/E6.json, aggregate.py]
\newcommand{\NumWFiveSphereLEightMmhmMed}{29{,}300}
% w5b_L8 / mmhm: nnz(R)+nnz(V) after the last update (mean of seeds) [configs/E6.json, aggregate.py]
\newcommand{\NumWFiveSphereLEightMmhmNnzEnd}{13{,}762{,}554}
% w5b_L8 / mmhm: nnz at the end / nnz after INIT [configs/E6.json, aggregate.py]
\newcommand{\NumWFiveSphereLEightMmhmNnzGrowth}{1.11}
% w5b_L8 / mmhm: max nnz(R)+nnz(V) over the run (mean of seeds) [configs/E6.json, aggregate.py]
\newcommand{\NumWFiveSphereLEightMmhmNnzMax}{13{,}762{,}554}
% w5b_L8 / mmhm: nnz(R)+nnz(V) after INIT (mean of seeds) [configs/E6.json, aggregate.py]
\newcommand{\NumWFiveSphereLEightMmhmNnzStart}{12{,}451{,}836}
% w5b_L8 / mmhm-norb: mean column additions per update (non-rebuild updates) [configs/E6.json, aggregate.py]
\newcommand{\NumWFiveSphereLEightMmhmNorbAdds}{0.500}
% w5b_L8 / mmhm-norb: amortized latency = total time incl. rebuilds / updates (us) [configs/E6.json, aggregate.py]
\newcommand{\NumWFiveSphereLEightMmhmNorbAmort}{29{,}400}
% w5b_L8 / mmhm-norb: mean footprint_sum per update (sum of |occ| over deletions) [configs/E6.json, aggregate.py]
\newcommand{\NumWFiveSphereLEightMmhmNorbFootprint}{0.500}
% w5b_L8 / mmhm-norb: INIT build time (ms), median of seeds [configs/E6.json, aggregate.py]
\newcommand{\NumWFiveSphereLEightMmhmNorbInitMs}{4{,}070}
% w5b_L8 / mmhm-norb: max latency (us) [configs/E6.json, aggregate.py]
\newcommand{\NumWFiveSphereLEightMmhmNorbMax}{364{,}000}
% w5b_L8 / mmhm-norb: median per-update latency (us), 3 seeds pooled [configs/E6.json, aggregate.py]
\newcommand{\NumWFiveSphereLEightMmhmNorbMed}{29{,}800}
% w5b_L8 / mmhm-norb: nnz(R)+nnz(V) after the last update (mean of seeds) [configs/E6.json, aggregate.py]
\newcommand{\NumWFiveSphereLEightMmhmNorbNnzEnd}{13{,}762{,}554}
% w5b_L8 / mmhm-norb: nnz at the end / nnz after INIT [configs/E6.json, aggregate.py]
\newcommand{\NumWFiveSphereLEightMmhmNorbNnzGrowth}{1.11}
% w5b_L8 / mmhm-norb: max nnz(R)+nnz(V) over the run (mean of seeds) [configs/E6.json, aggregate.py]
\newcommand{\NumWFiveSphereLEightMmhmNorbNnzMax}{13{,}762{,}554}
% w5b_L8 / mmhm-norb: nnz(R)+nnz(V) after INIT (mean of seeds) [configs/E6.json, aggregate.py]
\newcommand{\NumWFiveSphereLEightMmhmNorbNnzStart}{12{,}451{,}836}
% w5b_L8 / mmhm-norb: p95 latency (us) [configs/E6.json, aggregate.py]
\newcommand{\NumWFiveSphereLEightMmhmNorbPNinetyFive}{49{,}800}
% w5b_L8 / mmhm-norb: p99 latency (us) [configs/E6.json, aggregate.py]
\newcommand{\NumWFiveSphereLEightMmhmNorbPNinetyNine}{276{,}000}
% w5b_L8 / mmhm-norb: 25th percentile latency (us) [configs/E6.json, aggregate.py]
\newcommand{\NumWFiveSphereLEightMmhmNorbQOne}{4{,}380}
% w5b_L8 / mmhm-norb: 75th percentile latency (us) [configs/E6.json, aggregate.py]
\newcommand{\NumWFiveSphereLEightMmhmNorbQThree}{33{,}400}
% w5b_L8 / mmhm-norb: amortized latency of full-id / this method [configs/E6.json, aggregate.py]
\newcommand{\NumWFiveSphereLEightMmhmNorbSpeedupFullId}{6.04}
% w5b_L8 / mmhm-norb: median latency of full-id / this method [configs/E6.json, aggregate.py]
\newcommand{\NumWFiveSphereLEightMmhmNorbSpeedupFullIdMed}{6.32}
% w5b_L8 / mmhm-norb: amortized latency of gudhi (on its T or T' prefix) / this method [configs/E6.json, aggregate.py]
\newcommand{\NumWFiveSphereLEightMmhmNorbSpeedupGudhi}{49.5}
% w5b_L8 / mmhm-norb: updates timed per run (T or prefix T') [configs/E6.json, aggregate.py]
\newcommand{\NumWFiveSphereLEightMmhmNorbT}{2{,}000}
% w5b_L8 / mmhm: p95 latency (us) [configs/E6.json, aggregate.py]
\newcommand{\NumWFiveSphereLEightMmhmPNinetyFive}{49{,}200}
% w5b_L8 / mmhm: p99 latency (us) [configs/E6.json, aggregate.py]
\newcommand{\NumWFiveSphereLEightMmhmPNinetyNine}{270{,}000}
% w5b_L8 / mmhm: 25th percentile latency (us) [configs/E6.json, aggregate.py]
\newcommand{\NumWFiveSphereLEightMmhmQOne}{4{,}170}
% w5b_L8 / mmhm: 75th percentile latency (us) [configs/E6.json, aggregate.py]
\newcommand{\NumWFiveSphereLEightMmhmQThree}{33{,}000}
% w5b_L8 / mmhm: rebuilds per run (mean of seeds), default policy alpha=1, period off [configs/E6.json, aggregate.py]
\newcommand{\NumWFiveSphereLEightMmhmRebuilds}{0}
% w5b_L8 / mmhm: amortized latency of full-id / this method [configs/E6.json, aggregate.py]
\newcommand{\NumWFiveSphereLEightMmhmSpeedupFullId}{6.20}
% w5b_L8 / mmhm: median latency of full-id / this method [configs/E6.json, aggregate.py]
\newcommand{\NumWFiveSphereLEightMmhmSpeedupFullIdMed}{6.43}
% w5b_L8 / mmhm: amortized latency of gudhi (on its T or T' prefix) / this method [configs/E6.json, aggregate.py]
\newcommand{\NumWFiveSphereLEightMmhmSpeedupGudhi}{50.8}
% w5b_L8 / mmhm: updates timed per run (T or prefix T') [configs/E6.json, aggregate.py]
\newcommand{\NumWFiveSphereLEightMmhmT}{2{,}000}
% w5b_L8: mean atomic edits per update [configs/E6.json, aggregate.py]
\newcommand{\NumWFiveSphereLEightOps}{1.00}
% w5b_L8: |K_0| (mean of seeds) [configs/E6.json, aggregate.py]
\newcommand{\NumWFiveSphereLEightSize}{3{,}932{,}162}
% w5b_L5 / gudhi: amortized latency = total time incl. rebuilds / updates (us) [configs/E6.json, aggregate.py]
\newcommand{\NumWFiveSphereLFiveGudhiAmort}{12{,}600}
% w5b_L5 / gudhi: INIT build time (ms), median of seeds [configs/E6.json, aggregate.py]
\newcommand{\NumWFiveSphereLFiveGudhiInitMs}{31.6}
% w5b_L5 / gudhi: max latency (us) [configs/E6.json, aggregate.py]
\newcommand{\NumWFiveSphereLFiveGudhiMax}{29{,}400}
% w5b_L5 / gudhi: median per-update latency (us), 3 seeds pooled [configs/E6.json, aggregate.py]
\newcommand{\NumWFiveSphereLFiveGudhiMed}{12{,}700}
% w5b_L5 / gudhi: p95 latency (us) [configs/E6.json, aggregate.py]
\newcommand{\NumWFiveSphereLFiveGudhiPNinetyFive}{15{,}800}
% w5b_L5 / gudhi: p99 latency (us) [configs/E6.json, aggregate.py]
\newcommand{\NumWFiveSphereLFiveGudhiPNinetyNine}{18{,}800}
% w5b_L5 / gudhi: 25th percentile latency (us) [configs/E6.json, aggregate.py]
\newcommand{\NumWFiveSphereLFiveGudhiQOne}{11{,}500}
% w5b_L5 / gudhi: 75th percentile latency (us) [configs/E6.json, aggregate.py]
\newcommand{\NumWFiveSphereLFiveGudhiQThree}{13{,}200}
% w5b_L5 / gudhi: updates timed per run (T or prefix T') [configs/E6.json, aggregate.py]
\newcommand{\NumWFiveSphereLFiveGudhiT}{2{,}000}
% w5b_L7 / full-id: amortized latency = total time incl. rebuilds / updates (us) [configs/E6.json, aggregate.py]
\newcommand{\NumWFiveSphereLSevenFullIdAmort}{43{,}800}
% w5b_L7 / full-id: INIT build time (ms), median of seeds [configs/E6.json, aggregate.py]
\newcommand{\NumWFiveSphereLSevenFullIdInitMs}{528}
% w5b_L7 / full-id: max latency (us) [configs/E6.json, aggregate.py]
\newcommand{\NumWFiveSphereLSevenFullIdMax}{83{,}400}
% w5b_L7 / full-id: median per-update latency (us), 3 seeds pooled [configs/E6.json, aggregate.py]
\newcommand{\NumWFiveSphereLSevenFullIdMed}{46{,}100}
% w5b_L7 / full-id: p95 latency (us) [configs/E6.json, aggregate.py]
\newcommand{\NumWFiveSphereLSevenFullIdPNinetyFive}{54{,}400}
% w5b_L7 / full-id: p99 latency (us) [configs/E6.json, aggregate.py]
\newcommand{\NumWFiveSphereLSevenFullIdPNinetyNine}{64{,}900}
% w5b_L7 / full-id: 25th percentile latency (us) [configs/E6.json, aggregate.py]
\newcommand{\NumWFiveSphereLSevenFullIdQOne}{40{,}000}
% w5b_L7 / full-id: 75th percentile latency (us) [configs/E6.json, aggregate.py]
\newcommand{\NumWFiveSphereLSevenFullIdQThree}{48{,}400}
% w5b_L7 / full-id: amortized latency of gudhi (on its T or T' prefix) / this method [configs/E6.json, aggregate.py]
\newcommand{\NumWFiveSphereLSevenFullIdSpeedupGudhi}{6.16}
% w5b_L7 / full-id: updates timed per run (T or prefix T') [configs/E6.json, aggregate.py]
\newcommand{\NumWFiveSphereLSevenFullIdT}{2{,}000}
% w5b_L7 / gudhi: amortized latency = total time incl. rebuilds / updates (us) [configs/E6.json, aggregate.py]
\newcommand{\NumWFiveSphereLSevenGudhiAmort}{270{,}000}
% w5b_L7 / gudhi: INIT build time (ms), median of seeds [configs/E6.json, aggregate.py]
\newcommand{\NumWFiveSphereLSevenGudhiInitMs}{618}
% w5b_L7 / gudhi: max latency (us) [configs/E6.json, aggregate.py]
\newcommand{\NumWFiveSphereLSevenGudhiMax}{365{,}000}
% w5b_L7 / gudhi: median per-update latency (us), 3 seeds pooled [configs/E6.json, aggregate.py]
\newcommand{\NumWFiveSphereLSevenGudhiMed}{269{,}000}
% w5b_L7 / gudhi: p95 latency (us) [configs/E6.json, aggregate.py]
\newcommand{\NumWFiveSphereLSevenGudhiPNinetyFive}{310{,}000}
% w5b_L7 / gudhi: p99 latency (us) [configs/E6.json, aggregate.py]
\newcommand{\NumWFiveSphereLSevenGudhiPNinetyNine}{338{,}000}
% w5b_L7 / gudhi: 25th percentile latency (us) [configs/E6.json, aggregate.py]
\newcommand{\NumWFiveSphereLSevenGudhiQOne}{253{,}000}
% w5b_L7 / gudhi: 75th percentile latency (us) [configs/E6.json, aggregate.py]
\newcommand{\NumWFiveSphereLSevenGudhiQThree}{282{,}000}
% w5b_L7 / gudhi: amortized latency of full-id / this method [configs/E6.json, aggregate.py]
\newcommand{\NumWFiveSphereLSevenGudhiSpeedupFullId}{0.162}
% w5b_L7 / gudhi: median latency of full-id / this method [configs/E6.json, aggregate.py]
\newcommand{\NumWFiveSphereLSevenGudhiSpeedupFullIdMed}{0.171}
% w5b_L7 / gudhi: updates timed per run (T or prefix T') [configs/E6.json, aggregate.py]
\newcommand{\NumWFiveSphereLSevenGudhiT}{200}
% w5b_L7: mean atomic edits per update [configs/E6.json, aggregate.py]
\newcommand{\NumWFiveSphereLSevenOps}{1.00}
% w5b_L7: |K_0| (mean of seeds) [configs/E6.json, aggregate.py]
\newcommand{\NumWFiveSphereLSevenSize}{983{,}042}
% w5b_L6 / gudhi: amortized latency = total time incl. rebuilds / updates (us) [configs/E6.json, aggregate.py]
\newcommand{\NumWFiveSphereLSixGudhiAmort}{61{,}500}
% w5b_L6 / gudhi: INIT build time (ms), median of seeds [configs/E6.json, aggregate.py]
\newcommand{\NumWFiveSphereLSixGudhiInitMs}{151}
% w5b_L6 / gudhi: max latency (us) [configs/E6.json, aggregate.py]
\newcommand{\NumWFiveSphereLSixGudhiMax}{124{,}000}
% w5b_L6 / gudhi: median per-update latency (us), 3 seeds pooled [configs/E6.json, aggregate.py]
\newcommand{\NumWFiveSphereLSixGudhiMed}{60{,}900}
% w5b_L6 / gudhi: p95 latency (us) [configs/E6.json, aggregate.py]
\newcommand{\NumWFiveSphereLSixGudhiPNinetyFive}{76{,}000}
% w5b_L6 / gudhi: p99 latency (us) [configs/E6.json, aggregate.py]
\newcommand{\NumWFiveSphereLSixGudhiPNinetyNine}{87{,}800}
% w5b_L6 / gudhi: 25th percentile latency (us) [configs/E6.json, aggregate.py]
\newcommand{\NumWFiveSphereLSixGudhiQOne}{57{,}300}
% w5b_L6 / gudhi: 75th percentile latency (us) [configs/E6.json, aggregate.py]
\newcommand{\NumWFiveSphereLSixGudhiQThree}{64{,}000}
% w5b_L6 / gudhi: updates timed per run (T or prefix T') [configs/E6.json, aggregate.py]
\newcommand{\NumWFiveSphereLSixGudhiT}{2{,}000}
% max over W4 instances of beta_D / n at INIT [aggregate.py]
\newcommand{\NumWFourBetaDRatioMax}{6.3}
% min over W4 instances of beta_D / n at INIT [aggregate.py]
\newcommand{\NumWFourBetaDRatioMin}{6.3}
% w4_2d_n4000_D3: beta_D / n at INIT (mean of seeds) [aggregate.py (full-id betti)]
\newcommand{\NumWFourPlaneNFourThousandDThreeBetaDRatio}{6.33}
% w4_2d_n4000_D3: max beta_1 over all states of the 3 seeds [aggregate.py (full-id betti)]
\newcommand{\NumWFourPlaneNFourThousandDThreeBetaOneMax}{331}
% w4_2d_n4000_D3: min beta_1 over all states of the 3 seeds [aggregate.py (full-id betti)]
\newcommand{\NumWFourPlaneNFourThousandDThreeBetaOneMin}{308}
% w4_2d_n4000_D3: max beta_3 over all states of the 3 seeds [aggregate.py (full-id betti)]
\newcommand{\NumWFourPlaneNFourThousandDThreeBetaThreeMax}{25{,}967}
% w4_2d_n4000_D3: min beta_3 over all states of the 3 seeds [aggregate.py (full-id betti)]
\newcommand{\NumWFourPlaneNFourThousandDThreeBetaThreeMin}{24{,}744}
% w4_2d_n4000_D3: max beta_2 over all states of the 3 seeds [aggregate.py (full-id betti)]
\newcommand{\NumWFourPlaneNFourThousandDThreeBetaTwoMax}{1}
% w4_2d_n4000_D3: min beta_2 over all states of the 3 seeds [aggregate.py (full-id betti)]
\newcommand{\NumWFourPlaneNFourThousandDThreeBetaTwoMin}{0}
% w4_2d_n4000_D3: max beta_0 over all states of the 3 seeds [aggregate.py (full-id betti)]
\newcommand{\NumWFourPlaneNFourThousandDThreeBetaZeroMax}{3}
% w4_2d_n4000_D3: min beta_0 over all states of the 3 seeds [aggregate.py (full-id betti)]
\newcommand{\NumWFourPlaneNFourThousandDThreeBetaZeroMin}{2}
% w4_2d_n4000_D3 / full-id: amortized latency = total time incl. rebuilds / updates (us) [configs/E4.json, aggregate.py]
\newcommand{\NumWFourPlaneNFourThousandDThreeFullIdAmort}{4{,}830}
% w4_2d_n4000_D3 / full-id: INIT build time (ms), median of seeds [configs/E4.json, aggregate.py]
\newcommand{\NumWFourPlaneNFourThousandDThreeFullIdInitMs}{57.5}
% w4_2d_n4000_D3 / full-id: max latency (us) [configs/E4.json, aggregate.py]
\newcommand{\NumWFourPlaneNFourThousandDThreeFullIdMax}{12{,}000}
% w4_2d_n4000_D3 / full-id: median per-update latency (us), 1 seeds pooled [configs/E4.json, aggregate.py]
\newcommand{\NumWFourPlaneNFourThousandDThreeFullIdMed}{4{,}770}
% w4_2d_n4000_D3 / full-id: p95 latency (us) [configs/E4.json, aggregate.py]
\newcommand{\NumWFourPlaneNFourThousandDThreeFullIdPNinetyFive}{6{,}850}
% w4_2d_n4000_D3 / full-id: p99 latency (us) [configs/E4.json, aggregate.py]
\newcommand{\NumWFourPlaneNFourThousandDThreeFullIdPNinetyNine}{8{,}620}
% w4_2d_n4000_D3 / full-id: 25th percentile latency (us) [configs/E4.json, aggregate.py]
\newcommand{\NumWFourPlaneNFourThousandDThreeFullIdQOne}{4{,}090}
% w4_2d_n4000_D3 / full-id: 75th percentile latency (us) [configs/E4.json, aggregate.py]
\newcommand{\NumWFourPlaneNFourThousandDThreeFullIdQThree}{5{,}190}
% w4_2d_n4000_D3 / full-id: amortized latency of gudhi (on its T or T' prefix) / this method [configs/E4.json, aggregate.py]
\newcommand{\NumWFourPlaneNFourThousandDThreeFullIdSpeedupGudhi}{15.2}
% w4_2d_n4000_D3 / full-id: updates timed per run (T or prefix T') [configs/E4.json, aggregate.py]
\newcommand{\NumWFourPlaneNFourThousandDThreeFullIdT}{2{,}000}
% w4_2d_n4000_D3 / full-morse: amortized latency = total time incl. rebuilds / updates (us) [configs/E4.json, aggregate.py]
\newcommand{\NumWFourPlaneNFourThousandDThreeFullMorseAmort}{8{,}980}
% w4_2d_n4000_D3 / full-morse: INIT build time (ms), median of seeds [configs/E4.json, aggregate.py]
\newcommand{\NumWFourPlaneNFourThousandDThreeFullMorseInitMs}{62.8}
% w4_2d_n4000_D3 / full-morse: max latency (us) [configs/E4.json, aggregate.py]
\newcommand{\NumWFourPlaneNFourThousandDThreeFullMorseMax}{27{,}800}
% w4_2d_n4000_D3 / full-morse: median per-update latency (us), 3 seeds pooled [configs/E4.json, aggregate.py]
\newcommand{\NumWFourPlaneNFourThousandDThreeFullMorseMed}{8{,}810}
% w4_2d_n4000_D3: amortized latency full-morse / full-id [configs/E4.json, aggregate.py]
\newcommand{\NumWFourPlaneNFourThousandDThreeFullMorseOverFullId}{1.86}
% w4_2d_n4000_D3 / full-morse: p95 latency (us) [configs/E4.json, aggregate.py]
\newcommand{\NumWFourPlaneNFourThousandDThreeFullMorsePNinetyFive}{11{,}600}
% w4_2d_n4000_D3 / full-morse: p99 latency (us) [configs/E4.json, aggregate.py]
\newcommand{\NumWFourPlaneNFourThousandDThreeFullMorsePNinetyNine}{14{,}400}
% w4_2d_n4000_D3 / full-morse: 25th percentile latency (us) [configs/E4.json, aggregate.py]
\newcommand{\NumWFourPlaneNFourThousandDThreeFullMorseQOne}{7{,}980}
% w4_2d_n4000_D3 / full-morse: 75th percentile latency (us) [configs/E4.json, aggregate.py]
\newcommand{\NumWFourPlaneNFourThousandDThreeFullMorseQThree}{9{,}660}
% w4_2d_n4000_D3 / full-morse: amortized latency of full-id / this method [configs/E4.json, aggregate.py]
\newcommand{\NumWFourPlaneNFourThousandDThreeFullMorseSpeedupFullId}{0.538}
% w4_2d_n4000_D3 / full-morse: median latency of full-id / this method [configs/E4.json, aggregate.py]
\newcommand{\NumWFourPlaneNFourThousandDThreeFullMorseSpeedupFullIdMed}{0.542}
% w4_2d_n4000_D3 / full-morse: amortized latency of gudhi (on its T or T' prefix) / this method [configs/E4.json, aggregate.py]
\newcommand{\NumWFourPlaneNFourThousandDThreeFullMorseSpeedupGudhi}{8.15}
% w4_2d_n4000_D3 / full-morse: updates timed per run (T or prefix T') [configs/E4.json, aggregate.py]
\newcommand{\NumWFourPlaneNFourThousandDThreeFullMorseT}{2{,}000}
% w4_2d_n4000_D3 / gudhi: amortized latency = total time incl. rebuilds / updates (us) [configs/E4.json, aggregate.py]
\newcommand{\NumWFourPlaneNFourThousandDThreeGudhiAmort}{73{,}200}
% w4_2d_n4000_D3 / gudhi: INIT build time (ms), median of seeds [configs/E4.json, aggregate.py]
\newcommand{\NumWFourPlaneNFourThousandDThreeGudhiInitMs}{112}
% w4_2d_n4000_D3 / gudhi: max latency (us) [configs/E4.json, aggregate.py]
\newcommand{\NumWFourPlaneNFourThousandDThreeGudhiMax}{247{,}000}
% w4_2d_n4000_D3 / gudhi: median per-update latency (us), 3 seeds pooled [configs/E4.json, aggregate.py]
\newcommand{\NumWFourPlaneNFourThousandDThreeGudhiMed}{71{,}100}
% w4_2d_n4000_D3 / gudhi: p95 latency (us) [configs/E4.json, aggregate.py]
\newcommand{\NumWFourPlaneNFourThousandDThreeGudhiPNinetyFive}{94{,}500}
% w4_2d_n4000_D3 / gudhi: p99 latency (us) [configs/E4.json, aggregate.py]
\newcommand{\NumWFourPlaneNFourThousandDThreeGudhiPNinetyNine}{141{,}000}
% w4_2d_n4000_D3 / gudhi: 25th percentile latency (us) [configs/E4.json, aggregate.py]
\newcommand{\NumWFourPlaneNFourThousandDThreeGudhiQOne}{65{,}300}
% w4_2d_n4000_D3 / gudhi: 75th percentile latency (us) [configs/E4.json, aggregate.py]
\newcommand{\NumWFourPlaneNFourThousandDThreeGudhiQThree}{76{,}800}
% w4_2d_n4000_D3 / gudhi: amortized latency of full-id / this method [configs/E4.json, aggregate.py]
\newcommand{\NumWFourPlaneNFourThousandDThreeGudhiSpeedupFullId}{0.0660}
% w4_2d_n4000_D3 / gudhi: median latency of full-id / this method [configs/E4.json, aggregate.py]
\newcommand{\NumWFourPlaneNFourThousandDThreeGudhiSpeedupFullIdMed}{0.0671}
% w4_2d_n4000_D3 / gudhi: updates timed per run (T or prefix T') [configs/E4.json, aggregate.py]
\newcommand{\NumWFourPlaneNFourThousandDThreeGudhiT}{2{,}000}
% w4_2d_n4000_D3: mean atomic edits per update [configs/E4.json, aggregate.py]
\newcommand{\NumWFourPlaneNFourThousandDThreeOps}{189}
% w4_2d_n4000_D3: |K_0| (mean of seeds) [configs/E4.json, aggregate.py]
\newcommand{\NumWFourPlaneNFourThousandDThreeSize}{114{,}545}
% w4_2d_n4000_D2 / full-morse: amortized latency = total time incl. rebuilds / updates (us) [configs/E4.json, aggregate.py]
\newcommand{\NumWFourPlaneNFourThousandDTwoFullMorseAmort}{4{,}850}
% w4_2d_n4000_D2 / full-morse: INIT build time (ms), median of seeds [configs/E4.json, aggregate.py]
\newcommand{\NumWFourPlaneNFourThousandDTwoFullMorseInitMs}{30.5}
% w4_2d_n4000_D2 / full-morse: max latency (us) [configs/E4.json, aggregate.py]
\newcommand{\NumWFourPlaneNFourThousandDTwoFullMorseMax}{29{,}400}
% w4_2d_n4000_D2 / full-morse: median per-update latency (us), 3 seeds pooled [configs/E4.json, aggregate.py]
\newcommand{\NumWFourPlaneNFourThousandDTwoFullMorseMed}{4{,}720}
% w4_2d_n4000_D2 / full-morse: p95 latency (us) [configs/E4.json, aggregate.py]
\newcommand{\NumWFourPlaneNFourThousandDTwoFullMorsePNinetyFive}{6{,}350}
% w4_2d_n4000_D2 / full-morse: p99 latency (us) [configs/E4.json, aggregate.py]
\newcommand{\NumWFourPlaneNFourThousandDTwoFullMorsePNinetyNine}{7{,}510}
% w4_2d_n4000_D2 / full-morse: 25th percentile latency (us) [configs/E4.json, aggregate.py]
\newcommand{\NumWFourPlaneNFourThousandDTwoFullMorseQOne}{4{,}400}
% w4_2d_n4000_D2 / full-morse: 75th percentile latency (us) [configs/E4.json, aggregate.py]
\newcommand{\NumWFourPlaneNFourThousandDTwoFullMorseQThree}{5{,}010}
% w4_2d_n4000_D2 / full-morse: amortized latency of gudhi (on its T or T' prefix) / this method [configs/E4.json, aggregate.py]
\newcommand{\NumWFourPlaneNFourThousandDTwoFullMorseSpeedupGudhi}{7.15}
% w4_2d_n4000_D2 / full-morse: updates timed per run (T or prefix T') [configs/E4.json, aggregate.py]
\newcommand{\NumWFourPlaneNFourThousandDTwoFullMorseT}{2{,}000}
% w4_2d_n4000_D2 / gudhi: amortized latency = total time incl. rebuilds / updates (us) [configs/E4.json, aggregate.py]
\newcommand{\NumWFourPlaneNFourThousandDTwoGudhiAmort}{34{,}700}
% w4_2d_n4000_D2 / gudhi: INIT build time (ms), median of seeds [configs/E4.json, aggregate.py]
\newcommand{\NumWFourPlaneNFourThousandDTwoGudhiInitMs}{56.8}
% w4_2d_n4000_D2 / gudhi: max latency (us) [configs/E4.json, aggregate.py]
\newcommand{\NumWFourPlaneNFourThousandDTwoGudhiMax}{82{,}600}
% w4_2d_n4000_D2 / gudhi: median per-update latency (us), 3 seeds pooled [configs/E4.json, aggregate.py]
\newcommand{\NumWFourPlaneNFourThousandDTwoGudhiMed}{34{,}200}
% w4_2d_n4000_D2 / gudhi: p95 latency (us) [configs/E4.json, aggregate.py]
\newcommand{\NumWFourPlaneNFourThousandDTwoGudhiPNinetyFive}{43{,}400}
% w4_2d_n4000_D2 / gudhi: p99 latency (us) [configs/E4.json, aggregate.py]
\newcommand{\NumWFourPlaneNFourThousandDTwoGudhiPNinetyNine}{50{,}400}
% w4_2d_n4000_D2 / gudhi: 25th percentile latency (us) [configs/E4.json, aggregate.py]
\newcommand{\NumWFourPlaneNFourThousandDTwoGudhiQOne}{31{,}700}
% w4_2d_n4000_D2 / gudhi: 75th percentile latency (us) [configs/E4.json, aggregate.py]
\newcommand{\NumWFourPlaneNFourThousandDTwoGudhiQThree}{36{,}600}
% w4_2d_n4000_D2 / gudhi: updates timed per run (T or prefix T') [configs/E4.json, aggregate.py]
\newcommand{\NumWFourPlaneNFourThousandDTwoGudhiT}{2{,}000}
% w4_2d_n1000_D3 / full-morse: amortized latency = total time incl. rebuilds / updates (us) [configs/E4.json, aggregate.py]
\newcommand{\NumWFourPlaneNOneThousandDThreeFullMorseAmort}{1{,}810}
% w4_2d_n1000_D3 / full-morse: INIT build time (ms), median of seeds [configs/E4.json, aggregate.py]
\newcommand{\NumWFourPlaneNOneThousandDThreeFullMorseInitMs}{11.7}
% w4_2d_n1000_D3 / full-morse: max latency (us) [configs/E4.json, aggregate.py]
\newcommand{\NumWFourPlaneNOneThousandDThreeFullMorseMax}{16{,}300}
% w4_2d_n1000_D3 / full-morse: median per-update latency (us), 3 seeds pooled [configs/E4.json, aggregate.py]
\newcommand{\NumWFourPlaneNOneThousandDThreeFullMorseMed}{1{,}790}
% w4_2d_n1000_D3 / full-morse: p95 latency (us) [configs/E4.json, aggregate.py]
\newcommand{\NumWFourPlaneNOneThousandDThreeFullMorsePNinetyFive}{2{,}340}
% w4_2d_n1000_D3 / full-morse: p99 latency (us) [configs/E4.json, aggregate.py]
\newcommand{\NumWFourPlaneNOneThousandDThreeFullMorsePNinetyNine}{2{,}950}
% w4_2d_n1000_D3 / full-morse: 25th percentile latency (us) [configs/E4.json, aggregate.py]
\newcommand{\NumWFourPlaneNOneThousandDThreeFullMorseQOne}{1{,}590}
% w4_2d_n1000_D3 / full-morse: 75th percentile latency (us) [configs/E4.json, aggregate.py]
\newcommand{\NumWFourPlaneNOneThousandDThreeFullMorseQThree}{1{,}950}
% w4_2d_n1000_D3 / full-morse: amortized latency of gudhi (on its T or T' prefix) / this method [configs/E4.json, aggregate.py]
\newcommand{\NumWFourPlaneNOneThousandDThreeFullMorseSpeedupGudhi}{7.19}
% w4_2d_n1000_D3 / full-morse: updates timed per run (T or prefix T') [configs/E4.json, aggregate.py]
\newcommand{\NumWFourPlaneNOneThousandDThreeFullMorseT}{2{,}000}
% w4_2d_n1000_D3 / gudhi: amortized latency = total time incl. rebuilds / updates (us) [configs/E4.json, aggregate.py]
\newcommand{\NumWFourPlaneNOneThousandDThreeGudhiAmort}{13{,}000}
% w4_2d_n1000_D3 / gudhi: INIT build time (ms), median of seeds [configs/E4.json, aggregate.py]
\newcommand{\NumWFourPlaneNOneThousandDThreeGudhiInitMs}{22.5}
% w4_2d_n1000_D3 / gudhi: max latency (us) [configs/E4.json, aggregate.py]
\newcommand{\NumWFourPlaneNOneThousandDThreeGudhiMax}{38{,}700}
% w4_2d_n1000_D3 / gudhi: median per-update latency (us), 3 seeds pooled [configs/E4.json, aggregate.py]
\newcommand{\NumWFourPlaneNOneThousandDThreeGudhiMed}{12{,}800}
% w4_2d_n1000_D3 / gudhi: p95 latency (us) [configs/E4.json, aggregate.py]
\newcommand{\NumWFourPlaneNOneThousandDThreeGudhiPNinetyFive}{16{,}600}
% w4_2d_n1000_D3 / gudhi: p99 latency (us) [configs/E4.json, aggregate.py]
\newcommand{\NumWFourPlaneNOneThousandDThreeGudhiPNinetyNine}{20{,}800}
% w4_2d_n1000_D3 / gudhi: 25th percentile latency (us) [configs/E4.json, aggregate.py]
\newcommand{\NumWFourPlaneNOneThousandDThreeGudhiQOne}{12{,}100}
% w4_2d_n1000_D3 / gudhi: 75th percentile latency (us) [configs/E4.json, aggregate.py]
\newcommand{\NumWFourPlaneNOneThousandDThreeGudhiQThree}{13{,}500}
% w4_2d_n1000_D3 / gudhi: updates timed per run (T or prefix T') [configs/E4.json, aggregate.py]
\newcommand{\NumWFourPlaneNOneThousandDThreeGudhiT}{2{,}000}
% w4_2d_n1000_D2 / gudhi: amortized latency = total time incl. rebuilds / updates (us) [configs/E4.json, aggregate.py]
\newcommand{\NumWFourPlaneNOneThousandDTwoGudhiAmort}{6{,}280}
% w4_2d_n1000_D2 / gudhi: INIT build time (ms), median of seeds [configs/E4.json, aggregate.py]
\newcommand{\NumWFourPlaneNOneThousandDTwoGudhiInitMs}{11.4}
% w4_2d_n1000_D2 / gudhi: max latency (us) [configs/E4.json, aggregate.py]
\newcommand{\NumWFourPlaneNOneThousandDTwoGudhiMax}{32{,}300}
% w4_2d_n1000_D2 / gudhi: median per-update latency (us), 3 seeds pooled [configs/E4.json, aggregate.py]
\newcommand{\NumWFourPlaneNOneThousandDTwoGudhiMed}{6{,}160}
% w4_2d_n1000_D2 / gudhi: p95 latency (us) [configs/E4.json, aggregate.py]
\newcommand{\NumWFourPlaneNOneThousandDTwoGudhiPNinetyFive}{8{,}140}
% w4_2d_n1000_D2 / gudhi: p99 latency (us) [configs/E4.json, aggregate.py]
\newcommand{\NumWFourPlaneNOneThousandDTwoGudhiPNinetyNine}{10{,}000}
% w4_2d_n1000_D2 / gudhi: 25th percentile latency (us) [configs/E4.json, aggregate.py]
\newcommand{\NumWFourPlaneNOneThousandDTwoGudhiQOne}{5{,}900}
% w4_2d_n1000_D2 / gudhi: 75th percentile latency (us) [configs/E4.json, aggregate.py]
\newcommand{\NumWFourPlaneNOneThousandDTwoGudhiQThree}{6{,}520}
% w4_2d_n1000_D2 / gudhi: updates timed per run (T or prefix T') [configs/E4.json, aggregate.py]
\newcommand{\NumWFourPlaneNOneThousandDTwoGudhiT}{2{,}000}
% w4_2d_n16000_D2 / full-morse: amortized latency = total time incl. rebuilds / updates (us) [configs/E4.json, aggregate.py]
\newcommand{\NumWFourPlaneNSixteenThousandDTwoFullMorseAmort}{25{,}100}
% w4_2d_n16000_D2 / full-morse: INIT build time (ms), median of seeds [configs/E4.json, aggregate.py]
\newcommand{\NumWFourPlaneNSixteenThousandDTwoFullMorseInitMs}{148}
% w4_2d_n16000_D2 / full-morse: max latency (us) [configs/E4.json, aggregate.py]
\newcommand{\NumWFourPlaneNSixteenThousandDTwoFullMorseMax}{57{,}200}
% w4_2d_n16000_D2 / full-morse: median per-update latency (us), 3 seeds pooled [configs/E4.json, aggregate.py]
\newcommand{\NumWFourPlaneNSixteenThousandDTwoFullMorseMed}{24{,}300}
% w4_2d_n16000_D2 / full-morse: p95 latency (us) [configs/E4.json, aggregate.py]
\newcommand{\NumWFourPlaneNSixteenThousandDTwoFullMorsePNinetyFive}{33{,}400}
% w4_2d_n16000_D2 / full-morse: p99 latency (us) [configs/E4.json, aggregate.py]
\newcommand{\NumWFourPlaneNSixteenThousandDTwoFullMorsePNinetyNine}{42{,}200}
% w4_2d_n16000_D2 / full-morse: 25th percentile latency (us) [configs/E4.json, aggregate.py]
\newcommand{\NumWFourPlaneNSixteenThousandDTwoFullMorseQOne}{22{,}600}
% w4_2d_n16000_D2 / full-morse: 75th percentile latency (us) [configs/E4.json, aggregate.py]
\newcommand{\NumWFourPlaneNSixteenThousandDTwoFullMorseQThree}{26{,}500}
% w4_2d_n16000_D2 / full-morse: amortized latency of gudhi (on its T or T' prefix) / this method [configs/E4.json, aggregate.py]
\newcommand{\NumWFourPlaneNSixteenThousandDTwoFullMorseSpeedupGudhi}{7.81}
% w4_2d_n16000_D2 / full-morse: updates timed per run (T or prefix T') [configs/E4.json, aggregate.py]
\newcommand{\NumWFourPlaneNSixteenThousandDTwoFullMorseT}{2{,}000}
% w4_2d_n16000_D2 / gudhi: amortized latency = total time incl. rebuilds / updates (us) [configs/E4.json, aggregate.py]
\newcommand{\NumWFourPlaneNSixteenThousandDTwoGudhiAmort}{196{,}000}
% w4_2d_n16000_D2 / gudhi: INIT build time (ms), median of seeds [configs/E4.json, aggregate.py]
\newcommand{\NumWFourPlaneNSixteenThousandDTwoGudhiInitMs}{273}
% w4_2d_n16000_D2 / gudhi: max latency (us) [configs/E4.json, aggregate.py]
\newcommand{\NumWFourPlaneNSixteenThousandDTwoGudhiMax}{402{,}000}
% w4_2d_n16000_D2 / gudhi: median per-update latency (us), 3 seeds pooled [configs/E4.json, aggregate.py]
\newcommand{\NumWFourPlaneNSixteenThousandDTwoGudhiMed}{194{,}000}
% w4_2d_n16000_D2 / gudhi: p95 latency (us) [configs/E4.json, aggregate.py]
\newcommand{\NumWFourPlaneNSixteenThousandDTwoGudhiPNinetyFive}{232{,}000}
% w4_2d_n16000_D2 / gudhi: p99 latency (us) [configs/E4.json, aggregate.py]
\newcommand{\NumWFourPlaneNSixteenThousandDTwoGudhiPNinetyNine}{262{,}000}
% w4_2d_n16000_D2 / gudhi: 25th percentile latency (us) [configs/E4.json, aggregate.py]
\newcommand{\NumWFourPlaneNSixteenThousandDTwoGudhiQOne}{184{,}000}
% w4_2d_n16000_D2 / gudhi: 75th percentile latency (us) [configs/E4.json, aggregate.py]
\newcommand{\NumWFourPlaneNSixteenThousandDTwoGudhiQThree}{205{,}000}
% w4_2d_n16000_D2 / gudhi: updates timed per run (T or prefix T') [configs/E4.json, aggregate.py]
\newcommand{\NumWFourPlaneNSixteenThousandDTwoGudhiT}{2{,}000}
% w4_3d_n4000_D3 / full-morse: amortized latency = total time incl. rebuilds / updates (us) [configs/E4.json, aggregate.py]
\newcommand{\NumWFourSpaceNFourThousandDThreeFullMorseAmort}{6{,}290}
% w4_3d_n4000_D3 / full-morse: INIT build time (ms), median of seeds [configs/E4.json, aggregate.py]
\newcommand{\NumWFourSpaceNFourThousandDThreeFullMorseInitMs}{47.1}
% w4_3d_n4000_D3 / full-morse: max latency (us) [configs/E4.json, aggregate.py]
\newcommand{\NumWFourSpaceNFourThousandDThreeFullMorseMax}{13{,}900}
% w4_3d_n4000_D3 / full-morse: median per-update latency (us), 3 seeds pooled [configs/E4.json, aggregate.py]
\newcommand{\NumWFourSpaceNFourThousandDThreeFullMorseMed}{6{,}150}
% w4_3d_n4000_D3 / full-morse: p95 latency (us) [configs/E4.json, aggregate.py]
\newcommand{\NumWFourSpaceNFourThousandDThreeFullMorsePNinetyFive}{8{,}150}
% w4_3d_n4000_D3 / full-morse: p99 latency (us) [configs/E4.json, aggregate.py]
\newcommand{\NumWFourSpaceNFourThousandDThreeFullMorsePNinetyNine}{10{,}400}
% w4_3d_n4000_D3 / full-morse: 25th percentile latency (us) [configs/E4.json, aggregate.py]
\newcommand{\NumWFourSpaceNFourThousandDThreeFullMorseQOne}{5{,}570}
% w4_3d_n4000_D3 / full-morse: 75th percentile latency (us) [configs/E4.json, aggregate.py]
\newcommand{\NumWFourSpaceNFourThousandDThreeFullMorseQThree}{6{,}720}
% w4_3d_n4000_D3 / full-morse: amortized latency of gudhi (on its T or T' prefix) / this method [configs/E4.json, aggregate.py]
\newcommand{\NumWFourSpaceNFourThousandDThreeFullMorseSpeedupGudhi}{8.15}
% w4_3d_n4000_D3 / full-morse: updates timed per run (T or prefix T') [configs/E4.json, aggregate.py]
\newcommand{\NumWFourSpaceNFourThousandDThreeFullMorseT}{2{,}000}
% w4_3d_n4000_D3 / gudhi: amortized latency = total time incl. rebuilds / updates (us) [configs/E4.json, aggregate.py]
\newcommand{\NumWFourSpaceNFourThousandDThreeGudhiAmort}{51{,}300}
% w4_3d_n4000_D3 / gudhi: INIT build time (ms), median of seeds [configs/E4.json, aggregate.py]
\newcommand{\NumWFourSpaceNFourThousandDThreeGudhiInitMs}{75.5}
% w4_3d_n4000_D3 / gudhi: max latency (us) [configs/E4.json, aggregate.py]
\newcommand{\NumWFourSpaceNFourThousandDThreeGudhiMax}{100{,}000}
% w4_3d_n4000_D3 / gudhi: median per-update latency (us), 3 seeds pooled [configs/E4.json, aggregate.py]
\newcommand{\NumWFourSpaceNFourThousandDThreeGudhiMed}{50{,}400}
% w4_3d_n4000_D3 / gudhi: p95 latency (us) [configs/E4.json, aggregate.py]
\newcommand{\NumWFourSpaceNFourThousandDThreeGudhiPNinetyFive}{62{,}500}
% w4_3d_n4000_D3 / gudhi: p99 latency (us) [configs/E4.json, aggregate.py]
\newcommand{\NumWFourSpaceNFourThousandDThreeGudhiPNinetyNine}{76{,}400}
% w4_3d_n4000_D3 / gudhi: 25th percentile latency (us) [configs/E4.json, aggregate.py]
\newcommand{\NumWFourSpaceNFourThousandDThreeGudhiQOne}{46{,}900}
% w4_3d_n4000_D3 / gudhi: 75th percentile latency (us) [configs/E4.json, aggregate.py]
\newcommand{\NumWFourSpaceNFourThousandDThreeGudhiQThree}{54{,}200}
% w4_3d_n4000_D3 / gudhi: updates timed per run (T or prefix T') [configs/E4.json, aggregate.py]
\newcommand{\NumWFourSpaceNFourThousandDThreeGudhiT}{2{,}000}
% w4_3d_n1000_D3 / full-morse: amortized latency = total time incl. rebuilds / updates (us) [configs/E4.json, aggregate.py]
\newcommand{\NumWFourSpaceNOneThousandDThreeFullMorseAmort}{1{,}590}
% w4_3d_n1000_D3 / full-morse: INIT build time (ms), median of seeds [configs/E4.json, aggregate.py]
\newcommand{\NumWFourSpaceNOneThousandDThreeFullMorseInitMs}{9.65}
% w4_3d_n1000_D3 / full-morse: max latency (us) [configs/E4.json, aggregate.py]
\newcommand{\NumWFourSpaceNOneThousandDThreeFullMorseMax}{6{,}080}
% w4_3d_n1000_D3 / full-morse: median per-update latency (us), 3 seeds pooled [configs/E4.json, aggregate.py]
\newcommand{\NumWFourSpaceNOneThousandDThreeFullMorseMed}{1{,}550}
% w4_3d_n1000_D3 / full-morse: p95 latency (us) [configs/E4.json, aggregate.py]
\newcommand{\NumWFourSpaceNOneThousandDThreeFullMorsePNinetyFive}{2{,}010}
% w4_3d_n1000_D3 / full-morse: p99 latency (us) [configs/E4.json, aggregate.py]
\newcommand{\NumWFourSpaceNOneThousandDThreeFullMorsePNinetyNine}{2{,}320}
% w4_3d_n1000_D3 / full-morse: 25th percentile latency (us) [configs/E4.json, aggregate.py]
\newcommand{\NumWFourSpaceNOneThousandDThreeFullMorseQOne}{1{,}420}
% w4_3d_n1000_D3 / full-morse: 75th percentile latency (us) [configs/E4.json, aggregate.py]
\newcommand{\NumWFourSpaceNOneThousandDThreeFullMorseQThree}{1{,}700}
% w4_3d_n1000_D3 / full-morse: amortized latency of gudhi (on its T or T' prefix) / this method [configs/E4.json, aggregate.py]
\newcommand{\NumWFourSpaceNOneThousandDThreeFullMorseSpeedupGudhi}{6.91}
% w4_3d_n1000_D3 / full-morse: updates timed per run (T or prefix T') [configs/E4.json, aggregate.py]
\newcommand{\NumWFourSpaceNOneThousandDThreeFullMorseT}{2{,}000}
% w4_3d_n1000_D3 / gudhi: amortized latency = total time incl. rebuilds / updates (us) [configs/E4.json, aggregate.py]
\newcommand{\NumWFourSpaceNOneThousandDThreeGudhiAmort}{11{,}000}
% w4_3d_n1000_D3 / gudhi: INIT build time (ms), median of seeds [configs/E4.json, aggregate.py]
\newcommand{\NumWFourSpaceNOneThousandDThreeGudhiInitMs}{18.5}
% w4_3d_n1000_D3 / gudhi: max latency (us) [configs/E4.json, aggregate.py]
\newcommand{\NumWFourSpaceNOneThousandDThreeGudhiMax}{29{,}300}
% w4_3d_n1000_D3 / gudhi: median per-update latency (us), 3 seeds pooled [configs/E4.json, aggregate.py]
\newcommand{\NumWFourSpaceNOneThousandDThreeGudhiMed}{10{,}500}
% w4_3d_n1000_D3 / gudhi: p95 latency (us) [configs/E4.json, aggregate.py]
\newcommand{\NumWFourSpaceNOneThousandDThreeGudhiPNinetyFive}{14{,}900}
% w4_3d_n1000_D3 / gudhi: p99 latency (us) [configs/E4.json, aggregate.py]
\newcommand{\NumWFourSpaceNOneThousandDThreeGudhiPNinetyNine}{18{,}800}
% w4_3d_n1000_D3 / gudhi: 25th percentile latency (us) [configs/E4.json, aggregate.py]
\newcommand{\NumWFourSpaceNOneThousandDThreeGudhiQOne}{9{,}840}
% w4_3d_n1000_D3 / gudhi: 75th percentile latency (us) [configs/E4.json, aggregate.py]
\newcommand{\NumWFourSpaceNOneThousandDThreeGudhiQThree}{11{,}500}
% w4_3d_n1000_D3 / gudhi: updates timed per run (T or prefix T') [configs/E4.json, aggregate.py]
\newcommand{\NumWFourSpaceNOneThousandDThreeGudhiT}{2{,}000}
% w1_sphere_L5 / gudhi: amortized latency = total time incl. rebuilds / updates (us) [configs/E5.json, aggregate.py]
\newcommand{\NumWOneSphereLFiveGudhiAmort}{15{,}800}
% w1_sphere_L5 / gudhi: INIT build time (ms), median of seeds [configs/E5.json, aggregate.py]
\newcommand{\NumWOneSphereLFiveGudhiInitMs}{29.7}
% w1_sphere_L5 / gudhi: max latency (us) [configs/E5.json, aggregate.py]
\newcommand{\NumWOneSphereLFiveGudhiMax}{66{,}700}
% w1_sphere_L5 / gudhi: median per-update latency (us), 3 seeds pooled [configs/E5.json, aggregate.py]
\newcommand{\NumWOneSphereLFiveGudhiMed}{15{,}600}
% w1_sphere_L5 / gudhi: p95 latency (us) [configs/E5.json, aggregate.py]
\newcommand{\NumWOneSphereLFiveGudhiPNinetyFive}{21{,}000}
% w1_sphere_L5 / gudhi: p99 latency (us) [configs/E5.json, aggregate.py]
\newcommand{\NumWOneSphereLFiveGudhiPNinetyNine}{26{,}600}
% w1_sphere_L5 / gudhi: 25th percentile latency (us) [configs/E5.json, aggregate.py]
\newcommand{\NumWOneSphereLFiveGudhiQOne}{13{,}900}
% w1_sphere_L5 / gudhi: 75th percentile latency (us) [configs/E5.json, aggregate.py]
\newcommand{\NumWOneSphereLFiveGudhiQThree}{17{,}100}
% w1_sphere_L5 / gudhi: updates timed per run (T or prefix T') [configs/E5.json, aggregate.py]
\newcommand{\NumWOneSphereLFiveGudhiT}{2{,}000}
% w1_sphere_L4 / gudhi: amortized latency = total time incl. rebuilds / updates (us) [configs/E5.json, aggregate.py]
\newcommand{\NumWOneSphereLFourGudhiAmort}{2{,}910}
% w1_sphere_L4 / gudhi: INIT build time (ms), median of seeds [configs/E5.json, aggregate.py]
\newcommand{\NumWOneSphereLFourGudhiInitMs}{7.46}
% w1_sphere_L4 / gudhi: max latency (us) [configs/E5.json, aggregate.py]
\newcommand{\NumWOneSphereLFourGudhiMax}{18{,}800}
% w1_sphere_L4 / gudhi: median per-update latency (us), 3 seeds pooled [configs/E5.json, aggregate.py]
\newcommand{\NumWOneSphereLFourGudhiMed}{2{,}890}
% w1_sphere_L4 / gudhi: p95 latency (us) [configs/E5.json, aggregate.py]
\newcommand{\NumWOneSphereLFourGudhiPNinetyFive}{3{,}400}
% w1_sphere_L4 / gudhi: p99 latency (us) [configs/E5.json, aggregate.py]
\newcommand{\NumWOneSphereLFourGudhiPNinetyNine}{4{,}310}
% w1_sphere_L4 / gudhi: 25th percentile latency (us) [configs/E5.json, aggregate.py]
\newcommand{\NumWOneSphereLFourGudhiQOne}{2{,}820}
% w1_sphere_L4 / gudhi: 75th percentile latency (us) [configs/E5.json, aggregate.py]
\newcommand{\NumWOneSphereLFourGudhiQThree}{2{,}980}
% w1_sphere_L4 / gudhi: updates timed per run (T or prefix T') [configs/E5.json, aggregate.py]
\newcommand{\NumWOneSphereLFourGudhiT}{2{,}000}
% w1_sphere_L7: updates after which all holes are closed (beta_2 = 1), mean of seeds [aggregate.py (full-id betti)]
\newcommand{\NumWOneSphereLSevenClosedStates}{38.0}
% w1_sphere_L7 / full-id: amortized latency = total time incl. rebuilds / updates (us) [configs/E5.json, aggregate.py]
\newcommand{\NumWOneSphereLSevenFullIdAmort}{30{,}700}
% w1_sphere_L7 / full-id: INIT build time (ms), median of seeds [configs/E5.json, aggregate.py]
\newcommand{\NumWOneSphereLSevenFullIdInitMs}{562}
% w1_sphere_L7 / full-id: max latency (us) [configs/E5.json, aggregate.py]
\newcommand{\NumWOneSphereLSevenFullIdMax}{84{,}100}
% w1_sphere_L7 / full-id: median per-update latency (us), 3 seeds pooled [configs/E5.json, aggregate.py]
\newcommand{\NumWOneSphereLSevenFullIdMed}{29{,}100}
% w1_sphere_L7 / full-id: p95 latency (us) [configs/E5.json, aggregate.py]
\newcommand{\NumWOneSphereLSevenFullIdPNinetyFive}{44{,}700}
% w1_sphere_L7 / full-id: p99 latency (us) [configs/E5.json, aggregate.py]
\newcommand{\NumWOneSphereLSevenFullIdPNinetyNine}{52{,}200}
% w1_sphere_L7 / full-id: 25th percentile latency (us) [configs/E5.json, aggregate.py]
\newcommand{\NumWOneSphereLSevenFullIdQOne}{26{,}800}
% w1_sphere_L7 / full-id: 75th percentile latency (us) [configs/E5.json, aggregate.py]
\newcommand{\NumWOneSphereLSevenFullIdQThree}{32{,}400}
% w1_sphere_L7 / full-id: amortized latency of gudhi (on its T or T' prefix) / this method [configs/E5.json, aggregate.py]
\newcommand{\NumWOneSphereLSevenFullIdSpeedupGudhi}{8.87}
% w1_sphere_L7 / full-id: updates timed per run (T or prefix T') [configs/E5.json, aggregate.py]
\newcommand{\NumWOneSphereLSevenFullIdT}{2{,}000}
% w1_sphere_L7 / full-morse: amortized latency = total time incl. rebuilds / updates (us) [configs/E5.json, aggregate.py]
\newcommand{\NumWOneSphereLSevenFullMorseAmort}{61{,}500}
% w1_sphere_L7 / full-morse: INIT build time (ms), median of seeds [configs/E5.json, aggregate.py]
\newcommand{\NumWOneSphereLSevenFullMorseInitMs}{693}
% w1_sphere_L7 / full-morse: max latency (us) [configs/E5.json, aggregate.py]
\newcommand{\NumWOneSphereLSevenFullMorseMax}{163{,}000}
% w1_sphere_L7 / full-morse: median per-update latency (us), 3 seeds pooled [configs/E5.json, aggregate.py]
\newcommand{\NumWOneSphereLSevenFullMorseMed}{58{,}900}
% w1_sphere_L7: amortized latency full-morse / full-id [configs/E5.json, aggregate.py]
\newcommand{\NumWOneSphereLSevenFullMorseOverFullId}{2.00}
% w1_sphere_L7 / full-morse: p95 latency (us) [configs/E5.json, aggregate.py]
\newcommand{\NumWOneSphereLSevenFullMorsePNinetyFive}{82{,}000}
% w1_sphere_L7 / full-morse: p99 latency (us) [configs/E5.json, aggregate.py]
\newcommand{\NumWOneSphereLSevenFullMorsePNinetyNine}{114{,}000}
% w1_sphere_L7 / full-morse: 25th percentile latency (us) [configs/E5.json, aggregate.py]
\newcommand{\NumWOneSphereLSevenFullMorseQOne}{55{,}700}
% w1_sphere_L7 / full-morse: 75th percentile latency (us) [configs/E5.json, aggregate.py]
\newcommand{\NumWOneSphereLSevenFullMorseQThree}{63{,}000}
% w1_sphere_L7 / full-morse: amortized latency of full-id / this method [configs/E5.json, aggregate.py]
\newcommand{\NumWOneSphereLSevenFullMorseSpeedupFullId}{0.499}
% w1_sphere_L7 / full-morse: median latency of full-id / this method [configs/E5.json, aggregate.py]
\newcommand{\NumWOneSphereLSevenFullMorseSpeedupFullIdMed}{0.495}
% w1_sphere_L7 / full-morse: amortized latency of gudhi (on its T or T' prefix) / this method [configs/E5.json, aggregate.py]
\newcommand{\NumWOneSphereLSevenFullMorseSpeedupGudhi}{4.43}
% w1_sphere_L7 / full-morse: updates timed per run (T or prefix T') [configs/E5.json, aggregate.py]
\newcommand{\NumWOneSphereLSevenFullMorseT}{2{,}000}
% w1_sphere_L7 / gudhi: amortized latency = total time incl. rebuilds / updates (us) [configs/E5.json, aggregate.py]
\newcommand{\NumWOneSphereLSevenGudhiAmort}{272{,}000}
% w1_sphere_L7 / gudhi: INIT build time (ms), median of seeds [configs/E5.json, aggregate.py]
\newcommand{\NumWOneSphereLSevenGudhiInitMs}{584}
% w1_sphere_L7 / gudhi: max latency (us) [configs/E5.json, aggregate.py]
\newcommand{\NumWOneSphereLSevenGudhiMax}{370{,}000}
% w1_sphere_L7 / gudhi: median per-update latency (us), 3 seeds pooled [configs/E5.json, aggregate.py]
\newcommand{\NumWOneSphereLSevenGudhiMed}{272{,}000}
% w1_sphere_L7 / gudhi: p95 latency (us) [configs/E5.json, aggregate.py]
\newcommand{\NumWOneSphereLSevenGudhiPNinetyFive}{317{,}000}
% w1_sphere_L7 / gudhi: p99 latency (us) [configs/E5.json, aggregate.py]
\newcommand{\NumWOneSphereLSevenGudhiPNinetyNine}{333{,}000}
% w1_sphere_L7 / gudhi: 25th percentile latency (us) [configs/E5.json, aggregate.py]
\newcommand{\NumWOneSphereLSevenGudhiQOne}{258{,}000}
% w1_sphere_L7 / gudhi: 75th percentile latency (us) [configs/E5.json, aggregate.py]
\newcommand{\NumWOneSphereLSevenGudhiQThree}{285{,}000}
% w1_sphere_L7 / gudhi: amortized latency of full-id / this method [configs/E5.json, aggregate.py]
\newcommand{\NumWOneSphereLSevenGudhiSpeedupFullId}{0.113}
% w1_sphere_L7 / gudhi: median latency of full-id / this method [configs/E5.json, aggregate.py]
\newcommand{\NumWOneSphereLSevenGudhiSpeedupFullIdMed}{0.107}
% w1_sphere_L7 / gudhi: updates timed per run (T or prefix T') [configs/E5.json, aggregate.py]
\newcommand{\NumWOneSphereLSevenGudhiT}{200}
% w1_sphere_L7: mean atomic edits per update [configs/E5.json, aggregate.py]
\newcommand{\NumWOneSphereLSevenOps}{61.0}
% w1_sphere_L7: |K_0| (mean of seeds) [configs/E5.json, aggregate.py]
\newcommand{\NumWOneSphereLSevenSize}{983{,}042}
% w1_sphere_L6 / full-morse: amortized latency = total time incl. rebuilds / updates (us) [configs/E5.json, aggregate.py]
\newcommand{\NumWOneSphereLSixFullMorseAmort}{9{,}170}
% w1_sphere_L6 / full-morse: INIT build time (ms), median of seeds [configs/E5.json, aggregate.py]
\newcommand{\NumWOneSphereLSixFullMorseInitMs}{141}
% w1_sphere_L6 / full-morse: max latency (us) [configs/E5.json, aggregate.py]
\newcommand{\NumWOneSphereLSixFullMorseMax}{97{,}400}
% w1_sphere_L6 / full-morse: median per-update latency (us), 3 seeds pooled [configs/E5.json, aggregate.py]
\newcommand{\NumWOneSphereLSixFullMorseMed}{8{,}570}
% w1_sphere_L6 / full-morse: p95 latency (us) [configs/E5.json, aggregate.py]
\newcommand{\NumWOneSphereLSixFullMorsePNinetyFive}{12{,}300}
% w1_sphere_L6 / full-morse: p99 latency (us) [configs/E5.json, aggregate.py]
\newcommand{\NumWOneSphereLSixFullMorsePNinetyNine}{22{,}000}
% w1_sphere_L6 / full-morse: 25th percentile latency (us) [configs/E5.json, aggregate.py]
\newcommand{\NumWOneSphereLSixFullMorseQOne}{8{,}000}
% w1_sphere_L6 / full-morse: 75th percentile latency (us) [configs/E5.json, aggregate.py]
\newcommand{\NumWOneSphereLSixFullMorseQThree}{9{,}390}
% w1_sphere_L6 / full-morse: amortized latency of gudhi (on its T or T' prefix) / this method [configs/E5.json, aggregate.py]
\newcommand{\NumWOneSphereLSixFullMorseSpeedupGudhi}{7.23}
% w1_sphere_L6 / full-morse: updates timed per run (T or prefix T') [configs/E5.json, aggregate.py]
\newcommand{\NumWOneSphereLSixFullMorseT}{2{,}000}
% w1_sphere_L6 / gudhi: amortized latency = total time incl. rebuilds / updates (us) [configs/E5.json, aggregate.py]
\newcommand{\NumWOneSphereLSixGudhiAmort}{66{,}300}
% w1_sphere_L6 / gudhi: INIT build time (ms), median of seeds [configs/E5.json, aggregate.py]
\newcommand{\NumWOneSphereLSixGudhiInitMs}{132}
% w1_sphere_L6 / gudhi: max latency (us) [configs/E5.json, aggregate.py]
\newcommand{\NumWOneSphereLSixGudhiMax}{122{,}000}
% w1_sphere_L6 / gudhi: median per-update latency (us), 3 seeds pooled [configs/E5.json, aggregate.py]
\newcommand{\NumWOneSphereLSixGudhiMed}{65{,}300}
% w1_sphere_L6 / gudhi: p95 latency (us) [configs/E5.json, aggregate.py]
\newcommand{\NumWOneSphereLSixGudhiPNinetyFive}{83{,}500}
% w1_sphere_L6 / gudhi: p99 latency (us) [configs/E5.json, aggregate.py]
\newcommand{\NumWOneSphereLSixGudhiPNinetyNine}{101{,}000}
% w1_sphere_L6 / gudhi: 25th percentile latency (us) [configs/E5.json, aggregate.py]
\newcommand{\NumWOneSphereLSixGudhiQOne}{60{,}600}
% w1_sphere_L6 / gudhi: 75th percentile latency (us) [configs/E5.json, aggregate.py]
\newcommand{\NumWOneSphereLSixGudhiQThree}{69{,}700}
% w1_sphere_L6 / gudhi: updates timed per run (T or prefix T') [configs/E5.json, aggregate.py]
\newcommand{\NumWOneSphereLSixGudhiT}{2{,}000}
% w1_torus_500: updates after which all holes are closed (beta_2 = 1), mean of seeds [aggregate.py (full-id betti)]
\newcommand{\NumWOneTorusFiveHundredClosedStates}{37.7}
% w1_torus_500 / drb-id: mean column additions per update (non-rebuild updates) [configs/E5.json, aggregate.py]
\newcommand{\NumWOneTorusFiveHundredDrbIdAdds}{3{,}480}
% w1_torus_500 / drb-id: amortized latency = total time incl. rebuilds / updates (us) [configs/E5.json, aggregate.py]
\newcommand{\NumWOneTorusFiveHundredDrbIdAmort}{45{,}100}
% w1_torus_500 / drb-id: mean footprint_sum per update (sum of |occ| over deletions) [configs/E5.json, aggregate.py]
\newcommand{\NumWOneTorusFiveHundredDrbIdFootprint}{174}
% w1_torus_500 / drb-id: INIT build time (ms), median of seeds [configs/E5.json, aggregate.py]
\newcommand{\NumWOneTorusFiveHundredDrbIdInitMs}{1{,}010}
% w1_torus_500 / drb-id: max latency (us) [configs/E5.json, aggregate.py]
\newcommand{\NumWOneTorusFiveHundredDrbIdMax}{516{,}000}
% w1_torus_500 / drb-id: median per-update latency (us), 3 seeds pooled [configs/E5.json, aggregate.py]
\newcommand{\NumWOneTorusFiveHundredDrbIdMed}{10{,}900}
% w1_torus_500 / drb-id: nnz(R)+nnz(V) after the last update (mean of seeds) [configs/E5.json, aggregate.py]
\newcommand{\NumWOneTorusFiveHundredDrbIdNnzEnd}{11{,}710{,}711}
% w1_torus_500 / drb-id: nnz at the end / nnz after INIT [configs/E5.json, aggregate.py]
\newcommand{\NumWOneTorusFiveHundredDrbIdNnzGrowth}{2.46}
% w1_torus_500 / drb-id: max nnz(R)+nnz(V) over the run (mean of seeds) [configs/E5.json, aggregate.py]
\newcommand{\NumWOneTorusFiveHundredDrbIdNnzMax}{13{,}399{,}819}
% w1_torus_500 / drb-id: nnz(R)+nnz(V) after INIT (mean of seeds) [configs/E5.json, aggregate.py]
\newcommand{\NumWOneTorusFiveHundredDrbIdNnzStart}{4{,}752{,}487}
% w1_torus_500 / drb-id: p95 latency (us) [configs/E5.json, aggregate.py]
\newcommand{\NumWOneTorusFiveHundredDrbIdPNinetyFive}{246{,}000}
% w1_torus_500 / drb-id: p99 latency (us) [configs/E5.json, aggregate.py]
\newcommand{\NumWOneTorusFiveHundredDrbIdPNinetyNine}{313{,}000}
% w1_torus_500 / drb-id: 25th percentile latency (us) [configs/E5.json, aggregate.py]
\newcommand{\NumWOneTorusFiveHundredDrbIdQOne}{647}
% w1_torus_500 / drb-id: 75th percentile latency (us) [configs/E5.json, aggregate.py]
\newcommand{\NumWOneTorusFiveHundredDrbIdQThree}{21{,}100}
% w1_torus_500 / drb-id-rb: mean column additions per update (non-rebuild updates) [configs/E5.json, aggregate.py]
\newcommand{\NumWOneTorusFiveHundredDrbIdRbAdds}{2{,}810}
% w1_torus_500 / drb-id-rb: amortized latency = total time incl. rebuilds / updates (us) [configs/E5.json, aggregate.py]
\newcommand{\NumWOneTorusFiveHundredDrbIdRbAmort}{57{,}600}
% w1_torus_500 / drb-id-rb: mean footprint_sum per update (sum of |occ| over deletions) [configs/E5.json, aggregate.py]
\newcommand{\NumWOneTorusFiveHundredDrbIdRbFootprint}{125}
% w1_torus_500 / drb-id-rb: INIT build time (ms), median of seeds [configs/E5.json, aggregate.py]
\newcommand{\NumWOneTorusFiveHundredDrbIdRbInitMs}{986}
% w1_torus_500 / drb-id-rb: max latency (us) [configs/E5.json, aggregate.py]
\newcommand{\NumWOneTorusFiveHundredDrbIdRbMax}{576{,}000}
% w1_torus_500 / drb-id-rb: median per-update latency (us), 3 seeds pooled [configs/E5.json, aggregate.py]
\newcommand{\NumWOneTorusFiveHundredDrbIdRbMed}{13{,}000}
% w1_torus_500 / drb-id-rb: nnz(R)+nnz(V) after the last update (mean of seeds) [configs/E5.json, aggregate.py]
\newcommand{\NumWOneTorusFiveHundredDrbIdRbNnzEnd}{7{,}299{,}510}
% w1_torus_500 / drb-id-rb: nnz at the end / nnz after INIT [configs/E5.json, aggregate.py]
\newcommand{\NumWOneTorusFiveHundredDrbIdRbNnzGrowth}{1.54}
% w1_torus_500 / drb-id-rb: max nnz(R)+nnz(V) over the run (mean of seeds) [configs/E5.json, aggregate.py]
\newcommand{\NumWOneTorusFiveHundredDrbIdRbNnzMax}{10{,}610{,}871}
% w1_torus_500 / drb-id-rb: nnz(R)+nnz(V) after INIT (mean of seeds) [configs/E5.json, aggregate.py]
\newcommand{\NumWOneTorusFiveHundredDrbIdRbNnzStart}{4{,}752{,}487}
% w1_torus_500 / drb-id-rb: p95 latency (us) [configs/E5.json, aggregate.py]
\newcommand{\NumWOneTorusFiveHundredDrbIdRbPNinetyFive}{335{,}000}
% w1_torus_500 / drb-id-rb: p99 latency (us) [configs/E5.json, aggregate.py]
\newcommand{\NumWOneTorusFiveHundredDrbIdRbPNinetyNine}{420{,}000}
% w1_torus_500 / drb-id-rb: 25th percentile latency (us) [configs/E5.json, aggregate.py]
\newcommand{\NumWOneTorusFiveHundredDrbIdRbQOne}{558}
% w1_torus_500 / drb-id-rb: 75th percentile latency (us) [configs/E5.json, aggregate.py]
\newcommand{\NumWOneTorusFiveHundredDrbIdRbQThree}{35{,}600}
% w1_torus_500 / drb-id-rb: rebuilds per run (mean of seeds), default policy alpha=1, period off [configs/E5.json, aggregate.py]
\newcommand{\NumWOneTorusFiveHundredDrbIdRbRebuilds}{1.7}
% w1_torus_500 / drb-id-rb: amortized latency of full-id / this method [configs/E5.json, aggregate.py]
\newcommand{\NumWOneTorusFiveHundredDrbIdRbSpeedupFullId}{0.663}
% w1_torus_500 / drb-id-rb: median latency of full-id / this method [configs/E5.json, aggregate.py]
\newcommand{\NumWOneTorusFiveHundredDrbIdRbSpeedupFullIdMed}{2.68}
% w1_torus_500 / drb-id-rb: amortized latency of gudhi (on its T or T' prefix) / this method [configs/E5.json, aggregate.py]
\newcommand{\NumWOneTorusFiveHundredDrbIdRbSpeedupGudhi}{7.40}
% w1_torus_500 / drb-id-rb: updates timed per run (T or prefix T') [configs/E5.json, aggregate.py]
\newcommand{\NumWOneTorusFiveHundredDrbIdRbT}{2{,}000}
% w1_torus_500 / drb-id: amortized latency of full-id / this method [configs/E5.json, aggregate.py]
\newcommand{\NumWOneTorusFiveHundredDrbIdSpeedupFullId}{0.848}
% w1_torus_500 / drb-id: median latency of full-id / this method [configs/E5.json, aggregate.py]
\newcommand{\NumWOneTorusFiveHundredDrbIdSpeedupFullIdMed}{3.19}
% w1_torus_500 / drb-id: amortized latency of gudhi (on its T or T' prefix) / this method [configs/E5.json, aggregate.py]
\newcommand{\NumWOneTorusFiveHundredDrbIdSpeedupGudhi}{9.46}
% w1_torus_500 / drb-id: updates timed per run (T or prefix T') [configs/E5.json, aggregate.py]
\newcommand{\NumWOneTorusFiveHundredDrbIdT}{2{,}000}
% w1_torus_500 / euler: amortized latency = total time incl. rebuilds / updates (us) [configs/E5.json, aggregate.py]
\newcommand{\NumWOneTorusFiveHundredEulerAmort}{30.1}
% w1_torus_500 / euler: INIT build time (ms), median of seeds [configs/E5.json, aggregate.py]
\newcommand{\NumWOneTorusFiveHundredEulerInitMs}{855}
% w1_torus_500 / euler: max latency (us) [configs/E5.json, aggregate.py]
\newcommand{\NumWOneTorusFiveHundredEulerMax}{1{,}010}
% w1_torus_500 / euler: median per-update latency (us), 3 seeds pooled [configs/E5.json, aggregate.py]
\newcommand{\NumWOneTorusFiveHundredEulerMed}{32.1}
% w1_torus_500 / euler: p95 latency (us) [configs/E5.json, aggregate.py]
\newcommand{\NumWOneTorusFiveHundredEulerPNinetyFive}{50.9}
% w1_torus_500 / euler: p99 latency (us) [configs/E5.json, aggregate.py]
\newcommand{\NumWOneTorusFiveHundredEulerPNinetyNine}{74.5}
% w1_torus_500 / euler: 25th percentile latency (us) [configs/E5.json, aggregate.py]
\newcommand{\NumWOneTorusFiveHundredEulerQOne}{16.3}
% w1_torus_500 / euler: 75th percentile latency (us) [configs/E5.json, aggregate.py]
\newcommand{\NumWOneTorusFiveHundredEulerQThree}{38.9}
% w1_torus_500 / euler: amortized latency of full-id / this method [configs/E5.json, aggregate.py]
\newcommand{\NumWOneTorusFiveHundredEulerSpeedupFullId}{1{,}270}
% w1_torus_500 / euler: median latency of full-id / this method [configs/E5.json, aggregate.py]
\newcommand{\NumWOneTorusFiveHundredEulerSpeedupFullIdMed}{1{,}080}
% w1_torus_500 / euler: amortized latency of gudhi (on its T or T' prefix) / this method [configs/E5.json, aggregate.py]
\newcommand{\NumWOneTorusFiveHundredEulerSpeedupGudhi}{14{,}200}
% w1_torus_500 / euler: updates timed per run (T or prefix T') [configs/E5.json, aggregate.py]
\newcommand{\NumWOneTorusFiveHundredEulerT}{2{,}000}
% w1_torus_500 / full-id: amortized latency = total time incl. rebuilds / updates (us) [configs/E5.json, aggregate.py]
\newcommand{\NumWOneTorusFiveHundredFullIdAmort}{38{,}200}
% w1_torus_500 / full-id: INIT build time (ms), median of seeds [configs/E5.json, aggregate.py]
\newcommand{\NumWOneTorusFiveHundredFullIdInitMs}{894}
% w1_torus_500 / full-id: max latency (us) [configs/E5.json, aggregate.py]
\newcommand{\NumWOneTorusFiveHundredFullIdMax}{108{,}000}
% w1_torus_500 / full-id: median per-update latency (us), 3 seeds pooled [configs/E5.json, aggregate.py]
\newcommand{\NumWOneTorusFiveHundredFullIdMed}{34{,}700}
% w1_torus_500 / full-id: p95 latency (us) [configs/E5.json, aggregate.py]
\newcommand{\NumWOneTorusFiveHundredFullIdPNinetyFive}{65{,}800}
% w1_torus_500 / full-id: p99 latency (us) [configs/E5.json, aggregate.py]
\newcommand{\NumWOneTorusFiveHundredFullIdPNinetyNine}{79{,}000}
% w1_torus_500 / full-id: 25th percentile latency (us) [configs/E5.json, aggregate.py]
\newcommand{\NumWOneTorusFiveHundredFullIdQOne}{28{,}400}
% w1_torus_500 / full-id: 75th percentile latency (us) [configs/E5.json, aggregate.py]
\newcommand{\NumWOneTorusFiveHundredFullIdQThree}{44{,}900}
% w1_torus_500 / full-id: amortized latency of gudhi (on its T or T' prefix) / this method [configs/E5.json, aggregate.py]
\newcommand{\NumWOneTorusFiveHundredFullIdSpeedupGudhi}{11.2}
% w1_torus_500 / full-id: updates timed per run (T or prefix T') [configs/E5.json, aggregate.py]
\newcommand{\NumWOneTorusFiveHundredFullIdT}{2{,}000}
% w1_torus_500 / full-morse: amortized latency = total time incl. rebuilds / updates (us) [configs/E5.json, aggregate.py]
\newcommand{\NumWOneTorusFiveHundredFullMorseAmort}{98{,}200}
% w1_torus_500 / full-morse: INIT build time (ms), median of seeds [configs/E5.json, aggregate.py]
\newcommand{\NumWOneTorusFiveHundredFullMorseInitMs}{1{,}150}
% w1_torus_500 / full-morse: max latency (us) [configs/E5.json, aggregate.py]
\newcommand{\NumWOneTorusFiveHundredFullMorseMax}{371{,}000}
% w1_torus_500 / full-morse: median per-update latency (us), 3 seeds pooled [configs/E5.json, aggregate.py]
\newcommand{\NumWOneTorusFiveHundredFullMorseMed}{94{,}600}
% w1_torus_500: amortized latency full-morse / full-id [configs/E5.json, aggregate.py]
\newcommand{\NumWOneTorusFiveHundredFullMorseOverFullId}{2.57}
% w1_torus_500 / full-morse: p95 latency (us) [configs/E5.json, aggregate.py]
\newcommand{\NumWOneTorusFiveHundredFullMorsePNinetyFive}{125{,}000}
% w1_torus_500 / full-morse: p99 latency (us) [configs/E5.json, aggregate.py]
\newcommand{\NumWOneTorusFiveHundredFullMorsePNinetyNine}{183{,}000}
% w1_torus_500 / full-morse: 25th percentile latency (us) [configs/E5.json, aggregate.py]
\newcommand{\NumWOneTorusFiveHundredFullMorseQOne}{90{,}100}
% w1_torus_500 / full-morse: 75th percentile latency (us) [configs/E5.json, aggregate.py]
\newcommand{\NumWOneTorusFiveHundredFullMorseQThree}{100{,}000}
% w1_torus_500 / full-morse: amortized latency of full-id / this method [configs/E5.json, aggregate.py]
\newcommand{\NumWOneTorusFiveHundredFullMorseSpeedupFullId}{0.389}
% w1_torus_500 / full-morse: median latency of full-id / this method [configs/E5.json, aggregate.py]
\newcommand{\NumWOneTorusFiveHundredFullMorseSpeedupFullIdMed}{0.367}
% w1_torus_500 / full-morse: amortized latency of gudhi (on its T or T' prefix) / this method [configs/E5.json, aggregate.py]
\newcommand{\NumWOneTorusFiveHundredFullMorseSpeedupGudhi}{4.35}
% w1_torus_500 / full-morse: updates timed per run (T or prefix T') [configs/E5.json, aggregate.py]
\newcommand{\NumWOneTorusFiveHundredFullMorseT}{2{,}000}
% w1_torus_500 / gudhi: amortized latency = total time incl. rebuilds / updates (us) [configs/E5.json, aggregate.py]
\newcommand{\NumWOneTorusFiveHundredGudhiAmort}{427{,}000}
% w1_torus_500 / gudhi: INIT build time (ms), median of seeds [configs/E5.json, aggregate.py]
\newcommand{\NumWOneTorusFiveHundredGudhiInitMs}{910}
% w1_torus_500 / gudhi: max latency (us) [configs/E5.json, aggregate.py]
\newcommand{\NumWOneTorusFiveHundredGudhiMax}{567{,}000}
% w1_torus_500 / gudhi: median per-update latency (us), 3 seeds pooled [configs/E5.json, aggregate.py]
\newcommand{\NumWOneTorusFiveHundredGudhiMed}{421{,}000}
% w1_torus_500 / gudhi: p95 latency (us) [configs/E5.json, aggregate.py]
\newcommand{\NumWOneTorusFiveHundredGudhiPNinetyFive}{503{,}000}
% w1_torus_500 / gudhi: p99 latency (us) [configs/E5.json, aggregate.py]
\newcommand{\NumWOneTorusFiveHundredGudhiPNinetyNine}{546{,}000}
% w1_torus_500 / gudhi: 25th percentile latency (us) [configs/E5.json, aggregate.py]
\newcommand{\NumWOneTorusFiveHundredGudhiQOne}{397{,}000}
% w1_torus_500 / gudhi: 75th percentile latency (us) [configs/E5.json, aggregate.py]
\newcommand{\NumWOneTorusFiveHundredGudhiQThree}{450{,}000}
% w1_torus_500 / gudhi: amortized latency of full-id / this method [configs/E5.json, aggregate.py]
\newcommand{\NumWOneTorusFiveHundredGudhiSpeedupFullId}{0.0896}
% w1_torus_500 / gudhi: median latency of full-id / this method [configs/E5.json, aggregate.py]
\newcommand{\NumWOneTorusFiveHundredGudhiSpeedupFullIdMed}{0.0825}
% w1_torus_500 / gudhi: updates timed per run (T or prefix T') [configs/E5.json, aggregate.py]
\newcommand{\NumWOneTorusFiveHundredGudhiT}{200}
% w1_torus_500 / mmhm: mean column additions per update (non-rebuild updates) [configs/E5.json, aggregate.py]
\newcommand{\NumWOneTorusFiveHundredMmhmAdds}{2{,}060}
% w1_torus_500 / mmhm: amortized latency = total time incl. rebuilds / updates (us) [configs/E5.json, aggregate.py]
\newcommand{\NumWOneTorusFiveHundredMmhmAmort}{34{,}000}
% w1_torus_500 / mmhm: mean footprint_sum per update (sum of |occ| over deletions) [configs/E5.json, aggregate.py]
\newcommand{\NumWOneTorusFiveHundredMmhmFootprint}{125}
% w1_torus_500 / mmhm: INIT build time (ms), median of seeds [configs/E5.json, aggregate.py]
\newcommand{\NumWOneTorusFiveHundredMmhmInitMs}{1{,}220}
% w1_torus_500 / mmhm: max latency (us) [configs/E5.json, aggregate.py]
\newcommand{\NumWOneTorusFiveHundredMmhmMax}{583{,}000}
% w1_torus_500 / mmhm: median per-update latency (us), 3 seeds pooled [configs/E5.json, aggregate.py]
\newcommand{\NumWOneTorusFiveHundredMmhmMed}{11{,}500}
% w1_torus_500 / mmhm: nnz(R)+nnz(V) after the last update (mean of seeds) [configs/E5.json, aggregate.py]
\newcommand{\NumWOneTorusFiveHundredMmhmNnzEnd}{6{,}495{,}424}
% w1_torus_500 / mmhm: nnz at the end / nnz after INIT [configs/E5.json, aggregate.py]
\newcommand{\NumWOneTorusFiveHundredMmhmNnzGrowth}{1.37}
% w1_torus_500 / mmhm: max nnz(R)+nnz(V) over the run (mean of seeds) [configs/E5.json, aggregate.py]
\newcommand{\NumWOneTorusFiveHundredMmhmNnzMax}{9{,}228{,}533}
% w1_torus_500 / mmhm: nnz(R)+nnz(V) after INIT (mean of seeds) [configs/E5.json, aggregate.py]
\newcommand{\NumWOneTorusFiveHundredMmhmNnzStart}{4{,}751{,}324}
% w1_torus_500 / mmhm-norb: mean column additions per update (non-rebuild updates) [configs/E5.json, aggregate.py]
\newcommand{\NumWOneTorusFiveHundredMmhmNorbAdds}{1{,}940}
% w1_torus_500 / mmhm-norb: amortized latency = total time incl. rebuilds / updates (us) [configs/E5.json, aggregate.py]
\newcommand{\NumWOneTorusFiveHundredMmhmNorbAmort}{33{,}400}
% w1_torus_500 / mmhm-norb: mean footprint_sum per update (sum of |occ| over deletions) [configs/E5.json, aggregate.py]
\newcommand{\NumWOneTorusFiveHundredMmhmNorbFootprint}{155}
% w1_torus_500 / mmhm-norb: INIT build time (ms), median of seeds [configs/E5.json, aggregate.py]
\newcommand{\NumWOneTorusFiveHundredMmhmNorbInitMs}{1{,}130}
% w1_torus_500 / mmhm-norb: max latency (us) [configs/E5.json, aggregate.py]
\newcommand{\NumWOneTorusFiveHundredMmhmNorbMax}{511{,}000}
% w1_torus_500 / mmhm-norb: median per-update latency (us), 3 seeds pooled [configs/E5.json, aggregate.py]
\newcommand{\NumWOneTorusFiveHundredMmhmNorbMed}{11{,}700}
% w1_torus_500 / mmhm-norb: nnz(R)+nnz(V) after the last update (mean of seeds) [configs/E5.json, aggregate.py]
\newcommand{\NumWOneTorusFiveHundredMmhmNorbNnzEnd}{8{,}729{,}344}
% w1_torus_500 / mmhm-norb: nnz at the end / nnz after INIT [configs/E5.json, aggregate.py]
\newcommand{\NumWOneTorusFiveHundredMmhmNorbNnzGrowth}{1.84}
% w1_torus_500 / mmhm-norb: max nnz(R)+nnz(V) over the run (mean of seeds) [configs/E5.json, aggregate.py]
\newcommand{\NumWOneTorusFiveHundredMmhmNorbNnzMax}{10{,}329{,}205}
% w1_torus_500 / mmhm-norb: nnz(R)+nnz(V) after INIT (mean of seeds) [configs/E5.json, aggregate.py]
\newcommand{\NumWOneTorusFiveHundredMmhmNorbNnzStart}{4{,}751{,}324}
% w1_torus_500 / mmhm-norb: p95 latency (us) [configs/E5.json, aggregate.py]
\newcommand{\NumWOneTorusFiveHundredMmhmNorbPNinetyFive}{184{,}000}
% w1_torus_500 / mmhm-norb: p99 latency (us) [configs/E5.json, aggregate.py]
\newcommand{\NumWOneTorusFiveHundredMmhmNorbPNinetyNine}{401{,}000}
% w1_torus_500 / mmhm-norb: 25th percentile latency (us) [configs/E5.json, aggregate.py]
\newcommand{\NumWOneTorusFiveHundredMmhmNorbQOne}{386}
% w1_torus_500 / mmhm-norb: 75th percentile latency (us) [configs/E5.json, aggregate.py]
\newcommand{\NumWOneTorusFiveHundredMmhmNorbQThree}{21{,}300}
% w1_torus_500 / mmhm-norb: amortized latency of full-id / this method [configs/E5.json, aggregate.py]
\newcommand{\NumWOneTorusFiveHundredMmhmNorbSpeedupFullId}{1.14}
% w1_torus_500 / mmhm-norb: median latency of full-id / this method [configs/E5.json, aggregate.py]
\newcommand{\NumWOneTorusFiveHundredMmhmNorbSpeedupFullIdMed}{2.96}
% w1_torus_500 / mmhm-norb: amortized latency of gudhi (on its T or T' prefix) / this method [configs/E5.json, aggregate.py]
\newcommand{\NumWOneTorusFiveHundredMmhmNorbSpeedupGudhi}{12.8}
% w1_torus_500 / mmhm-norb: updates timed per run (T or prefix T') [configs/E5.json, aggregate.py]
\newcommand{\NumWOneTorusFiveHundredMmhmNorbT}{2{,}000}
% w1_torus_500 / mmhm: p95 latency (us) [configs/E5.json, aggregate.py]
\newcommand{\NumWOneTorusFiveHundredMmhmPNinetyFive}{183{,}000}
% w1_torus_500 / mmhm: p99 latency (us) [configs/E5.json, aggregate.py]
\newcommand{\NumWOneTorusFiveHundredMmhmPNinetyNine}{397{,}000}
% w1_torus_500 / mmhm: 25th percentile latency (us) [configs/E5.json, aggregate.py]
\newcommand{\NumWOneTorusFiveHundredMmhmQOne}{381}
% w1_torus_500 / mmhm: 75th percentile latency (us) [configs/E5.json, aggregate.py]
\newcommand{\NumWOneTorusFiveHundredMmhmQThree}{22{,}300}
% w1_torus_500 / mmhm: rebuilds per run (mean of seeds), default policy alpha=1, period off [configs/E5.json, aggregate.py]
\newcommand{\NumWOneTorusFiveHundredMmhmRebuilds}{1.0}
% w1_torus_500 / mmhm: amortized latency of full-id / this method [configs/E5.json, aggregate.py]
\newcommand{\NumWOneTorusFiveHundredMmhmSpeedupFullId}{1.13}
% w1_torus_500 / mmhm: median latency of full-id / this method [configs/E5.json, aggregate.py]
\newcommand{\NumWOneTorusFiveHundredMmhmSpeedupFullIdMed}{3.02}
% w1_torus_500 / mmhm: amortized latency of gudhi (on its T or T' prefix) / this method [configs/E5.json, aggregate.py]
\newcommand{\NumWOneTorusFiveHundredMmhmSpeedupGudhi}{12.6}
% w1_torus_500 / mmhm: updates timed per run (T or prefix T') [configs/E5.json, aggregate.py]
\newcommand{\NumWOneTorusFiveHundredMmhmT}{2{,}000}
% w1_torus_500: mean atomic edits per update [configs/E5.json, aggregate.py]
\newcommand{\NumWOneTorusFiveHundredOps}{61.0}
% w1_torus_500: |K_0| (mean of seeds) [configs/E5.json, aggregate.py]
\newcommand{\NumWOneTorusFiveHundredSize}{1{,}500{,}000}
% w1_torus_128 / gudhi: amortized latency = total time incl. rebuilds / updates (us) [configs/E5.json, aggregate.py]
\newcommand{\NumWOneTorusOneHundredTwentyEightGudhiAmort}{24{,}400}
% w1_torus_128 / gudhi: INIT build time (ms), median of seeds [configs/E5.json, aggregate.py]
\newcommand{\NumWOneTorusOneHundredTwentyEightGudhiInitMs}{52.0}
% w1_torus_128 / gudhi: max latency (us) [configs/E5.json, aggregate.py]
\newcommand{\NumWOneTorusOneHundredTwentyEightGudhiMax}{45{,}300}
% w1_torus_128 / gudhi: median per-update latency (us), 3 seeds pooled [configs/E5.json, aggregate.py]
\newcommand{\NumWOneTorusOneHundredTwentyEightGudhiMed}{23{,}800}
% w1_torus_128 / gudhi: p95 latency (us) [configs/E5.json, aggregate.py]
\newcommand{\NumWOneTorusOneHundredTwentyEightGudhiPNinetyFive}{31{,}700}
% w1_torus_128 / gudhi: p99 latency (us) [configs/E5.json, aggregate.py]
\newcommand{\NumWOneTorusOneHundredTwentyEightGudhiPNinetyNine}{38{,}500}
% w1_torus_128 / gudhi: 25th percentile latency (us) [configs/E5.json, aggregate.py]
\newcommand{\NumWOneTorusOneHundredTwentyEightGudhiQOne}{22{,}000}
% w1_torus_128 / gudhi: 75th percentile latency (us) [configs/E5.json, aggregate.py]
\newcommand{\NumWOneTorusOneHundredTwentyEightGudhiQThree}{25{,}800}
% w1_torus_128 / gudhi: updates timed per run (T or prefix T') [configs/E5.json, aggregate.py]
\newcommand{\NumWOneTorusOneHundredTwentyEightGudhiT}{2{,}000}
% w1_torus_64 / gudhi: amortized latency = total time incl. rebuilds / updates (us) [configs/E5.json, aggregate.py]
\newcommand{\NumWOneTorusSixtyFourGudhiAmort}{4{,}760}
% w1_torus_64 / gudhi: INIT build time (ms), median of seeds [configs/E5.json, aggregate.py]
\newcommand{\NumWOneTorusSixtyFourGudhiInitMs}{11.9}
% w1_torus_64 / gudhi: max latency (us) [configs/E5.json, aggregate.py]
\newcommand{\NumWOneTorusSixtyFourGudhiMax}{11{,}800}
% w1_torus_64 / gudhi: median per-update latency (us), 3 seeds pooled [configs/E5.json, aggregate.py]
\newcommand{\NumWOneTorusSixtyFourGudhiMed}{4{,}660}
% w1_torus_64 / gudhi: p95 latency (us) [configs/E5.json, aggregate.py]
\newcommand{\NumWOneTorusSixtyFourGudhiPNinetyFive}{6{,}270}
% w1_torus_64 / gudhi: p99 latency (us) [configs/E5.json, aggregate.py]
\newcommand{\NumWOneTorusSixtyFourGudhiPNinetyNine}{7{,}840}
% w1_torus_64 / gudhi: 25th percentile latency (us) [configs/E5.json, aggregate.py]
\newcommand{\NumWOneTorusSixtyFourGudhiQOne}{4{,}340}
% w1_torus_64 / gudhi: 75th percentile latency (us) [configs/E5.json, aggregate.py]
\newcommand{\NumWOneTorusSixtyFourGudhiQThree}{4{,}950}
% w1_torus_64 / gudhi: updates timed per run (T or prefix T') [configs/E5.json, aggregate.py]
\newcommand{\NumWOneTorusSixtyFourGudhiT}{2{,}000}
% w1_torus_256 / full-morse: amortized latency = total time incl. rebuilds / updates (us) [configs/E5.json, aggregate.py]
\newcommand{\NumWOneTorusTwoHundredFiftySixFullMorseAmort}{16{,}300}
% w1_torus_256 / full-morse: INIT build time (ms), median of seeds [configs/E5.json, aggregate.py]
\newcommand{\NumWOneTorusTwoHundredFiftySixFullMorseInitMs}{251}
% w1_torus_256 / full-morse: max latency (us) [configs/E5.json, aggregate.py]
\newcommand{\NumWOneTorusTwoHundredFiftySixFullMorseMax}{49{,}600}
% w1_torus_256 / full-morse: median per-update latency (us), 3 seeds pooled [configs/E5.json, aggregate.py]
\newcommand{\NumWOneTorusTwoHundredFiftySixFullMorseMed}{15{,}300}
% w1_torus_256 / full-morse: p95 latency (us) [configs/E5.json, aggregate.py]
\newcommand{\NumWOneTorusTwoHundredFiftySixFullMorsePNinetyFive}{23{,}800}
% w1_torus_256 / full-morse: p99 latency (us) [configs/E5.json, aggregate.py]
\newcommand{\NumWOneTorusTwoHundredFiftySixFullMorsePNinetyNine}{38{,}300}
% w1_torus_256 / full-morse: 25th percentile latency (us) [configs/E5.json, aggregate.py]
\newcommand{\NumWOneTorusTwoHundredFiftySixFullMorseQOne}{14{,}100}
% w1_torus_256 / full-morse: 75th percentile latency (us) [configs/E5.json, aggregate.py]
\newcommand{\NumWOneTorusTwoHundredFiftySixFullMorseQThree}{16{,}800}
% w1_torus_256 / full-morse: amortized latency of gudhi (on its T or T' prefix) / this method [configs/E5.json, aggregate.py]
\newcommand{\NumWOneTorusTwoHundredFiftySixFullMorseSpeedupGudhi}{5.64}
% w1_torus_256 / full-morse: updates timed per run (T or prefix T') [configs/E5.json, aggregate.py]
\newcommand{\NumWOneTorusTwoHundredFiftySixFullMorseT}{2{,}000}
% w1_torus_256 / gudhi: amortized latency = total time incl. rebuilds / updates (us) [configs/E5.json, aggregate.py]
\newcommand{\NumWOneTorusTwoHundredFiftySixGudhiAmort}{91{,}800}
% w1_torus_256 / gudhi: INIT build time (ms), median of seeds [configs/E5.json, aggregate.py]
\newcommand{\NumWOneTorusTwoHundredFiftySixGudhiInitMs}{213}
% w1_torus_256 / gudhi: max latency (us) [configs/E5.json, aggregate.py]
\newcommand{\NumWOneTorusTwoHundredFiftySixGudhiMax}{163{,}000}
% w1_torus_256 / gudhi: median per-update latency (us), 3 seeds pooled [configs/E5.json, aggregate.py]
\newcommand{\NumWOneTorusTwoHundredFiftySixGudhiMed}{91{,}000}
% w1_torus_256 / gudhi: p95 latency (us) [configs/E5.json, aggregate.py]
\newcommand{\NumWOneTorusTwoHundredFiftySixGudhiPNinetyFive}{113{,}000}
% w1_torus_256 / gudhi: p99 latency (us) [configs/E5.json, aggregate.py]
\newcommand{\NumWOneTorusTwoHundredFiftySixGudhiPNinetyNine}{127{,}000}
% w1_torus_256 / gudhi: 25th percentile latency (us) [configs/E5.json, aggregate.py]
\newcommand{\NumWOneTorusTwoHundredFiftySixGudhiQOne}{85{,}500}
% w1_torus_256 / gudhi: 75th percentile latency (us) [configs/E5.json, aggregate.py]
\newcommand{\NumWOneTorusTwoHundredFiftySixGudhiQThree}{95{,}100}
% w1_torus_256 / gudhi: updates timed per run (T or prefix T') [configs/E5.json, aggregate.py]
\newcommand{\NumWOneTorusTwoHundredFiftySixGudhiT}{200}
% w3_N14_k2 / full-id: amortized latency = total time incl. rebuilds / updates (us) [configs/E3.json, aggregate.py]
\newcommand{\NumWThreeNFourteenFullIdAmort}{2{,}320}
% w3_N14_k2 / full-id: INIT build time (ms), median of seeds [configs/E3.json, aggregate.py]
\newcommand{\NumWThreeNFourteenFullIdInitMs}{31.9}
% w3_N14_k2 / full-id: max latency (us) [configs/E3.json, aggregate.py]
\newcommand{\NumWThreeNFourteenFullIdMax}{7{,}470}
% w3_N14_k2 / full-id: median per-update latency (us), 3 seeds pooled [configs/E3.json, aggregate.py]
\newcommand{\NumWThreeNFourteenFullIdMed}{2{,}290}
% w3_N14_k2 / full-id: p95 latency (us) [configs/E3.json, aggregate.py]
\newcommand{\NumWThreeNFourteenFullIdPNinetyFive}{3{,}040}
% w3_N14_k2 / full-id: p99 latency (us) [configs/E3.json, aggregate.py]
\newcommand{\NumWThreeNFourteenFullIdPNinetyNine}{3{,}590}
% w3_N14_k2 / full-id: 25th percentile latency (us) [configs/E3.json, aggregate.py]
\newcommand{\NumWThreeNFourteenFullIdQOne}{2{,}090}
% w3_N14_k2 / full-id: 75th percentile latency (us) [configs/E3.json, aggregate.py]
\newcommand{\NumWThreeNFourteenFullIdQThree}{2{,}510}
% w3_N14_k2 / full-id: amortized latency of gudhi (on its T or T' prefix) / this method [configs/E3.json, aggregate.py]
\newcommand{\NumWThreeNFourteenFullIdSpeedupGudhi}{12.1}
% w3_N14_k2 / full-id: updates timed per run (T or prefix T') [configs/E3.json, aggregate.py]
\newcommand{\NumWThreeNFourteenFullIdT}{2{,}000}
% w3_N14_k2 / full-morse: amortized latency = total time incl. rebuilds / updates (us) [configs/E3.json, aggregate.py]
\newcommand{\NumWThreeNFourteenFullMorseAmort}{4{,}320}
% w3_N14_k2 / full-morse: INIT build time (ms), median of seeds [configs/E3.json, aggregate.py]
\newcommand{\NumWThreeNFourteenFullMorseInitMs}{36.2}
% w3_N14_k2 / full-morse: max latency (us) [configs/E3.json, aggregate.py]
\newcommand{\NumWThreeNFourteenFullMorseMax}{10{,}400}
% w3_N14_k2 / full-morse: median per-update latency (us), 3 seeds pooled [configs/E3.json, aggregate.py]
\newcommand{\NumWThreeNFourteenFullMorseMed}{4{,}240}
% w3_N14_k2: amortized latency full-morse / full-id [configs/E3.json, aggregate.py]
\newcommand{\NumWThreeNFourteenFullMorseOverFullId}{1.86}
% w3_N14_k2 / full-morse: p95 latency (us) [configs/E3.json, aggregate.py]
\newcommand{\NumWThreeNFourteenFullMorsePNinetyFive}{5{,}720}
% w3_N14_k2 / full-morse: p99 latency (us) [configs/E3.json, aggregate.py]
\newcommand{\NumWThreeNFourteenFullMorsePNinetyNine}{7{,}000}
% w3_N14_k2 / full-morse: 25th percentile latency (us) [configs/E3.json, aggregate.py]
\newcommand{\NumWThreeNFourteenFullMorseQOne}{3{,}850}
% w3_N14_k2 / full-morse: 75th percentile latency (us) [configs/E3.json, aggregate.py]
\newcommand{\NumWThreeNFourteenFullMorseQThree}{4{,}630}
% w3_N14_k2 / full-morse: amortized latency of full-id / this method [configs/E3.json, aggregate.py]
\newcommand{\NumWThreeNFourteenFullMorseSpeedupFullId}{0.536}
% w3_N14_k2 / full-morse: median latency of full-id / this method [configs/E3.json, aggregate.py]
\newcommand{\NumWThreeNFourteenFullMorseSpeedupFullIdMed}{0.541}
% w3_N14_k2 / full-morse: amortized latency of gudhi (on its T or T' prefix) / this method [configs/E3.json, aggregate.py]
\newcommand{\NumWThreeNFourteenFullMorseSpeedupGudhi}{6.50}
% w3_N14_k2 / full-morse: updates timed per run (T or prefix T') [configs/E3.json, aggregate.py]
\newcommand{\NumWThreeNFourteenFullMorseT}{2{,}000}
% w3_N14_k2 / gudhi: amortized latency = total time incl. rebuilds / updates (us) [configs/E3.json, aggregate.py]
\newcommand{\NumWThreeNFourteenGudhiAmort}{28{,}100}
% w3_N14_k2 / gudhi: INIT build time (ms), median of seeds [configs/E3.json, aggregate.py]
\newcommand{\NumWThreeNFourteenGudhiInitMs}{64.0}
% w3_N14_k2 / gudhi: max latency (us) [configs/E3.json, aggregate.py]
\newcommand{\NumWThreeNFourteenGudhiMax}{91{,}100}
% w3_N14_k2 / gudhi: median per-update latency (us), 3 seeds pooled [configs/E3.json, aggregate.py]
\newcommand{\NumWThreeNFourteenGudhiMed}{27{,}300}
% w3_N14_k2 / gudhi: p95 latency (us) [configs/E3.json, aggregate.py]
\newcommand{\NumWThreeNFourteenGudhiPNinetyFive}{38{,}200}
% w3_N14_k2 / gudhi: p99 latency (us) [configs/E3.json, aggregate.py]
\newcommand{\NumWThreeNFourteenGudhiPNinetyNine}{46{,}900}
% w3_N14_k2 / gudhi: 25th percentile latency (us) [configs/E3.json, aggregate.py]
\newcommand{\NumWThreeNFourteenGudhiQOne}{25{,}200}
% w3_N14_k2 / gudhi: 75th percentile latency (us) [configs/E3.json, aggregate.py]
\newcommand{\NumWThreeNFourteenGudhiQThree}{29{,}500}
% w3_N14_k2 / gudhi: amortized latency of full-id / this method [configs/E3.json, aggregate.py]
\newcommand{\NumWThreeNFourteenGudhiSpeedupFullId}{0.0825}
% w3_N14_k2 / gudhi: median latency of full-id / this method [configs/E3.json, aggregate.py]
\newcommand{\NumWThreeNFourteenGudhiSpeedupFullIdMed}{0.0840}
% w3_N14_k2 / gudhi: updates timed per run (T or prefix T') [configs/E3.json, aggregate.py]
\newcommand{\NumWThreeNFourteenGudhiT}{2{,}000}
% w3_N14_k2: mean atomic edits per update [configs/E3.json, aggregate.py]
\newcommand{\NumWThreeNFourteenOps}{654}
% w3_N14_k2: |K_0| (mean of seeds) [configs/E3.json, aggregate.py]
\newcommand{\NumWThreeNFourteenSize}{74{,}957}
% w3_N20_k2 / drb-id: mean column additions per update (non-rebuild updates) [configs/E3.json, aggregate.py]
\newcommand{\NumWThreeNTwentyDrbIdAdds}{7{,}220}
% w3_N20_k2 / drb-id: amortized latency = total time incl. rebuilds / updates (us) [configs/E3.json, aggregate.py]
\newcommand{\NumWThreeNTwentyDrbIdAmort}{6{,}530}
% w3_N20_k2 / drb-id: mean footprint_sum per update (sum of |occ| over deletions) [configs/E3.json, aggregate.py]
\newcommand{\NumWThreeNTwentyDrbIdFootprint}{4{,}200}
% w3_N20_k2 / drb-id: INIT build time (ms), median of seeds [configs/E3.json, aggregate.py]
\newcommand{\NumWThreeNTwentyDrbIdInitMs}{116}
% w3_N20_k2 / drb-id: max latency (us) [configs/E3.json, aggregate.py]
\newcommand{\NumWThreeNTwentyDrbIdMax}{179{,}000}
% w3_N20_k2 / drb-id: median per-update latency (us), 2 seeds pooled [configs/E3.json, aggregate.py]
\newcommand{\NumWThreeNTwentyDrbIdMed}{3{,}460}
% w3_N20_k2 / drb-id: nnz(R)+nnz(V) after the last update (mean of seeds) [configs/E3.json, aggregate.py]
\newcommand{\NumWThreeNTwentyDrbIdNnzEnd}{4{,}680{,}898}
% w3_N20_k2 / drb-id: nnz at the end / nnz after INIT [configs/E3.json, aggregate.py]
\newcommand{\NumWThreeNTwentyDrbIdNnzGrowth}{5.76}
% w3_N20_k2 / drb-id: max nnz(R)+nnz(V) over the run (mean of seeds) [configs/E3.json, aggregate.py]
\newcommand{\NumWThreeNTwentyDrbIdNnzMax}{5{,}366{,}799}
% w3_N20_k2 / drb-id: nnz(R)+nnz(V) after INIT (mean of seeds) [configs/E3.json, aggregate.py]
\newcommand{\NumWThreeNTwentyDrbIdNnzStart}{812{,}580}
% w3_N20_k2 / drb-id: p95 latency (us) [configs/E3.json, aggregate.py]
\newcommand{\NumWThreeNTwentyDrbIdPNinetyFive}{21{,}800}
% w3_N20_k2 / drb-id: p99 latency (us) [configs/E3.json, aggregate.py]
\newcommand{\NumWThreeNTwentyDrbIdPNinetyNine}{47{,}000}
% w3_N20_k2 / drb-id: 25th percentile latency (us) [configs/E3.json, aggregate.py]
\newcommand{\NumWThreeNTwentyDrbIdQOne}{1{,}900}
% w3_N20_k2 / drb-id: 75th percentile latency (us) [configs/E3.json, aggregate.py]
\newcommand{\NumWThreeNTwentyDrbIdQThree}{7{,}030}
% w3_N20_k2 / drb-id-rb: mean column additions per update (non-rebuild updates) [configs/E3.json, aggregate.py]
\newcommand{\NumWThreeNTwentyDrbIdRbAdds}{5{,}780}
% w3_N20_k2 / drb-id-rb: amortized latency = total time incl. rebuilds / updates (us) [configs/E3.json, aggregate.py]
\newcommand{\NumWThreeNTwentyDrbIdRbAmort}{2{,}160}
% w3_N20_k2 / drb-id-rb: mean footprint_sum per update (sum of |occ| over deletions) [configs/E3.json, aggregate.py]
\newcommand{\NumWThreeNTwentyDrbIdRbFootprint}{2{,}210}
% w3_N20_k2 / drb-id-rb: INIT build time (ms), median of seeds [configs/E3.json, aggregate.py]
\newcommand{\NumWThreeNTwentyDrbIdRbInitMs}{114}
% w3_N20_k2 / drb-id-rb: max latency (us) [configs/E3.json, aggregate.py]
\newcommand{\NumWThreeNTwentyDrbIdRbMax}{46{,}600}
% w3_N20_k2 / drb-id-rb: median per-update latency (us), 2 seeds pooled [configs/E3.json, aggregate.py]
\newcommand{\NumWThreeNTwentyDrbIdRbMed}{1{,}450}
% w3_N20_k2 / drb-id-rb: nnz(R)+nnz(V) after the last update (mean of seeds) [configs/E3.json, aggregate.py]
\newcommand{\NumWThreeNTwentyDrbIdRbNnzEnd}{1{,}674{,}160}
% w3_N20_k2 / drb-id-rb: nnz at the end / nnz after INIT [configs/E3.json, aggregate.py]
\newcommand{\NumWThreeNTwentyDrbIdRbNnzGrowth}{2.06}
% w3_N20_k2 / drb-id-rb: max nnz(R)+nnz(V) over the run (mean of seeds) [configs/E3.json, aggregate.py]
\newcommand{\NumWThreeNTwentyDrbIdRbNnzMax}{1{,}750{,}786}
% w3_N20_k2 / drb-id-rb: nnz(R)+nnz(V) after INIT (mean of seeds) [configs/E3.json, aggregate.py]
\newcommand{\NumWThreeNTwentyDrbIdRbNnzStart}{812{,}580}
% w3_N20_k2 / drb-id-rb: p95 latency (us) [configs/E3.json, aggregate.py]
\newcommand{\NumWThreeNTwentyDrbIdRbPNinetyFive}{5{,}400}
% w3_N20_k2 / drb-id-rb: p99 latency (us) [configs/E3.json, aggregate.py]
\newcommand{\NumWThreeNTwentyDrbIdRbPNinetyNine}{9{,}290}
% w3_N20_k2 / drb-id-rb: 25th percentile latency (us) [configs/E3.json, aggregate.py]
\newcommand{\NumWThreeNTwentyDrbIdRbQOne}{989}
% w3_N20_k2 / drb-id-rb: 75th percentile latency (us) [configs/E3.json, aggregate.py]
\newcommand{\NumWThreeNTwentyDrbIdRbQThree}{2{,}430}
% w3_N20_k2 / drb-id-rb: rebuilds per run (mean of seeds), default policy alpha=1, period off [configs/E3.json, aggregate.py]
\newcommand{\NumWThreeNTwentyDrbIdRbRebuilds}{10}
% w3_N20_k2 / drb-id-rb: amortized latency of full-id / this method [configs/E3.json, aggregate.py]
\newcommand{\NumWThreeNTwentyDrbIdRbSpeedupFullId}{3.32}
% w3_N20_k2 / drb-id-rb: median latency of full-id / this method [configs/E3.json, aggregate.py]
\newcommand{\NumWThreeNTwentyDrbIdRbSpeedupFullIdMed}{5.03}
% w3_N20_k2 / drb-id-rb: amortized latency of gudhi (on its T or T' prefix) / this method [configs/E3.json, aggregate.py]
\newcommand{\NumWThreeNTwentyDrbIdRbSpeedupGudhi}{42.0}
% w3_N20_k2 / drb-id-rb: updates timed per run (T or prefix T') [configs/E3.json, aggregate.py]
\newcommand{\NumWThreeNTwentyDrbIdRbT}{2{,}000}
% w3_N20_k2 / drb-id: amortized latency of full-id / this method [configs/E3.json, aggregate.py]
\newcommand{\NumWThreeNTwentyDrbIdSpeedupFullId}{1.10}
% w3_N20_k2 / drb-id: median latency of full-id / this method [configs/E3.json, aggregate.py]
\newcommand{\NumWThreeNTwentyDrbIdSpeedupFullIdMed}{2.11}
% w3_N20_k2 / drb-id: amortized latency of gudhi (on its T or T' prefix) / this method [configs/E3.json, aggregate.py]
\newcommand{\NumWThreeNTwentyDrbIdSpeedupGudhi}{13.9}
% w3_N20_k2 / drb-id: updates timed per run (T or prefix T') [configs/E3.json, aggregate.py]
\newcommand{\NumWThreeNTwentyDrbIdT}{2{,}000}
% w3_N20_k2 / full-id: amortized latency = total time incl. rebuilds / updates (us) [configs/E3.json, aggregate.py]
\newcommand{\NumWThreeNTwentyFullIdAmort}{7{,}170}
% w3_N20_k2 / full-id: INIT build time (ms), median of seeds [configs/E3.json, aggregate.py]
\newcommand{\NumWThreeNTwentyFullIdInitMs}{105}
% w3_N20_k2 / full-id: max latency (us) [configs/E3.json, aggregate.py]
\newcommand{\NumWThreeNTwentyFullIdMax}{55{,}300}
% w3_N20_k2 / full-id: median per-update latency (us), 3 seeds pooled [configs/E3.json, aggregate.py]
\newcommand{\NumWThreeNTwentyFullIdMed}{7{,}300}
% w3_N20_k2 / full-id: p95 latency (us) [configs/E3.json, aggregate.py]
\newcommand{\NumWThreeNTwentyFullIdPNinetyFive}{9{,}860}
% w3_N20_k2 / full-id: p99 latency (us) [configs/E3.json, aggregate.py]
\newcommand{\NumWThreeNTwentyFullIdPNinetyNine}{12{,}500}
% w3_N20_k2 / full-id: 25th percentile latency (us) [configs/E3.json, aggregate.py]
\newcommand{\NumWThreeNTwentyFullIdQOne}{6{,}030}
% w3_N20_k2 / full-id: 75th percentile latency (us) [configs/E3.json, aggregate.py]
\newcommand{\NumWThreeNTwentyFullIdQThree}{8{,}330}
% w3_N20_k2 / full-id: amortized latency of gudhi (on its T or T' prefix) / this method [configs/E3.json, aggregate.py]
\newcommand{\NumWThreeNTwentyFullIdSpeedupGudhi}{12.6}
% w3_N20_k2 / full-id: updates timed per run (T or prefix T') [configs/E3.json, aggregate.py]
\newcommand{\NumWThreeNTwentyFullIdT}{2{,}000}
% w3_N20_k2 / full-morse: amortized latency = total time incl. rebuilds / updates (us) [configs/E3.json, aggregate.py]
\newcommand{\NumWThreeNTwentyFullMorseAmort}{15{,}400}
% w3_N20_k2 / full-morse: INIT build time (ms), median of seeds [configs/E3.json, aggregate.py]
\newcommand{\NumWThreeNTwentyFullMorseInitMs}{125}
% w3_N20_k2 / full-morse: max latency (us) [configs/E3.json, aggregate.py]
\newcommand{\NumWThreeNTwentyFullMorseMax}{50{,}200}
% w3_N20_k2 / full-morse: median per-update latency (us), 3 seeds pooled [configs/E3.json, aggregate.py]
\newcommand{\NumWThreeNTwentyFullMorseMed}{15{,}200}
% w3_N20_k2: amortized latency full-morse / full-id [configs/E3.json, aggregate.py]
\newcommand{\NumWThreeNTwentyFullMorseOverFullId}{2.15}
% w3_N20_k2 / full-morse: p95 latency (us) [configs/E3.json, aggregate.py]
\newcommand{\NumWThreeNTwentyFullMorsePNinetyFive}{20{,}600}
% w3_N20_k2 / full-morse: p99 latency (us) [configs/E3.json, aggregate.py]
\newcommand{\NumWThreeNTwentyFullMorsePNinetyNine}{29{,}000}
% w3_N20_k2 / full-morse: 25th percentile latency (us) [configs/E3.json, aggregate.py]
\newcommand{\NumWThreeNTwentyFullMorseQOne}{13{,}500}
% w3_N20_k2 / full-morse: 75th percentile latency (us) [configs/E3.json, aggregate.py]
\newcommand{\NumWThreeNTwentyFullMorseQThree}{17{,}000}
% w3_N20_k2 / full-morse: amortized latency of full-id / this method [configs/E3.json, aggregate.py]
\newcommand{\NumWThreeNTwentyFullMorseSpeedupFullId}{0.465}
% w3_N20_k2 / full-morse: median latency of full-id / this method [configs/E3.json, aggregate.py]
\newcommand{\NumWThreeNTwentyFullMorseSpeedupFullIdMed}{0.481}
% w3_N20_k2 / full-morse: amortized latency of gudhi (on its T or T' prefix) / this method [configs/E3.json, aggregate.py]
\newcommand{\NumWThreeNTwentyFullMorseSpeedupGudhi}{5.87}
% w3_N20_k2 / full-morse: updates timed per run (T or prefix T') [configs/E3.json, aggregate.py]
\newcommand{\NumWThreeNTwentyFullMorseT}{2{,}000}
% w3_N20_k2 / gudhi: amortized latency = total time incl. rebuilds / updates (us) [configs/E3.json, aggregate.py]
\newcommand{\NumWThreeNTwentyGudhiAmort}{90{,}500}
% w3_N20_k2 / gudhi: INIT build time (ms), median of seeds [configs/E3.json, aggregate.py]
\newcommand{\NumWThreeNTwentyGudhiInitMs}{237}
% w3_N20_k2 / gudhi: max latency (us) [configs/E3.json, aggregate.py]
\newcommand{\NumWThreeNTwentyGudhiMax}{286{,}000}
% w3_N20_k2 / gudhi: median per-update latency (us), 3 seeds pooled [configs/E3.json, aggregate.py]
\newcommand{\NumWThreeNTwentyGudhiMed}{90{,}100}
% w3_N20_k2 / gudhi: p95 latency (us) [configs/E3.json, aggregate.py]
\newcommand{\NumWThreeNTwentyGudhiPNinetyFive}{121{,}000}
% w3_N20_k2 / gudhi: p99 latency (us) [configs/E3.json, aggregate.py]
\newcommand{\NumWThreeNTwentyGudhiPNinetyNine}{143{,}000}
% w3_N20_k2 / gudhi: 25th percentile latency (us) [configs/E3.json, aggregate.py]
\newcommand{\NumWThreeNTwentyGudhiQOne}{79{,}200}
% w3_N20_k2 / gudhi: 75th percentile latency (us) [configs/E3.json, aggregate.py]
\newcommand{\NumWThreeNTwentyGudhiQThree}{100{,}000}
% w3_N20_k2 / gudhi: amortized latency of full-id / this method [configs/E3.json, aggregate.py]
\newcommand{\NumWThreeNTwentyGudhiSpeedupFullId}{0.0791}
% w3_N20_k2 / gudhi: median latency of full-id / this method [configs/E3.json, aggregate.py]
\newcommand{\NumWThreeNTwentyGudhiSpeedupFullIdMed}{0.0811}
% w3_N20_k2 / gudhi: updates timed per run (T or prefix T') [configs/E3.json, aggregate.py]
\newcommand{\NumWThreeNTwentyGudhiT}{2{,}000}
% w3_N20_k2 / mmhm: mean column additions per update (non-rebuild updates) [configs/E3.json, aggregate.py]
\newcommand{\NumWThreeNTwentyMmhmAdds}{4{,}730}
% w3_N20_k2 / mmhm: amortized latency = total time incl. rebuilds / updates (us) [configs/E3.json, aggregate.py]
\newcommand{\NumWThreeNTwentyMmhmAmort}{1{,}550}
% w3_N20_k2 / mmhm: mean footprint_sum per update (sum of |occ| over deletions) [configs/E3.json, aggregate.py]
\newcommand{\NumWThreeNTwentyMmhmFootprint}{2{,}120}
% w3_N20_k2 / mmhm: INIT build time (ms), median of seeds [configs/E3.json, aggregate.py]
\newcommand{\NumWThreeNTwentyMmhmInitMs}{122}
% w3_N20_k2 / mmhm: max latency (us) [configs/E3.json, aggregate.py]
\newcommand{\NumWThreeNTwentyMmhmMax}{88{,}400}
% w3_N20_k2 / mmhm: median per-update latency (us), 2 seeds pooled [configs/E3.json, aggregate.py]
\newcommand{\NumWThreeNTwentyMmhmMed}{1{,}030}
% w3_N20_k2 / mmhm: nnz(R)+nnz(V) after the last update (mean of seeds) [configs/E3.json, aggregate.py]
\newcommand{\NumWThreeNTwentyMmhmNnzEnd}{1{,}478{,}382}
% w3_N20_k2 / mmhm: nnz at the end / nnz after INIT [configs/E3.json, aggregate.py]
\newcommand{\NumWThreeNTwentyMmhmNnzGrowth}{1.82}
% w3_N20_k2 / mmhm: max nnz(R)+nnz(V) over the run (mean of seeds) [configs/E3.json, aggregate.py]
\newcommand{\NumWThreeNTwentyMmhmNnzMax}{1{,}619{,}300}
% w3_N20_k2 / mmhm: nnz(R)+nnz(V) after INIT (mean of seeds) [configs/E3.json, aggregate.py]
\newcommand{\NumWThreeNTwentyMmhmNnzStart}{812{,}580}
% w3_N20_k2 / mmhm-norb: mean column additions per update (non-rebuild updates) [configs/E3.json, aggregate.py]
\newcommand{\NumWThreeNTwentyMmhmNorbAdds}{5{,}440}
% w3_N20_k2 / mmhm-norb: amortized latency = total time incl. rebuilds / updates (us) [configs/E3.json, aggregate.py]
\newcommand{\NumWThreeNTwentyMmhmNorbAmort}{3{,}810}
% w3_N20_k2 / mmhm-norb: mean footprint_sum per update (sum of |occ| over deletions) [configs/E3.json, aggregate.py]
\newcommand{\NumWThreeNTwentyMmhmNorbFootprint}{3{,}010}
% w3_N20_k2 / mmhm-norb: INIT build time (ms), median of seeds [configs/E3.json, aggregate.py]
\newcommand{\NumWThreeNTwentyMmhmNorbInitMs}{132}
% w3_N20_k2 / mmhm-norb: max latency (us) [configs/E3.json, aggregate.py]
\newcommand{\NumWThreeNTwentyMmhmNorbMax}{99{,}600}
% w3_N20_k2 / mmhm-norb: median per-update latency (us), 2 seeds pooled [configs/E3.json, aggregate.py]
\newcommand{\NumWThreeNTwentyMmhmNorbMed}{2{,}190}
% w3_N20_k2 / mmhm-norb: nnz(R)+nnz(V) after the last update (mean of seeds) [configs/E3.json, aggregate.py]
\newcommand{\NumWThreeNTwentyMmhmNorbNnzEnd}{3{,}780{,}667}
% w3_N20_k2 / mmhm-norb: nnz at the end / nnz after INIT [configs/E3.json, aggregate.py]
\newcommand{\NumWThreeNTwentyMmhmNorbNnzGrowth}{4.65}
% w3_N20_k2 / mmhm-norb: max nnz(R)+nnz(V) over the run (mean of seeds) [configs/E3.json, aggregate.py]
\newcommand{\NumWThreeNTwentyMmhmNorbNnzMax}{4{,}037{,}664}
% w3_N20_k2 / mmhm-norb: nnz(R)+nnz(V) after INIT (mean of seeds) [configs/E3.json, aggregate.py]
\newcommand{\NumWThreeNTwentyMmhmNorbNnzStart}{812{,}580}
% w3_N20_k2 / mmhm-norb: p95 latency (us) [configs/E3.json, aggregate.py]
\newcommand{\NumWThreeNTwentyMmhmNorbPNinetyFive}{11{,}900}
% w3_N20_k2 / mmhm-norb: p99 latency (us) [configs/E3.json, aggregate.py]
\newcommand{\NumWThreeNTwentyMmhmNorbPNinetyNine}{24{,}800}
% w3_N20_k2 / mmhm-norb: 25th percentile latency (us) [configs/E3.json, aggregate.py]
\newcommand{\NumWThreeNTwentyMmhmNorbQOne}{1{,}290}
% w3_N20_k2 / mmhm-norb: 75th percentile latency (us) [configs/E3.json, aggregate.py]
\newcommand{\NumWThreeNTwentyMmhmNorbQThree}{4{,}070}
% w3_N20_k2 / mmhm-norb: amortized latency of full-id / this method [configs/E3.json, aggregate.py]
\newcommand{\NumWThreeNTwentyMmhmNorbSpeedupFullId}{1.88}
% w3_N20_k2 / mmhm-norb: median latency of full-id / this method [configs/E3.json, aggregate.py]
\newcommand{\NumWThreeNTwentyMmhmNorbSpeedupFullIdMed}{3.33}
% w3_N20_k2 / mmhm-norb: amortized latency of gudhi (on its T or T' prefix) / this method [configs/E3.json, aggregate.py]
\newcommand{\NumWThreeNTwentyMmhmNorbSpeedupGudhi}{23.8}
% w3_N20_k2 / mmhm-norb: updates timed per run (T or prefix T') [configs/E3.json, aggregate.py]
\newcommand{\NumWThreeNTwentyMmhmNorbT}{2{,}000}
% w3_N20_k2 / mmhm: p95 latency (us) [configs/E3.json, aggregate.py]
\newcommand{\NumWThreeNTwentyMmhmPNinetyFive}{3{,}400}
% w3_N20_k2 / mmhm: p99 latency (us) [configs/E3.json, aggregate.py]
\newcommand{\NumWThreeNTwentyMmhmPNinetyNine}{6{,}170}
% w3_N20_k2 / mmhm: 25th percentile latency (us) [configs/E3.json, aggregate.py]
\newcommand{\NumWThreeNTwentyMmhmQOne}{782}
% w3_N20_k2 / mmhm: 75th percentile latency (us) [configs/E3.json, aggregate.py]
\newcommand{\NumWThreeNTwentyMmhmQThree}{1{,}620}
% w3_N20_k2 / mmhm: rebuilds per run (mean of seeds), default policy alpha=1, period off [configs/E3.json, aggregate.py]
\newcommand{\NumWThreeNTwentyMmhmRebuilds}{6.0}
% w3_N20_k2 / mmhm: amortized latency of full-id / this method [configs/E3.json, aggregate.py]
\newcommand{\NumWThreeNTwentyMmhmSpeedupFullId}{4.62}
% w3_N20_k2 / mmhm: median latency of full-id / this method [configs/E3.json, aggregate.py]
\newcommand{\NumWThreeNTwentyMmhmSpeedupFullIdMed}{7.08}
% w3_N20_k2 / mmhm: amortized latency of gudhi (on its T or T' prefix) / this method [configs/E3.json, aggregate.py]
\newcommand{\NumWThreeNTwentyMmhmSpeedupGudhi}{58.4}
% w3_N20_k2 / mmhm: updates timed per run (T or prefix T') [configs/E3.json, aggregate.py]
\newcommand{\NumWThreeNTwentyMmhmT}{2{,}000}
% w3_N20_k2: mean atomic edits per update [configs/E3.json, aggregate.py]
\newcommand{\NumWThreeNTwentyOps}{672}
% w3_N27_k2 / drb-id: mean column additions per update (non-rebuild updates) [configs/E3.json, aggregate.py]
\newcommand{\NumWThreeNTwentySevenDrbIdAdds}{8{,}620}
% w3_N27_k2 / drb-id: amortized latency = total time incl. rebuilds / updates (us) [configs/E3.json, aggregate.py]
\newcommand{\NumWThreeNTwentySevenDrbIdAmort}{4{,}600}
% w3_N27_k2 / drb-id: mean footprint_sum per update (sum of |occ| over deletions) [configs/E3.json, aggregate.py]
\newcommand{\NumWThreeNTwentySevenDrbIdFootprint}{4{,}470}
% w3_N27_k2 / drb-id: INIT build time (ms), median of seeds [configs/E3.json, aggregate.py]
\newcommand{\NumWThreeNTwentySevenDrbIdInitMs}{370}
% w3_N27_k2 / drb-id: max latency (us) [configs/E3.json, aggregate.py]
\newcommand{\NumWThreeNTwentySevenDrbIdMax}{117{,}000}
% w3_N27_k2 / drb-id: median per-update latency (us), 3 seeds pooled [configs/E3.json, aggregate.py]
\newcommand{\NumWThreeNTwentySevenDrbIdMed}{2{,}690}
% w3_N27_k2 / drb-id: nnz(R)+nnz(V) after the last update (mean of seeds) [configs/E3.json, aggregate.py]
\newcommand{\NumWThreeNTwentySevenDrbIdNnzEnd}{9{,}410{,}451}
% w3_N27_k2 / drb-id: nnz at the end / nnz after INIT [configs/E3.json, aggregate.py]
\newcommand{\NumWThreeNTwentySevenDrbIdNnzGrowth}{4.74}
% w3_N27_k2 / drb-id: max nnz(R)+nnz(V) over the run (mean of seeds) [configs/E3.json, aggregate.py]
\newcommand{\NumWThreeNTwentySevenDrbIdNnzMax}{9{,}476{,}459}
% w3_N27_k2 / drb-id: nnz(R)+nnz(V) after INIT (mean of seeds) [configs/E3.json, aggregate.py]
\newcommand{\NumWThreeNTwentySevenDrbIdNnzStart}{1{,}986{,}039}
% w3_N27_k2 / drb-id: p95 latency (us) [configs/E3.json, aggregate.py]
\newcommand{\NumWThreeNTwentySevenDrbIdPNinetyFive}{14{,}100}
% w3_N27_k2 / drb-id: p99 latency (us) [configs/E3.json, aggregate.py]
\newcommand{\NumWThreeNTwentySevenDrbIdPNinetyNine}{29{,}600}
% w3_N27_k2 / drb-id: 25th percentile latency (us) [configs/E3.json, aggregate.py]
\newcommand{\NumWThreeNTwentySevenDrbIdQOne}{1{,}560}
% w3_N27_k2 / drb-id: 75th percentile latency (us) [configs/E3.json, aggregate.py]
\newcommand{\NumWThreeNTwentySevenDrbIdQThree}{4{,}990}
% w3_N27_k2 / drb-id-rb: mean column additions per update (non-rebuild updates) [configs/E3.json, aggregate.py]
\newcommand{\NumWThreeNTwentySevenDrbIdRbAdds}{6{,}890}
% w3_N27_k2 / drb-id-rb: amortized latency = total time incl. rebuilds / updates (us) [configs/E3.json, aggregate.py]
\newcommand{\NumWThreeNTwentySevenDrbIdRbAmort}{2{,}500}
% w3_N27_k2 / drb-id-rb: mean footprint_sum per update (sum of |occ| over deletions) [configs/E3.json, aggregate.py]
\newcommand{\NumWThreeNTwentySevenDrbIdRbFootprint}{2{,}510}
% w3_N27_k2 / drb-id-rb: INIT build time (ms), median of seeds [configs/E3.json, aggregate.py]
\newcommand{\NumWThreeNTwentySevenDrbIdRbInitMs}{387}
% w3_N27_k2 / drb-id-rb: max latency (us) [configs/E3.json, aggregate.py]
\newcommand{\NumWThreeNTwentySevenDrbIdRbMax}{111{,}000}
% w3_N27_k2 / drb-id-rb: median per-update latency (us), 3 seeds pooled [configs/E3.json, aggregate.py]
\newcommand{\NumWThreeNTwentySevenDrbIdRbMed}{1{,}620}
% w3_N27_k2 / drb-id-rb: nnz(R)+nnz(V) after the last update (mean of seeds) [configs/E3.json, aggregate.py]
\newcommand{\NumWThreeNTwentySevenDrbIdRbNnzEnd}{3{,}674{,}303}
% w3_N27_k2 / drb-id-rb: nnz at the end / nnz after INIT [configs/E3.json, aggregate.py]
\newcommand{\NumWThreeNTwentySevenDrbIdRbNnzGrowth}{1.85}
% w3_N27_k2 / drb-id-rb: max nnz(R)+nnz(V) over the run (mean of seeds) [configs/E3.json, aggregate.py]
\newcommand{\NumWThreeNTwentySevenDrbIdRbNnzMax}{4{,}223{,}413}
% w3_N27_k2 / drb-id-rb: nnz(R)+nnz(V) after INIT (mean of seeds) [configs/E3.json, aggregate.py]
\newcommand{\NumWThreeNTwentySevenDrbIdRbNnzStart}{1{,}986{,}039}
% w3_N27_k2 / drb-id-rb: p95 latency (us) [configs/E3.json, aggregate.py]
\newcommand{\NumWThreeNTwentySevenDrbIdRbPNinetyFive}{6{,}460}
% w3_N27_k2 / drb-id-rb: p99 latency (us) [configs/E3.json, aggregate.py]
\newcommand{\NumWThreeNTwentySevenDrbIdRbPNinetyNine}{11{,}700}
% w3_N27_k2 / drb-id-rb: 25th percentile latency (us) [configs/E3.json, aggregate.py]
\newcommand{\NumWThreeNTwentySevenDrbIdRbQOne}{1{,}170}
% w3_N27_k2 / drb-id-rb: 75th percentile latency (us) [configs/E3.json, aggregate.py]
\newcommand{\NumWThreeNTwentySevenDrbIdRbQThree}{2{,}580}
% w3_N27_k2 / drb-id-rb: rebuilds per run (mean of seeds), default policy alpha=1, period off [configs/E3.json, aggregate.py]
\newcommand{\NumWThreeNTwentySevenDrbIdRbRebuilds}{4.3}
% w3_N27_k2 / drb-id-rb: amortized latency of full-id / this method [configs/E3.json, aggregate.py]
\newcommand{\NumWThreeNTwentySevenDrbIdRbSpeedupFullId}{6.85}
% w3_N27_k2 / drb-id-rb: median latency of full-id / this method [configs/E3.json, aggregate.py]
\newcommand{\NumWThreeNTwentySevenDrbIdRbSpeedupFullIdMed}{10.9}
% w3_N27_k2 / drb-id-rb: amortized latency of gudhi (on its T or T' prefix) / this method [configs/E3.json, aggregate.py]
\newcommand{\NumWThreeNTwentySevenDrbIdRbSpeedupGudhi}{67.0}
% w3_N27_k2 / drb-id-rb: updates timed per run (T or prefix T') [configs/E3.json, aggregate.py]
\newcommand{\NumWThreeNTwentySevenDrbIdRbT}{2{,}000}
% w3_N27_k2 / drb-id: amortized latency of full-id / this method [configs/E3.json, aggregate.py]
\newcommand{\NumWThreeNTwentySevenDrbIdSpeedupFullId}{3.72}
% w3_N27_k2 / drb-id: median latency of full-id / this method [configs/E3.json, aggregate.py]
\newcommand{\NumWThreeNTwentySevenDrbIdSpeedupFullIdMed}{6.54}
% w3_N27_k2 / drb-id: amortized latency of gudhi (on its T or T' prefix) / this method [configs/E3.json, aggregate.py]
\newcommand{\NumWThreeNTwentySevenDrbIdSpeedupGudhi}{36.4}
% w3_N27_k2 / drb-id: updates timed per run (T or prefix T') [configs/E3.json, aggregate.py]
\newcommand{\NumWThreeNTwentySevenDrbIdT}{2{,}000}
% w3_N27_k2 / full-id: amortized latency = total time incl. rebuilds / updates (us) [configs/E3.json, aggregate.py]
\newcommand{\NumWThreeNTwentySevenFullIdAmort}{17{,}100}
% w3_N27_k2 / full-id: INIT build time (ms), median of seeds [configs/E3.json, aggregate.py]
\newcommand{\NumWThreeNTwentySevenFullIdInitMs}{382}
% w3_N27_k2 / full-id: max latency (us) [configs/E3.json, aggregate.py]
\newcommand{\NumWThreeNTwentySevenFullIdMax}{99{,}700}
% w3_N27_k2 / full-id: median per-update latency (us), 3 seeds pooled [configs/E3.json, aggregate.py]
\newcommand{\NumWThreeNTwentySevenFullIdMed}{17{,}600}
% w3_N27_k2 / full-id: p95 latency (us) [configs/E3.json, aggregate.py]
\newcommand{\NumWThreeNTwentySevenFullIdPNinetyFive}{25{,}500}
% w3_N27_k2 / full-id: p99 latency (us) [configs/E3.json, aggregate.py]
\newcommand{\NumWThreeNTwentySevenFullIdPNinetyNine}{30{,}800}
% w3_N27_k2 / full-id: 25th percentile latency (us) [configs/E3.json, aggregate.py]
\newcommand{\NumWThreeNTwentySevenFullIdQOne}{13{,}100}
% w3_N27_k2 / full-id: 75th percentile latency (us) [configs/E3.json, aggregate.py]
\newcommand{\NumWThreeNTwentySevenFullIdQThree}{20{,}700}
% w3_N27_k2 / full-id: amortized latency of gudhi (on its T or T' prefix) / this method [configs/E3.json, aggregate.py]
\newcommand{\NumWThreeNTwentySevenFullIdSpeedupGudhi}{9.78}
% w3_N27_k2 / full-id: updates timed per run (T or prefix T') [configs/E3.json, aggregate.py]
\newcommand{\NumWThreeNTwentySevenFullIdT}{2{,}000}
% w3_N27_k2 / full-morse: amortized latency = total time incl. rebuilds / updates (us) [configs/E3.json, aggregate.py]
\newcommand{\NumWThreeNTwentySevenFullMorseAmort}{39{,}200}
% w3_N27_k2 / full-morse: INIT build time (ms), median of seeds [configs/E3.json, aggregate.py]
\newcommand{\NumWThreeNTwentySevenFullMorseInitMs}{413}
% w3_N27_k2 / full-morse: max latency (us) [configs/E3.json, aggregate.py]
\newcommand{\NumWThreeNTwentySevenFullMorseMax}{89{,}500}
% w3_N27_k2 / full-morse: median per-update latency (us), 3 seeds pooled [configs/E3.json, aggregate.py]
\newcommand{\NumWThreeNTwentySevenFullMorseMed}{38{,}600}
% w3_N27_k2: amortized latency full-morse / full-id [configs/E3.json, aggregate.py]
\newcommand{\NumWThreeNTwentySevenFullMorseOverFullId}{2.29}
% w3_N27_k2 / full-morse: p95 latency (us) [configs/E3.json, aggregate.py]
\newcommand{\NumWThreeNTwentySevenFullMorsePNinetyFive}{51{,}400}
% w3_N27_k2 / full-morse: p99 latency (us) [configs/E3.json, aggregate.py]
\newcommand{\NumWThreeNTwentySevenFullMorsePNinetyNine}{69{,}800}
% w3_N27_k2 / full-morse: 25th percentile latency (us) [configs/E3.json, aggregate.py]
\newcommand{\NumWThreeNTwentySevenFullMorseQOne}{33{,}300}
% w3_N27_k2 / full-morse: 75th percentile latency (us) [configs/E3.json, aggregate.py]
\newcommand{\NumWThreeNTwentySevenFullMorseQThree}{43{,}300}
% w3_N27_k2 / full-morse: amortized latency of full-id / this method [configs/E3.json, aggregate.py]
\newcommand{\NumWThreeNTwentySevenFullMorseSpeedupFullId}{0.437}
% w3_N27_k2 / full-morse: median latency of full-id / this method [configs/E3.json, aggregate.py]
\newcommand{\NumWThreeNTwentySevenFullMorseSpeedupFullIdMed}{0.455}
% w3_N27_k2 / full-morse: amortized latency of gudhi (on its T or T' prefix) / this method [configs/E3.json, aggregate.py]
\newcommand{\NumWThreeNTwentySevenFullMorseSpeedupGudhi}{4.28}
% w3_N27_k2 / full-morse: updates timed per run (T or prefix T') [configs/E3.json, aggregate.py]
\newcommand{\NumWThreeNTwentySevenFullMorseT}{2{,}000}
% w3_N27_k2 / gudhi: amortized latency = total time incl. rebuilds / updates (us) [configs/E3.json, aggregate.py]
\newcommand{\NumWThreeNTwentySevenGudhiAmort}{168{,}000}
% w3_N27_k2 / gudhi: INIT build time (ms), median of seeds [configs/E3.json, aggregate.py]
\newcommand{\NumWThreeNTwentySevenGudhiInitMs}{339}
% w3_N27_k2 / gudhi: max latency (us) [configs/E3.json, aggregate.py]
\newcommand{\NumWThreeNTwentySevenGudhiMax}{244{,}000}
% w3_N27_k2 / gudhi: median per-update latency (us), 3 seeds pooled [configs/E3.json, aggregate.py]
\newcommand{\NumWThreeNTwentySevenGudhiMed}{164{,}000}
% w3_N27_k2 / gudhi: p95 latency (us) [configs/E3.json, aggregate.py]
\newcommand{\NumWThreeNTwentySevenGudhiPNinetyFive}{207{,}000}
% w3_N27_k2 / gudhi: p99 latency (us) [configs/E3.json, aggregate.py]
\newcommand{\NumWThreeNTwentySevenGudhiPNinetyNine}{232{,}000}
% w3_N27_k2 / gudhi: 25th percentile latency (us) [configs/E3.json, aggregate.py]
\newcommand{\NumWThreeNTwentySevenGudhiQOne}{156{,}000}
% w3_N27_k2 / gudhi: 75th percentile latency (us) [configs/E3.json, aggregate.py]
\newcommand{\NumWThreeNTwentySevenGudhiQThree}{175{,}000}
% w3_N27_k2 / gudhi: amortized latency of full-id / this method [configs/E3.json, aggregate.py]
\newcommand{\NumWThreeNTwentySevenGudhiSpeedupFullId}{0.102}
% w3_N27_k2 / gudhi: median latency of full-id / this method [configs/E3.json, aggregate.py]
\newcommand{\NumWThreeNTwentySevenGudhiSpeedupFullIdMed}{0.107}
% w3_N27_k2 / gudhi: updates timed per run (T or prefix T') [configs/E3.json, aggregate.py]
\newcommand{\NumWThreeNTwentySevenGudhiT}{200}
% w3_N27_k2 / mmhm: mean column additions per update (non-rebuild updates) [configs/E3.json, aggregate.py]
\newcommand{\NumWThreeNTwentySevenMmhmAdds}{5{,}460}
% w3_N27_k2 / mmhm: amortized latency = total time incl. rebuilds / updates (us) [configs/E3.json, aggregate.py]
\newcommand{\NumWThreeNTwentySevenMmhmAmort}{1{,}810}
% w3_N27_k2 / mmhm: mean footprint_sum per update (sum of |occ| over deletions) [configs/E3.json, aggregate.py]
\newcommand{\NumWThreeNTwentySevenMmhmFootprint}{2{,}230}
% w3_N27_k2 / mmhm: INIT build time (ms), median of seeds [configs/E3.json, aggregate.py]
\newcommand{\NumWThreeNTwentySevenMmhmInitMs}{460}
% w3_N27_k2 / mmhm: max latency (us) [configs/E3.json, aggregate.py]
\newcommand{\NumWThreeNTwentySevenMmhmMax}{187{,}000}
% w3_N27_k2 / mmhm: median per-update latency (us), 3 seeds pooled [configs/E3.json, aggregate.py]
\newcommand{\NumWThreeNTwentySevenMmhmMed}{1{,}170}
% w3_N27_k2 / mmhm: nnz(R)+nnz(V) after the last update (mean of seeds) [configs/E3.json, aggregate.py]
\newcommand{\NumWThreeNTwentySevenMmhmNnzEnd}{3{,}173{,}529}
% w3_N27_k2 / mmhm: nnz at the end / nnz after INIT [configs/E3.json, aggregate.py]
\newcommand{\NumWThreeNTwentySevenMmhmNnzGrowth}{1.60}
% w3_N27_k2 / mmhm: max nnz(R)+nnz(V) over the run (mean of seeds) [configs/E3.json, aggregate.py]
\newcommand{\NumWThreeNTwentySevenMmhmNnzMax}{3{,}968{,}745}
% w3_N27_k2 / mmhm: nnz(R)+nnz(V) after INIT (mean of seeds) [configs/E3.json, aggregate.py]
\newcommand{\NumWThreeNTwentySevenMmhmNnzStart}{1{,}986{,}039}
% w3_N27_k2 / mmhm-norb: mean column additions per update (non-rebuild updates) [configs/E3.json, aggregate.py]
\newcommand{\NumWThreeNTwentySevenMmhmNorbAdds}{6{,}510}
% w3_N27_k2 / mmhm-norb: amortized latency = total time incl. rebuilds / updates (us) [configs/E3.json, aggregate.py]
\newcommand{\NumWThreeNTwentySevenMmhmNorbAmort}{2{,}920}
% w3_N27_k2 / mmhm-norb: mean footprint_sum per update (sum of |occ| over deletions) [configs/E3.json, aggregate.py]
\newcommand{\NumWThreeNTwentySevenMmhmNorbFootprint}{3{,}530}
% w3_N27_k2 / mmhm-norb: INIT build time (ms), median of seeds [configs/E3.json, aggregate.py]
\newcommand{\NumWThreeNTwentySevenMmhmNorbInitMs}{439}
% w3_N27_k2 / mmhm-norb: max latency (us) [configs/E3.json, aggregate.py]
\newcommand{\NumWThreeNTwentySevenMmhmNorbMax}{99{,}200}
% w3_N27_k2 / mmhm-norb: median per-update latency (us), 3 seeds pooled [configs/E3.json, aggregate.py]
\newcommand{\NumWThreeNTwentySevenMmhmNorbMed}{1{,}790}
% w3_N27_k2 / mmhm-norb: nnz(R)+nnz(V) after the last update (mean of seeds) [configs/E3.json, aggregate.py]
\newcommand{\NumWThreeNTwentySevenMmhmNorbNnzEnd}{7{,}841{,}439}
% w3_N27_k2 / mmhm-norb: nnz at the end / nnz after INIT [configs/E3.json, aggregate.py]
\newcommand{\NumWThreeNTwentySevenMmhmNorbNnzGrowth}{3.95}
% w3_N27_k2 / mmhm-norb: max nnz(R)+nnz(V) over the run (mean of seeds) [configs/E3.json, aggregate.py]
\newcommand{\NumWThreeNTwentySevenMmhmNorbNnzMax}{7{,}875{,}919}
% w3_N27_k2 / mmhm-norb: nnz(R)+nnz(V) after INIT (mean of seeds) [configs/E3.json, aggregate.py]
\newcommand{\NumWThreeNTwentySevenMmhmNorbNnzStart}{1{,}986{,}039}
% w3_N27_k2 / mmhm-norb: p95 latency (us) [configs/E3.json, aggregate.py]
\newcommand{\NumWThreeNTwentySevenMmhmNorbPNinetyFive}{8{,}760}
% w3_N27_k2 / mmhm-norb: p99 latency (us) [configs/E3.json, aggregate.py]
\newcommand{\NumWThreeNTwentySevenMmhmNorbPNinetyNine}{18{,}600}
% w3_N27_k2 / mmhm-norb: 25th percentile latency (us) [configs/E3.json, aggregate.py]
\newcommand{\NumWThreeNTwentySevenMmhmNorbQOne}{1{,}130}
% w3_N27_k2 / mmhm-norb: 75th percentile latency (us) [configs/E3.json, aggregate.py]
\newcommand{\NumWThreeNTwentySevenMmhmNorbQThree}{3{,}140}
% w3_N27_k2 / mmhm-norb: amortized latency of full-id / this method [configs/E3.json, aggregate.py]
\newcommand{\NumWThreeNTwentySevenMmhmNorbSpeedupFullId}{5.86}
% w3_N27_k2 / mmhm-norb: median latency of full-id / this method [configs/E3.json, aggregate.py]
\newcommand{\NumWThreeNTwentySevenMmhmNorbSpeedupFullIdMed}{9.80}
% w3_N27_k2 / mmhm-norb: amortized latency of gudhi (on its T or T' prefix) / this method [configs/E3.json, aggregate.py]
\newcommand{\NumWThreeNTwentySevenMmhmNorbSpeedupGudhi}{57.3}
% w3_N27_k2 / mmhm-norb: updates timed per run (T or prefix T') [configs/E3.json, aggregate.py]
\newcommand{\NumWThreeNTwentySevenMmhmNorbT}{2{,}000}
% w3_N27_k2 / mmhm: p95 latency (us) [configs/E3.json, aggregate.py]
\newcommand{\NumWThreeNTwentySevenMmhmPNinetyFive}{4{,}160}
% w3_N27_k2 / mmhm: p99 latency (us) [configs/E3.json, aggregate.py]
\newcommand{\NumWThreeNTwentySevenMmhmPNinetyNine}{7{,}400}
% w3_N27_k2 / mmhm: 25th percentile latency (us) [configs/E3.json, aggregate.py]
\newcommand{\NumWThreeNTwentySevenMmhmQOne}{886}
% w3_N27_k2 / mmhm: 75th percentile latency (us) [configs/E3.json, aggregate.py]
\newcommand{\NumWThreeNTwentySevenMmhmQThree}{1{,}760}
% w3_N27_k2 / mmhm: rebuilds per run (mean of seeds), default policy alpha=1, period off [configs/E3.json, aggregate.py]
\newcommand{\NumWThreeNTwentySevenMmhmRebuilds}{3.0}
% w3_N27_k2 / mmhm: amortized latency of full-id / this method [configs/E3.json, aggregate.py]
\newcommand{\NumWThreeNTwentySevenMmhmSpeedupFullId}{9.46}
% w3_N27_k2 / mmhm: median latency of full-id / this method [configs/E3.json, aggregate.py]
\newcommand{\NumWThreeNTwentySevenMmhmSpeedupFullIdMed}{15.1}
% w3_N27_k2 / mmhm: amortized latency of gudhi (on its T or T' prefix) / this method [configs/E3.json, aggregate.py]
\newcommand{\NumWThreeNTwentySevenMmhmSpeedupGudhi}{92.5}
% w3_N27_k2 / mmhm: updates timed per run (T or prefix T') [configs/E3.json, aggregate.py]
\newcommand{\NumWThreeNTwentySevenMmhmT}{2{,}000}
% w3_N27_k2: mean atomic edits per update [configs/E3.json, aggregate.py]
\newcommand{\NumWThreeNTwentySevenOps}{683}
% w3_N27_k2: |K_0| (mean of seeds) [configs/E3.json, aggregate.py]
\newcommand{\NumWThreeNTwentySevenSize}{525{,}043}
% w3_N20_k2: |K_0| (mean of seeds) [configs/E3.json, aggregate.py]
\newcommand{\NumWThreeNTwentySize}{215{,}321}
% w2_N15_s8 / full-morse: amortized latency = total time incl. rebuilds / updates (us) [configs/E1.json, aggregate.py]
\newcommand{\NumWTwoNFifteenSEightFullMorseAmort}{3{,}170}
% w2_N15_s8 / full-morse: INIT build time (ms), median of seeds [configs/E1.json, aggregate.py]
\newcommand{\NumWTwoNFifteenSEightFullMorseInitMs}{47.9}
% w2_N15_s8 / full-morse: max latency (us) [configs/E1.json, aggregate.py]
\newcommand{\NumWTwoNFifteenSEightFullMorseMax}{7{,}030}
% w2_N15_s8 / full-morse: median per-update latency (us), 3 seeds pooled [configs/E1.json, aggregate.py]
\newcommand{\NumWTwoNFifteenSEightFullMorseMed}{3{,}070}
% w2_N15_s8 / full-morse: p95 latency (us) [configs/E1.json, aggregate.py]
\newcommand{\NumWTwoNFifteenSEightFullMorsePNinetyFive}{4{,}190}
% w2_N15_s8 / full-morse: p99 latency (us) [configs/E1.json, aggregate.py]
\newcommand{\NumWTwoNFifteenSEightFullMorsePNinetyNine}{5{,}040}
% w2_N15_s8 / full-morse: 25th percentile latency (us) [configs/E1.json, aggregate.py]
\newcommand{\NumWTwoNFifteenSEightFullMorseQOne}{2{,}920}
% w2_N15_s8 / full-morse: 75th percentile latency (us) [configs/E1.json, aggregate.py]
\newcommand{\NumWTwoNFifteenSEightFullMorseQThree}{3{,}240}
% w2_N15_s8 / full-morse: amortized latency of gudhi (on its T or T' prefix) / this method [configs/E1.json, aggregate.py]
\newcommand{\NumWTwoNFifteenSEightFullMorseSpeedupGudhi}{9.17}
% w2_N15_s8 / full-morse: updates timed per run (T or prefix T') [configs/E1.json, aggregate.py]
\newcommand{\NumWTwoNFifteenSEightFullMorseT}{2{,}000}
% w2_N15_s8 / gudhi: amortized latency = total time incl. rebuilds / updates (us) [configs/E1.json, aggregate.py]
\newcommand{\NumWTwoNFifteenSEightGudhiAmort}{29{,}100}
% w2_N15_s8 / gudhi: INIT build time (ms), median of seeds [configs/E1.json, aggregate.py]
\newcommand{\NumWTwoNFifteenSEightGudhiInitMs}{57.3}
% w2_N15_s8 / gudhi: max latency (us) [configs/E1.json, aggregate.py]
\newcommand{\NumWTwoNFifteenSEightGudhiMax}{107{,}000}
% w2_N15_s8 / gudhi: median per-update latency (us), 3 seeds pooled [configs/E1.json, aggregate.py]
\newcommand{\NumWTwoNFifteenSEightGudhiMed}{28{,}700}
% w2_N15_s8 / gudhi: p95 latency (us) [configs/E1.json, aggregate.py]
\newcommand{\NumWTwoNFifteenSEightGudhiPNinetyFive}{37{,}500}
% w2_N15_s8 / gudhi: p99 latency (us) [configs/E1.json, aggregate.py]
\newcommand{\NumWTwoNFifteenSEightGudhiPNinetyNine}{46{,}400}
% w2_N15_s8 / gudhi: 25th percentile latency (us) [configs/E1.json, aggregate.py]
\newcommand{\NumWTwoNFifteenSEightGudhiQOne}{26{,}000}
% w2_N15_s8 / gudhi: 75th percentile latency (us) [configs/E1.json, aggregate.py]
\newcommand{\NumWTwoNFifteenSEightGudhiQThree}{31{,}200}
% w2_N15_s8 / gudhi: updates timed per run (T or prefix T') [configs/E1.json, aggregate.py]
\newcommand{\NumWTwoNFifteenSEightGudhiT}{2{,}000}
% w2_N7_s8 / full-morse: amortized latency = total time incl. rebuilds / updates (us) [configs/E1.json, aggregate.py]
\newcommand{\NumWTwoNSevenSEightFullMorseAmort}{350}
% w2_N7_s8 / full-morse: INIT build time (ms), median of seeds [configs/E1.json, aggregate.py]
\newcommand{\NumWTwoNSevenSEightFullMorseInitMs}{3.54}
% w2_N7_s8 / full-morse: max latency (us) [configs/E1.json, aggregate.py]
\newcommand{\NumWTwoNSevenSEightFullMorseMax}{1{,}850}
% w2_N7_s8 / full-morse: median per-update latency (us), 3 seeds pooled [configs/E1.json, aggregate.py]
\newcommand{\NumWTwoNSevenSEightFullMorseMed}{342}
% w2_N7_s8 / full-morse: p95 latency (us) [configs/E1.json, aggregate.py]
\newcommand{\NumWTwoNSevenSEightFullMorsePNinetyFive}{517}
% w2_N7_s8 / full-morse: p99 latency (us) [configs/E1.json, aggregate.py]
\newcommand{\NumWTwoNSevenSEightFullMorsePNinetyNine}{610}
% w2_N7_s8 / full-morse: 25th percentile latency (us) [configs/E1.json, aggregate.py]
\newcommand{\NumWTwoNSevenSEightFullMorseQOne}{299}
% w2_N7_s8 / full-morse: 75th percentile latency (us) [configs/E1.json, aggregate.py]
\newcommand{\NumWTwoNSevenSEightFullMorseQThree}{374}
% w2_N7_s8 / full-morse: amortized latency of gudhi (on its T or T' prefix) / this method [configs/E1.json, aggregate.py]
\newcommand{\NumWTwoNSevenSEightFullMorseSpeedupGudhi}{7.92}
% w2_N7_s8 / full-morse: updates timed per run (T or prefix T') [configs/E1.json, aggregate.py]
\newcommand{\NumWTwoNSevenSEightFullMorseT}{2{,}000}
% w2_N7_s8 / gudhi: amortized latency = total time incl. rebuilds / updates (us) [configs/E1.json, aggregate.py]
\newcommand{\NumWTwoNSevenSEightGudhiAmort}{2{,}770}
% w2_N7_s8 / gudhi: INIT build time (ms), median of seeds [configs/E1.json, aggregate.py]
\newcommand{\NumWTwoNSevenSEightGudhiInitMs}{6.19}
% w2_N7_s8 / gudhi: max latency (us) [configs/E1.json, aggregate.py]
\newcommand{\NumWTwoNSevenSEightGudhiMax}{14{,}500}
% w2_N7_s8 / gudhi: median per-update latency (us), 3 seeds pooled [configs/E1.json, aggregate.py]
\newcommand{\NumWTwoNSevenSEightGudhiMed}{2{,}700}
% w2_N7_s8 / gudhi: p95 latency (us) [configs/E1.json, aggregate.py]
\newcommand{\NumWTwoNSevenSEightGudhiPNinetyFive}{3{,}720}
% w2_N7_s8 / gudhi: p99 latency (us) [configs/E1.json, aggregate.py]
\newcommand{\NumWTwoNSevenSEightGudhiPNinetyNine}{4{,}550}
% w2_N7_s8 / gudhi: 25th percentile latency (us) [configs/E1.json, aggregate.py]
\newcommand{\NumWTwoNSevenSEightGudhiQOne}{2{,}550}
% w2_N7_s8 / gudhi: 75th percentile latency (us) [configs/E1.json, aggregate.py]
\newcommand{\NumWTwoNSevenSEightGudhiQThree}{2{,}870}
% w2_N7_s8 / gudhi: updates timed per run (T or prefix T') [configs/E1.json, aggregate.py]
\newcommand{\NumWTwoNSevenSEightGudhiT}{2{,}000}
% w2_N10_s8 / full-morse: amortized latency = total time incl. rebuilds / updates (us) [configs/E1.json, aggregate.py]
\newcommand{\NumWTwoNTenSEightFullMorseAmort}{944}
% w2_N10_s8 / full-morse: INIT build time (ms), median of seeds [configs/E1.json, aggregate.py]
\newcommand{\NumWTwoNTenSEightFullMorseInitMs}{10.6}
% w2_N10_s8 / full-morse: max latency (us) [configs/E1.json, aggregate.py]
\newcommand{\NumWTwoNTenSEightFullMorseMax}{5{,}520}
% w2_N10_s8 / full-morse: median per-update latency (us), 3 seeds pooled [configs/E1.json, aggregate.py]
\newcommand{\NumWTwoNTenSEightFullMorseMed}{907}
% w2_N10_s8 / full-morse: p95 latency (us) [configs/E1.json, aggregate.py]
\newcommand{\NumWTwoNTenSEightFullMorsePNinetyFive}{1{,}410}
% w2_N10_s8 / full-morse: p99 latency (us) [configs/E1.json, aggregate.py]
\newcommand{\NumWTwoNTenSEightFullMorsePNinetyNine}{1{,}760}
% w2_N10_s8 / full-morse: 25th percentile latency (us) [configs/E1.json, aggregate.py]
\newcommand{\NumWTwoNTenSEightFullMorseQOne}{819}
% w2_N10_s8 / full-morse: 75th percentile latency (us) [configs/E1.json, aggregate.py]
\newcommand{\NumWTwoNTenSEightFullMorseQThree}{973}
% w2_N10_s8 / full-morse: amortized latency of gudhi (on its T or T' prefix) / this method [configs/E1.json, aggregate.py]
\newcommand{\NumWTwoNTenSEightFullMorseSpeedupGudhi}{8.68}
% w2_N10_s8 / full-morse: updates timed per run (T or prefix T') [configs/E1.json, aggregate.py]
\newcommand{\NumWTwoNTenSEightFullMorseT}{2{,}000}
% w2_N10_s8 / gudhi: amortized latency = total time incl. rebuilds / updates (us) [configs/E1.json, aggregate.py]
\newcommand{\NumWTwoNTenSEightGudhiAmort}{8{,}190}
% w2_N10_s8 / gudhi: INIT build time (ms), median of seeds [configs/E1.json, aggregate.py]
\newcommand{\NumWTwoNTenSEightGudhiInitMs}{17.1}
% w2_N10_s8 / gudhi: max latency (us) [configs/E1.json, aggregate.py]
\newcommand{\NumWTwoNTenSEightGudhiMax}{43{,}300}
% w2_N10_s8 / gudhi: median per-update latency (us), 3 seeds pooled [configs/E1.json, aggregate.py]
\newcommand{\NumWTwoNTenSEightGudhiMed}{7{,}970}
% w2_N10_s8 / gudhi: p95 latency (us) [configs/E1.json, aggregate.py]
\newcommand{\NumWTwoNTenSEightGudhiPNinetyFive}{10{,}700}
% w2_N10_s8 / gudhi: p99 latency (us) [configs/E1.json, aggregate.py]
\newcommand{\NumWTwoNTenSEightGudhiPNinetyNine}{13{,}500}
% w2_N10_s8 / gudhi: 25th percentile latency (us) [configs/E1.json, aggregate.py]
\newcommand{\NumWTwoNTenSEightGudhiQOne}{7{,}410}
% w2_N10_s8 / gudhi: 75th percentile latency (us) [configs/E1.json, aggregate.py]
\newcommand{\NumWTwoNTenSEightGudhiQThree}{8{,}670}
% w2_N10_s8 / gudhi: updates timed per run (T or prefix T') [configs/E1.json, aggregate.py]
\newcommand{\NumWTwoNTenSEightGudhiT}{2{,}000}
% w2_N33_s8 / full-id: amortized latency = total time incl. rebuilds / updates (us) [configs/E1.json, aggregate.py]
\newcommand{\NumWTwoNThirtyThreeSEightFullIdAmort}{14{,}900}
% w2_N33_s8 / full-id: INIT build time (ms), median of seeds [configs/E1.json, aggregate.py]
\newcommand{\NumWTwoNThirtyThreeSEightFullIdInitMs}{667}
% w2_N33_s8 / full-id: max latency (us) [configs/E1.json, aggregate.py]
\newcommand{\NumWTwoNThirtyThreeSEightFullIdMax}{37{,}600}
% w2_N33_s8 / full-id: median per-update latency (us), 3 seeds pooled [configs/E1.json, aggregate.py]
\newcommand{\NumWTwoNThirtyThreeSEightFullIdMed}{14{,}600}
% w2_N33_s8 / full-id: p95 latency (us) [configs/E1.json, aggregate.py]
\newcommand{\NumWTwoNThirtyThreeSEightFullIdPNinetyFive}{18{,}300}
% w2_N33_s8 / full-id: p99 latency (us) [configs/E1.json, aggregate.py]
\newcommand{\NumWTwoNThirtyThreeSEightFullIdPNinetyNine}{21{,}600}
% w2_N33_s8 / full-id: 25th percentile latency (us) [configs/E1.json, aggregate.py]
\newcommand{\NumWTwoNThirtyThreeSEightFullIdQOne}{13{,}600}
% w2_N33_s8 / full-id: 75th percentile latency (us) [configs/E1.json, aggregate.py]
\newcommand{\NumWTwoNThirtyThreeSEightFullIdQThree}{15{,}800}
% w2_N33_s8 / full-id: amortized latency of gudhi (on its T or T' prefix) / this method [configs/E1.json, aggregate.py]
\newcommand{\NumWTwoNThirtyThreeSEightFullIdSpeedupGudhi}{20.6}
% w2_N33_s8 / full-id: updates timed per run (T or prefix T') [configs/E1.json, aggregate.py]
\newcommand{\NumWTwoNThirtyThreeSEightFullIdT}{2{,}000}
% w2_N33_s8 / full-morse: amortized latency = total time incl. rebuilds / updates (us) [configs/E1.json, aggregate.py]
\newcommand{\NumWTwoNThirtyThreeSEightFullMorseAmort}{62{,}800}
% w2_N33_s8 / full-morse: INIT build time (ms), median of seeds [configs/E1.json, aggregate.py]
\newcommand{\NumWTwoNThirtyThreeSEightFullMorseInitMs}{739}
% w2_N33_s8 / full-morse: max latency (us) [configs/E1.json, aggregate.py]
\newcommand{\NumWTwoNThirtyThreeSEightFullMorseMax}{113{,}000}
% w2_N33_s8 / full-morse: median per-update latency (us), 3 seeds pooled [configs/E1.json, aggregate.py]
\newcommand{\NumWTwoNThirtyThreeSEightFullMorseMed}{61{,}200}
% w2_N33_s8: amortized latency full-morse / full-id [configs/E1.json, aggregate.py]
\newcommand{\NumWTwoNThirtyThreeSEightFullMorseOverFullId}{4.22}
% w2_N33_s8 / full-morse: p95 latency (us) [configs/E1.json, aggregate.py]
\newcommand{\NumWTwoNThirtyThreeSEightFullMorsePNinetyFive}{80{,}400}
% w2_N33_s8 / full-morse: p99 latency (us) [configs/E1.json, aggregate.py]
\newcommand{\NumWTwoNThirtyThreeSEightFullMorsePNinetyNine}{97{,}900}
% w2_N33_s8 / full-morse: 25th percentile latency (us) [configs/E1.json, aggregate.py]
\newcommand{\NumWTwoNThirtyThreeSEightFullMorseQOne}{57{,}900}
% w2_N33_s8 / full-morse: 75th percentile latency (us) [configs/E1.json, aggregate.py]
\newcommand{\NumWTwoNThirtyThreeSEightFullMorseQThree}{64{,}400}
% w2_N33_s8 / full-morse: amortized latency of full-id / this method [configs/E1.json, aggregate.py]
\newcommand{\NumWTwoNThirtyThreeSEightFullMorseSpeedupFullId}{0.237}
% w2_N33_s8 / full-morse: median latency of full-id / this method [configs/E1.json, aggregate.py]
\newcommand{\NumWTwoNThirtyThreeSEightFullMorseSpeedupFullIdMed}{0.239}
% w2_N33_s8 / full-morse: amortized latency of gudhi (on its T or T' prefix) / this method [configs/E1.json, aggregate.py]
\newcommand{\NumWTwoNThirtyThreeSEightFullMorseSpeedupGudhi}{4.89}
% w2_N33_s8 / full-morse: updates timed per run (T or prefix T') [configs/E1.json, aggregate.py]
\newcommand{\NumWTwoNThirtyThreeSEightFullMorseT}{2{,}000}
% w2_N33_s8 / gudhi: amortized latency = total time incl. rebuilds / updates (us) [configs/E1.json, aggregate.py]
\newcommand{\NumWTwoNThirtyThreeSEightGudhiAmort}{307{,}000}
% w2_N33_s8 / gudhi: INIT build time (ms), median of seeds [configs/E1.json, aggregate.py]
\newcommand{\NumWTwoNThirtyThreeSEightGudhiInitMs}{605}
% w2_N33_s8 / gudhi: max latency (us) [configs/E1.json, aggregate.py]
\newcommand{\NumWTwoNThirtyThreeSEightGudhiMax}{414{,}000}
% w2_N33_s8 / gudhi: median per-update latency (us), 3 seeds pooled [configs/E1.json, aggregate.py]
\newcommand{\NumWTwoNThirtyThreeSEightGudhiMed}{305{,}000}
% w2_N33_s8 / gudhi: p95 latency (us) [configs/E1.json, aggregate.py]
\newcommand{\NumWTwoNThirtyThreeSEightGudhiPNinetyFive}{354{,}000}
% w2_N33_s8 / gudhi: p99 latency (us) [configs/E1.json, aggregate.py]
\newcommand{\NumWTwoNThirtyThreeSEightGudhiPNinetyNine}{386{,}000}
% w2_N33_s8 / gudhi: 25th percentile latency (us) [configs/E1.json, aggregate.py]
\newcommand{\NumWTwoNThirtyThreeSEightGudhiQOne}{291{,}000}
% w2_N33_s8 / gudhi: 75th percentile latency (us) [configs/E1.json, aggregate.py]
\newcommand{\NumWTwoNThirtyThreeSEightGudhiQThree}{320{,}000}
% w2_N33_s8 / gudhi: amortized latency of full-id / this method [configs/E1.json, aggregate.py]
\newcommand{\NumWTwoNThirtyThreeSEightGudhiSpeedupFullId}{0.0485}
% w2_N33_s8 / gudhi: median latency of full-id / this method [configs/E1.json, aggregate.py]
\newcommand{\NumWTwoNThirtyThreeSEightGudhiSpeedupFullIdMed}{0.0479}
% w2_N33_s8 / gudhi: updates timed per run (T or prefix T') [configs/E1.json, aggregate.py]
\newcommand{\NumWTwoNThirtyThreeSEightGudhiT}{200}
% w2_N33_s8: mean atomic edits per update [configs/E1.json, aggregate.py]
\newcommand{\NumWTwoNThirtyThreeSEightOps}{128}
% w2_N33_s8: |K_0| (mean of seeds) [configs/E1.json, aggregate.py]
\newcommand{\NumWTwoNThirtyThreeSEightSize}{955{,}059}
% w2_N22_s8 / full-id: amortized latency = total time incl. rebuilds / updates (us) [configs/E1.json, aggregate.py]
\newcommand{\NumWTwoNTwentyTwoSEightFullIdAmort}{3{,}700}
% w2_N22_s8 / full-id: INIT build time (ms), median of seeds [configs/E1.json, aggregate.py]
\newcommand{\NumWTwoNTwentyTwoSEightFullIdInitMs}{154}
% w2_N22_s8 / full-id: max latency (us) [configs/E1.json, aggregate.py]
\newcommand{\NumWTwoNTwentyTwoSEightFullIdMax}{14{,}500}
% w2_N22_s8 / full-id: median per-update latency (us), 3 seeds pooled [configs/E1.json, aggregate.py]
\newcommand{\NumWTwoNTwentyTwoSEightFullIdMed}{3{,}590}
% w2_N22_s8 / full-id: p95 latency (us) [configs/E1.json, aggregate.py]
\newcommand{\NumWTwoNTwentyTwoSEightFullIdPNinetyFive}{4{,}840}
% w2_N22_s8 / full-id: p99 latency (us) [configs/E1.json, aggregate.py]
\newcommand{\NumWTwoNTwentyTwoSEightFullIdPNinetyNine}{5{,}730}
% w2_N22_s8 / full-id: 25th percentile latency (us) [configs/E1.json, aggregate.py]
\newcommand{\NumWTwoNTwentyTwoSEightFullIdQOne}{3{,}390}
% w2_N22_s8 / full-id: 75th percentile latency (us) [configs/E1.json, aggregate.py]
\newcommand{\NumWTwoNTwentyTwoSEightFullIdQThree}{3{,}790}
% w2_N22_s8 / full-id: amortized latency of gudhi (on its T or T' prefix) / this method [configs/E1.json, aggregate.py]
\newcommand{\NumWTwoNTwentyTwoSEightFullIdSpeedupGudhi}{26.4}
% w2_N22_s8 / full-id: updates timed per run (T or prefix T') [configs/E1.json, aggregate.py]
\newcommand{\NumWTwoNTwentyTwoSEightFullIdT}{2{,}000}
% w2_N22_s8 / full-morse: amortized latency = total time incl. rebuilds / updates (us) [configs/E1.json, aggregate.py]
\newcommand{\NumWTwoNTwentyTwoSEightFullMorseAmort}{11{,}100}
% w2_N22_s8 / full-morse: INIT build time (ms), median of seeds [configs/E1.json, aggregate.py]
\newcommand{\NumWTwoNTwentyTwoSEightFullMorseInitMs}{192}
% w2_N22_s8 / full-morse: max latency (us) [configs/E1.json, aggregate.py]
\newcommand{\NumWTwoNTwentyTwoSEightFullMorseMax}{41{,}400}
% w2_N22_s8 / full-morse: median per-update latency (us), 3 seeds pooled [configs/E1.json, aggregate.py]
\newcommand{\NumWTwoNTwentyTwoSEightFullMorseMed}{10{,}800}
% w2_N22_s8: amortized latency full-morse / full-id [configs/E2.json, aggregate.py]
\newcommand{\NumWTwoNTwentyTwoSEightFullMorseOverFullId}{3.01}
% w2_N22_s8 / full-morse: p95 latency (us) [configs/E1.json, aggregate.py]
\newcommand{\NumWTwoNTwentyTwoSEightFullMorsePNinetyFive}{14{,}800}
% w2_N22_s8 / full-morse: p99 latency (us) [configs/E1.json, aggregate.py]
\newcommand{\NumWTwoNTwentyTwoSEightFullMorsePNinetyNine}{19{,}000}
% w2_N22_s8 / full-morse: 25th percentile latency (us) [configs/E1.json, aggregate.py]
\newcommand{\NumWTwoNTwentyTwoSEightFullMorseQOne}{9{,}980}
% w2_N22_s8 / full-morse: 75th percentile latency (us) [configs/E1.json, aggregate.py]
\newcommand{\NumWTwoNTwentyTwoSEightFullMorseQThree}{11{,}500}
% w2_N22_s8 / full-morse: amortized latency of full-id / this method [configs/E1.json, aggregate.py]
\newcommand{\NumWTwoNTwentyTwoSEightFullMorseSpeedupFullId}{0.332}
% w2_N22_s8 / full-morse: median latency of full-id / this method [configs/E1.json, aggregate.py]
\newcommand{\NumWTwoNTwentyTwoSEightFullMorseSpeedupFullIdMed}{0.332}
% w2_N22_s8 / full-morse: amortized latency of gudhi (on its T or T' prefix) / this method [configs/E1.json, aggregate.py]
\newcommand{\NumWTwoNTwentyTwoSEightFullMorseSpeedupGudhi}{8.76}
% w2_N22_s8 / full-morse: updates timed per run (T or prefix T') [configs/E1.json, aggregate.py]
\newcommand{\NumWTwoNTwentyTwoSEightFullMorseT}{2{,}000}
% w2_N22_s8 / gudhi: amortized latency = total time incl. rebuilds / updates (us) [configs/E1.json, aggregate.py]
\newcommand{\NumWTwoNTwentyTwoSEightGudhiAmort}{97{,}600}
% w2_N22_s8 / gudhi: INIT build time (ms), median of seeds [configs/E1.json, aggregate.py]
\newcommand{\NumWTwoNTwentyTwoSEightGudhiInitMs}{209}
% w2_N22_s8 / gudhi: max latency (us) [configs/E1.json, aggregate.py]
\newcommand{\NumWTwoNTwentyTwoSEightGudhiMax}{475{,}000}
% w2_N22_s8 / gudhi: median per-update latency (us), 3 seeds pooled [configs/E1.json, aggregate.py]
\newcommand{\NumWTwoNTwentyTwoSEightGudhiMed}{94{,}100}
% w2_N22_s8 / gudhi: p95 latency (us) [configs/E1.json, aggregate.py]
\newcommand{\NumWTwoNTwentyTwoSEightGudhiPNinetyFive}{128{,}000}
% w2_N22_s8 / gudhi: p99 latency (us) [configs/E1.json, aggregate.py]
\newcommand{\NumWTwoNTwentyTwoSEightGudhiPNinetyNine}{184{,}000}
% w2_N22_s8 / gudhi: 25th percentile latency (us) [configs/E1.json, aggregate.py]
\newcommand{\NumWTwoNTwentyTwoSEightGudhiQOne}{86{,}800}
% w2_N22_s8 / gudhi: 75th percentile latency (us) [configs/E1.json, aggregate.py]
\newcommand{\NumWTwoNTwentyTwoSEightGudhiQThree}{102{,}000}
% w2_N22_s8 / gudhi: amortized latency of full-id / this method [configs/E1.json, aggregate.py]
\newcommand{\NumWTwoNTwentyTwoSEightGudhiSpeedupFullId}{0.0379}
% w2_N22_s8 / gudhi: median latency of full-id / this method [configs/E1.json, aggregate.py]
\newcommand{\NumWTwoNTwentyTwoSEightGudhiSpeedupFullIdMed}{0.0382}
% w2_N22_s8 / gudhi: updates timed per run (T or prefix T') [configs/E1.json, aggregate.py]
\newcommand{\NumWTwoNTwentyTwoSEightGudhiT}{2{,}000}
% w2_N22_s8: mean atomic edits per update [configs/E1.json, aggregate.py]
\newcommand{\NumWTwoNTwentyTwoSEightOps}{128}
% w2_N22_s8: |K_0| (mean of seeds) [configs/E1.json, aggregate.py]
\newcommand{\NumWTwoNTwentyTwoSEightSize}{286{,}589}
% w2_N22_s4 / full-morse: amortized latency = total time incl. rebuilds / updates (us) [configs/E2.json, aggregate.py]
\newcommand{\NumWTwoNTwentyTwoSFourFullMorseAmort}{10{,}700}
% w2_N22_s4 / full-morse: INIT build time (ms), median of seeds [configs/E2.json, aggregate.py]
\newcommand{\NumWTwoNTwentyTwoSFourFullMorseInitMs}{199}
% w2_N22_s4 / full-morse: max latency (us) [configs/E2.json, aggregate.py]
\newcommand{\NumWTwoNTwentyTwoSFourFullMorseMax}{28{,}100}
% w2_N22_s4 / full-morse: median per-update latency (us), 3 seeds pooled [configs/E2.json, aggregate.py]
\newcommand{\NumWTwoNTwentyTwoSFourFullMorseMed}{10{,}100}
% w2_N22_s4 / full-morse: p95 latency (us) [configs/E2.json, aggregate.py]
\newcommand{\NumWTwoNTwentyTwoSFourFullMorsePNinetyFive}{14{,}200}
% w2_N22_s4 / full-morse: p99 latency (us) [configs/E2.json, aggregate.py]
\newcommand{\NumWTwoNTwentyTwoSFourFullMorsePNinetyNine}{17{,}700}
% w2_N22_s4 / full-morse: 25th percentile latency (us) [configs/E2.json, aggregate.py]
\newcommand{\NumWTwoNTwentyTwoSFourFullMorseQOne}{9{,}640}
% w2_N22_s4 / full-morse: 75th percentile latency (us) [configs/E2.json, aggregate.py]
\newcommand{\NumWTwoNTwentyTwoSFourFullMorseQThree}{11{,}100}
% w2_N22_s4 / full-morse: updates timed per run (T or prefix T') [configs/E2.json, aggregate.py]
\newcommand{\NumWTwoNTwentyTwoSFourFullMorseT}{2{,}000}
% w2_N22_s1 / full-morse: amortized latency = total time incl. rebuilds / updates (us) [configs/E2.json, aggregate.py]
\newcommand{\NumWTwoNTwentyTwoSOneFullMorseAmort}{10{,}800}
% w2_N22_s1 / full-morse: INIT build time (ms), median of seeds [configs/E2.json, aggregate.py]
\newcommand{\NumWTwoNTwentyTwoSOneFullMorseInitMs}{190}
% w2_N22_s1 / full-morse: max latency (us) [configs/E2.json, aggregate.py]
\newcommand{\NumWTwoNTwentyTwoSOneFullMorseMax}{31{,}900}
% w2_N22_s1 / full-morse: median per-update latency (us), 3 seeds pooled [configs/E2.json, aggregate.py]
\newcommand{\NumWTwoNTwentyTwoSOneFullMorseMed}{10{,}500}
% w2_N22_s1 / full-morse: p95 latency (us) [configs/E2.json, aggregate.py]
\newcommand{\NumWTwoNTwentyTwoSOneFullMorsePNinetyFive}{14{,}400}
% w2_N22_s1 / full-morse: p99 latency (us) [configs/E2.json, aggregate.py]
\newcommand{\NumWTwoNTwentyTwoSOneFullMorsePNinetyNine}{18{,}900}
% w2_N22_s1 / full-morse: 25th percentile latency (us) [configs/E2.json, aggregate.py]
\newcommand{\NumWTwoNTwentyTwoSOneFullMorseQOne}{9{,}650}
% w2_N22_s1 / full-morse: 75th percentile latency (us) [configs/E2.json, aggregate.py]
\newcommand{\NumWTwoNTwentyTwoSOneFullMorseQThree}{11{,}200}
% w2_N22_s1 / full-morse: updates timed per run (T or prefix T') [configs/E2.json, aggregate.py]
\newcommand{\NumWTwoNTwentyTwoSOneFullMorseT}{2{,}000}
% w2_N22_s16 / full-id: amortized latency = total time incl. rebuilds / updates (us) [configs/E2.json, aggregate.py]
\newcommand{\NumWTwoNTwentyTwoSSixteenFullIdAmort}{3{,}750}
% w2_N22_s16 / full-id: INIT build time (ms), median of seeds [configs/E2.json, aggregate.py]
\newcommand{\NumWTwoNTwentyTwoSSixteenFullIdInitMs}{164}
% w2_N22_s16 / full-id: max latency (us) [configs/E2.json, aggregate.py]
\newcommand{\NumWTwoNTwentyTwoSSixteenFullIdMax}{8{,}050}
% w2_N22_s16 / full-id: median per-update latency (us), 3 seeds pooled [configs/E2.json, aggregate.py]
\newcommand{\NumWTwoNTwentyTwoSSixteenFullIdMed}{3{,}690}
% w2_N22_s16 / full-id: p95 latency (us) [configs/E2.json, aggregate.py]
\newcommand{\NumWTwoNTwentyTwoSSixteenFullIdPNinetyFive}{4{,}870}
% w2_N22_s16 / full-id: p99 latency (us) [configs/E2.json, aggregate.py]
\newcommand{\NumWTwoNTwentyTwoSSixteenFullIdPNinetyNine}{6{,}030}
% w2_N22_s16 / full-id: 25th percentile latency (us) [configs/E2.json, aggregate.py]
\newcommand{\NumWTwoNTwentyTwoSSixteenFullIdQOne}{3{,}440}
% w2_N22_s16 / full-id: 75th percentile latency (us) [configs/E2.json, aggregate.py]
\newcommand{\NumWTwoNTwentyTwoSSixteenFullIdQThree}{3{,}910}
% w2_N22_s16 / full-id: updates timed per run (T or prefix T') [configs/E2.json, aggregate.py]
\newcommand{\NumWTwoNTwentyTwoSSixteenFullIdT}{2{,}000}
% w2_N22_s16 / full-morse: amortized latency = total time incl. rebuilds / updates (us) [configs/E2.json, aggregate.py]
\newcommand{\NumWTwoNTwentyTwoSSixteenFullMorseAmort}{12{,}000}
% w2_N22_s16 / full-morse: INIT build time (ms), median of seeds [configs/E2.json, aggregate.py]
\newcommand{\NumWTwoNTwentyTwoSSixteenFullMorseInitMs}{179}
% w2_N22_s16 / full-morse: max latency (us) [configs/E2.json, aggregate.py]
\newcommand{\NumWTwoNTwentyTwoSSixteenFullMorseMax}{46{,}400}
% w2_N22_s16 / full-morse: median per-update latency (us), 3 seeds pooled [configs/E2.json, aggregate.py]
\newcommand{\NumWTwoNTwentyTwoSSixteenFullMorseMed}{11{,}500}
% w2_N22_s16: amortized latency full-morse / full-id [configs/E2.json, aggregate.py]
\newcommand{\NumWTwoNTwentyTwoSSixteenFullMorseOverFullId}{3.19}
% w2_N22_s16 / full-morse: p95 latency (us) [configs/E2.json, aggregate.py]
\newcommand{\NumWTwoNTwentyTwoSSixteenFullMorsePNinetyFive}{16{,}800}
% w2_N22_s16 / full-morse: p99 latency (us) [configs/E2.json, aggregate.py]
\newcommand{\NumWTwoNTwentyTwoSSixteenFullMorsePNinetyNine}{20{,}600}
% w2_N22_s16 / full-morse: 25th percentile latency (us) [configs/E2.json, aggregate.py]
\newcommand{\NumWTwoNTwentyTwoSSixteenFullMorseQOne}{10{,}600}
% w2_N22_s16 / full-morse: 75th percentile latency (us) [configs/E2.json, aggregate.py]
\newcommand{\NumWTwoNTwentyTwoSSixteenFullMorseQThree}{12{,}400}
% w2_N22_s16 / full-morse: amortized latency of full-id / this method [configs/E2.json, aggregate.py]
\newcommand{\NumWTwoNTwentyTwoSSixteenFullMorseSpeedupFullId}{0.313}
% w2_N22_s16 / full-morse: median latency of full-id / this method [configs/E2.json, aggregate.py]
\newcommand{\NumWTwoNTwentyTwoSSixteenFullMorseSpeedupFullIdMed}{0.320}
% w2_N22_s16 / full-morse: updates timed per run (T or prefix T') [configs/E2.json, aggregate.py]
\newcommand{\NumWTwoNTwentyTwoSSixteenFullMorseT}{2{,}000}
% w2_N22_s16: mean atomic edits per update [configs/E2.json, aggregate.py]
\newcommand{\NumWTwoNTwentyTwoSSixteenOps}{256}
% w2_N22_s16: |K_0| (mean of seeds) [configs/E2.json, aggregate.py]
\newcommand{\NumWTwoNTwentyTwoSSixteenSize}{287{,}485}
% w2_N22_s64 / drb-id: mean column additions per update (non-rebuild updates) [configs/E2.json, aggregate.py]
\newcommand{\NumWTwoNTwentyTwoSSixtyFourDrbIdAdds}{3{,}130}
% w2_N22_s64 / drb-id: amortized latency = total time incl. rebuilds / updates (us) [configs/E2.json, aggregate.py]
\newcommand{\NumWTwoNTwentyTwoSSixtyFourDrbIdAmort}{895}
% w2_N22_s64 / drb-id: mean footprint_sum per update (sum of |occ| over deletions) [configs/E2.json, aggregate.py]
\newcommand{\NumWTwoNTwentyTwoSSixtyFourDrbIdFootprint}{903}
% w2_N22_s64 / drb-id: INIT build time (ms), median of seeds [configs/E2.json, aggregate.py]
\newcommand{\NumWTwoNTwentyTwoSSixtyFourDrbIdInitMs}{183}
% w2_N22_s64 / drb-id: max latency (us) [configs/E2.json, aggregate.py]
\newcommand{\NumWTwoNTwentyTwoSSixtyFourDrbIdMax}{14{,}600}
% w2_N22_s64 / drb-id: median per-update latency (us), 3 seeds pooled [configs/E2.json, aggregate.py]
\newcommand{\NumWTwoNTwentyTwoSSixtyFourDrbIdMed}{789}
% w2_N22_s64 / drb-id: nnz(R)+nnz(V) after the last update (mean of seeds) [configs/E2.json, aggregate.py]
\newcommand{\NumWTwoNTwentyTwoSSixtyFourDrbIdNnzEnd}{1{,}131{,}865}
% w2_N22_s64 / drb-id: nnz at the end / nnz after INIT [configs/E2.json, aggregate.py]
\newcommand{\NumWTwoNTwentyTwoSSixtyFourDrbIdNnzGrowth}{1.02}
% w2_N22_s64 / drb-id: max nnz(R)+nnz(V) over the run (mean of seeds) [configs/E2.json, aggregate.py]
\newcommand{\NumWTwoNTwentyTwoSSixtyFourDrbIdNnzMax}{1{,}196{,}415}
% w2_N22_s64 / drb-id: nnz(R)+nnz(V) after INIT (mean of seeds) [configs/E2.json, aggregate.py]
\newcommand{\NumWTwoNTwentyTwoSSixtyFourDrbIdNnzStart}{1{,}106{,}686}
% w2_N22_s64 / drb-id: p95 latency (us) [configs/E2.json, aggregate.py]
\newcommand{\NumWTwoNTwentyTwoSSixtyFourDrbIdPNinetyFive}{1{,}530}
% w2_N22_s64 / drb-id: p99 latency (us) [configs/E2.json, aggregate.py]
\newcommand{\NumWTwoNTwentyTwoSSixtyFourDrbIdPNinetyNine}{2{,}250}
% w2_N22_s64 / drb-id: 25th percentile latency (us) [configs/E2.json, aggregate.py]
\newcommand{\NumWTwoNTwentyTwoSSixtyFourDrbIdQOne}{659}
% w2_N22_s64 / drb-id: 75th percentile latency (us) [configs/E2.json, aggregate.py]
\newcommand{\NumWTwoNTwentyTwoSSixtyFourDrbIdQThree}{1{,}010}
% w2_N22_s64 / drb-id-rb: mean column additions per update (non-rebuild updates) [configs/E2.json, aggregate.py]
\newcommand{\NumWTwoNTwentyTwoSSixtyFourDrbIdRbAdds}{3{,}130}
% w2_N22_s64 / drb-id-rb: amortized latency = total time incl. rebuilds / updates (us) [configs/E2.json, aggregate.py]
\newcommand{\NumWTwoNTwentyTwoSSixtyFourDrbIdRbAmort}{805}
% w2_N22_s64 / drb-id-rb: mean footprint_sum per update (sum of |occ| over deletions) [configs/E2.json, aggregate.py]
\newcommand{\NumWTwoNTwentyTwoSSixtyFourDrbIdRbFootprint}{903}
% w2_N22_s64 / drb-id-rb: INIT build time (ms), median of seeds [configs/E2.json, aggregate.py]
\newcommand{\NumWTwoNTwentyTwoSSixtyFourDrbIdRbInitMs}{165}
% w2_N22_s64 / drb-id-rb: max latency (us) [configs/E2.json, aggregate.py]
\newcommand{\NumWTwoNTwentyTwoSSixtyFourDrbIdRbMax}{14{,}500}
% w2_N22_s64 / drb-id-rb: median per-update latency (us), 3 seeds pooled [configs/E2.json, aggregate.py]
\newcommand{\NumWTwoNTwentyTwoSSixtyFourDrbIdRbMed}{724}
% w2_N22_s64 / drb-id-rb: nnz(R)+nnz(V) after the last update (mean of seeds) [configs/E2.json, aggregate.py]
\newcommand{\NumWTwoNTwentyTwoSSixtyFourDrbIdRbNnzEnd}{1{,}131{,}865}
% w2_N22_s64 / drb-id-rb: nnz at the end / nnz after INIT [configs/E2.json, aggregate.py]
\newcommand{\NumWTwoNTwentyTwoSSixtyFourDrbIdRbNnzGrowth}{1.02}
% w2_N22_s64 / drb-id-rb: max nnz(R)+nnz(V) over the run (mean of seeds) [configs/E2.json, aggregate.py]
\newcommand{\NumWTwoNTwentyTwoSSixtyFourDrbIdRbNnzMax}{1{,}196{,}415}
% w2_N22_s64 / drb-id-rb: nnz(R)+nnz(V) after INIT (mean of seeds) [configs/E2.json, aggregate.py]
\newcommand{\NumWTwoNTwentyTwoSSixtyFourDrbIdRbNnzStart}{1{,}106{,}686}
% w2_N22_s64 / drb-id-rb: p95 latency (us) [configs/E2.json, aggregate.py]
\newcommand{\NumWTwoNTwentyTwoSSixtyFourDrbIdRbPNinetyFive}{1{,}330}
% w2_N22_s64 / drb-id-rb: p99 latency (us) [configs/E2.json, aggregate.py]
\newcommand{\NumWTwoNTwentyTwoSSixtyFourDrbIdRbPNinetyNine}{1{,}670}
% w2_N22_s64 / drb-id-rb: 25th percentile latency (us) [configs/E2.json, aggregate.py]
\newcommand{\NumWTwoNTwentyTwoSSixtyFourDrbIdRbQOne}{613}
% w2_N22_s64 / drb-id-rb: 75th percentile latency (us) [configs/E2.json, aggregate.py]
\newcommand{\NumWTwoNTwentyTwoSSixtyFourDrbIdRbQThree}{916}
% w2_N22_s64 / drb-id-rb: rebuilds per run (mean of seeds), default policy alpha=1, period off [configs/E2.json, aggregate.py]
\newcommand{\NumWTwoNTwentyTwoSSixtyFourDrbIdRbRebuilds}{0}
% w2_N22_s64 / drb-id-rb: amortized latency of full-id / this method [configs/E2.json, aggregate.py]
\newcommand{\NumWTwoNTwentyTwoSSixtyFourDrbIdRbSpeedupFullId}{5.85}
% w2_N22_s64 / drb-id-rb: median latency of full-id / this method [configs/E2.json, aggregate.py]
\newcommand{\NumWTwoNTwentyTwoSSixtyFourDrbIdRbSpeedupFullIdMed}{6.30}
% w2_N22_s64 / drb-id-rb: updates timed per run (T or prefix T') [configs/E2.json, aggregate.py]
\newcommand{\NumWTwoNTwentyTwoSSixtyFourDrbIdRbT}{2{,}000}
% w2_N22_s64 / drb-id: amortized latency of full-id / this method [configs/E2.json, aggregate.py]
\newcommand{\NumWTwoNTwentyTwoSSixtyFourDrbIdSpeedupFullId}{5.26}
% w2_N22_s64 / drb-id: median latency of full-id / this method [configs/E2.json, aggregate.py]
\newcommand{\NumWTwoNTwentyTwoSSixtyFourDrbIdSpeedupFullIdMed}{5.78}
% w2_N22_s64 / drb-id: updates timed per run (T or prefix T') [configs/E2.json, aggregate.py]
\newcommand{\NumWTwoNTwentyTwoSSixtyFourDrbIdT}{2{,}000}
% w2_N22_s64 / full-id: amortized latency = total time incl. rebuilds / updates (us) [configs/E2.json, aggregate.py]
\newcommand{\NumWTwoNTwentyTwoSSixtyFourFullIdAmort}{4{,}710}
% w2_N22_s64 / full-id: INIT build time (ms), median of seeds [configs/E2.json, aggregate.py]
\newcommand{\NumWTwoNTwentyTwoSSixtyFourFullIdInitMs}{199}
% w2_N22_s64 / full-id: max latency (us) [configs/E2.json, aggregate.py]
\newcommand{\NumWTwoNTwentyTwoSSixtyFourFullIdMax}{12{,}700}
% w2_N22_s64 / full-id: median per-update latency (us), 3 seeds pooled [configs/E2.json, aggregate.py]
\newcommand{\NumWTwoNTwentyTwoSSixtyFourFullIdMed}{4{,}560}
% w2_N22_s64 / full-id: p95 latency (us) [configs/E2.json, aggregate.py]
\newcommand{\NumWTwoNTwentyTwoSSixtyFourFullIdPNinetyFive}{6{,}330}
% w2_N22_s64 / full-id: p99 latency (us) [configs/E2.json, aggregate.py]
\newcommand{\NumWTwoNTwentyTwoSSixtyFourFullIdPNinetyNine}{7{,}240}
% w2_N22_s64 / full-id: 25th percentile latency (us) [configs/E2.json, aggregate.py]
\newcommand{\NumWTwoNTwentyTwoSSixtyFourFullIdQOne}{4{,}240}
% w2_N22_s64 / full-id: 75th percentile latency (us) [configs/E2.json, aggregate.py]
\newcommand{\NumWTwoNTwentyTwoSSixtyFourFullIdQThree}{4{,}940}
% w2_N22_s64 / full-id: updates timed per run (T or prefix T') [configs/E2.json, aggregate.py]
\newcommand{\NumWTwoNTwentyTwoSSixtyFourFullIdT}{2{,}000}
% w2_N22_s64 / full-morse: amortized latency = total time incl. rebuilds / updates (us) [configs/E2.json, aggregate.py]
\newcommand{\NumWTwoNTwentyTwoSSixtyFourFullMorseAmort}{14{,}800}
% w2_N22_s64 / full-morse: INIT build time (ms), median of seeds [configs/E2.json, aggregate.py]
\newcommand{\NumWTwoNTwentyTwoSSixtyFourFullMorseInitMs}{219}
% w2_N22_s64 / full-morse: max latency (us) [configs/E2.json, aggregate.py]
\newcommand{\NumWTwoNTwentyTwoSSixtyFourFullMorseMax}{79{,}400}
% w2_N22_s64 / full-morse: median per-update latency (us), 3 seeds pooled [configs/E2.json, aggregate.py]
\newcommand{\NumWTwoNTwentyTwoSSixtyFourFullMorseMed}{14{,}100}
% w2_N22_s64: amortized latency full-morse / full-id [configs/E2.json, aggregate.py]
\newcommand{\NumWTwoNTwentyTwoSSixtyFourFullMorseOverFullId}{3.15}
% w2_N22_s64 / full-morse: p95 latency (us) [configs/E2.json, aggregate.py]
\newcommand{\NumWTwoNTwentyTwoSSixtyFourFullMorsePNinetyFive}{20{,}700}
% w2_N22_s64 / full-morse: p99 latency (us) [configs/E2.json, aggregate.py]
\newcommand{\NumWTwoNTwentyTwoSSixtyFourFullMorsePNinetyNine}{28{,}600}
% w2_N22_s64 / full-morse: 25th percentile latency (us) [configs/E2.json, aggregate.py]
\newcommand{\NumWTwoNTwentyTwoSSixtyFourFullMorseQOne}{13{,}200}
% w2_N22_s64 / full-morse: 75th percentile latency (us) [configs/E2.json, aggregate.py]
\newcommand{\NumWTwoNTwentyTwoSSixtyFourFullMorseQThree}{15{,}100}
% w2_N22_s64 / full-morse: amortized latency of full-id / this method [configs/E2.json, aggregate.py]
\newcommand{\NumWTwoNTwentyTwoSSixtyFourFullMorseSpeedupFullId}{0.317}
% w2_N22_s64 / full-morse: median latency of full-id / this method [configs/E2.json, aggregate.py]
\newcommand{\NumWTwoNTwentyTwoSSixtyFourFullMorseSpeedupFullIdMed}{0.322}
% w2_N22_s64 / full-morse: updates timed per run (T or prefix T') [configs/E2.json, aggregate.py]
\newcommand{\NumWTwoNTwentyTwoSSixtyFourFullMorseT}{2{,}000}
% w2_N22_s64 / mmhm: mean column additions per update (non-rebuild updates) [configs/E2.json, aggregate.py]
\newcommand{\NumWTwoNTwentyTwoSSixtyFourMmhmAdds}{2{,}600}
% w2_N22_s64 / mmhm: amortized latency = total time incl. rebuilds / updates (us) [configs/E2.json, aggregate.py]
\newcommand{\NumWTwoNTwentyTwoSSixtyFourMmhmAmort}{981}
% w2_N22_s64 / mmhm: mean footprint_sum per update (sum of |occ| over deletions) [configs/E2.json, aggregate.py]
\newcommand{\NumWTwoNTwentyTwoSSixtyFourMmhmFootprint}{966}
% w2_N22_s64 / mmhm: INIT build time (ms), median of seeds [configs/E2.json, aggregate.py]
\newcommand{\NumWTwoNTwentyTwoSSixtyFourMmhmInitMs}{202}
% w2_N22_s64 / mmhm: max latency (us) [configs/E2.json, aggregate.py]
\newcommand{\NumWTwoNTwentyTwoSSixtyFourMmhmMax}{38{,}300}
% w2_N22_s64 / mmhm: median per-update latency (us), 3 seeds pooled [configs/E2.json, aggregate.py]
\newcommand{\NumWTwoNTwentyTwoSSixtyFourMmhmMed}{884}
% w2_N22_s64 / mmhm: nnz(R)+nnz(V) after the last update (mean of seeds) [configs/E2.json, aggregate.py]
\newcommand{\NumWTwoNTwentyTwoSSixtyFourMmhmNnzEnd}{1{,}122{,}427}
% w2_N22_s64 / mmhm: nnz at the end / nnz after INIT [configs/E2.json, aggregate.py]
\newcommand{\NumWTwoNTwentyTwoSSixtyFourMmhmNnzGrowth}{1.02}
% w2_N22_s64 / mmhm: max nnz(R)+nnz(V) over the run (mean of seeds) [configs/E2.json, aggregate.py]
\newcommand{\NumWTwoNTwentyTwoSSixtyFourMmhmNnzMax}{1{,}175{,}666}
% w2_N22_s64 / mmhm: nnz(R)+nnz(V) after INIT (mean of seeds) [configs/E2.json, aggregate.py]
\newcommand{\NumWTwoNTwentyTwoSSixtyFourMmhmNnzStart}{1{,}105{,}146}
% w2_N22_s64 / mmhm-norb: mean column additions per update (non-rebuild updates) [configs/E2.json, aggregate.py]
\newcommand{\NumWTwoNTwentyTwoSSixtyFourMmhmNorbAdds}{2{,}600}
% w2_N22_s64 / mmhm-norb: amortized latency = total time incl. rebuilds / updates (us) [configs/E2.json, aggregate.py]
\newcommand{\NumWTwoNTwentyTwoSSixtyFourMmhmNorbAmort}{862}
% w2_N22_s64 / mmhm-norb: mean footprint_sum per update (sum of |occ| over deletions) [configs/E2.json, aggregate.py]
\newcommand{\NumWTwoNTwentyTwoSSixtyFourMmhmNorbFootprint}{966}
% w2_N22_s64 / mmhm-norb: INIT build time (ms), median of seeds [configs/E2.json, aggregate.py]
\newcommand{\NumWTwoNTwentyTwoSSixtyFourMmhmNorbInitMs}{204}
% w2_N22_s64 / mmhm-norb: max latency (us) [configs/E2.json, aggregate.py]
\newcommand{\NumWTwoNTwentyTwoSSixtyFourMmhmNorbMax}{13{,}100}
% w2_N22_s64 / mmhm-norb: median per-update latency (us), 3 seeds pooled [configs/E2.json, aggregate.py]
\newcommand{\NumWTwoNTwentyTwoSSixtyFourMmhmNorbMed}{814}
% w2_N22_s64 / mmhm-norb: nnz(R)+nnz(V) after the last update (mean of seeds) [configs/E2.json, aggregate.py]
\newcommand{\NumWTwoNTwentyTwoSSixtyFourMmhmNorbNnzEnd}{1{,}122{,}427}
% w2_N22_s64 / mmhm-norb: nnz at the end / nnz after INIT [configs/E2.json, aggregate.py]
\newcommand{\NumWTwoNTwentyTwoSSixtyFourMmhmNorbNnzGrowth}{1.02}
% w2_N22_s64 / mmhm-norb: max nnz(R)+nnz(V) over the run (mean of seeds) [configs/E2.json, aggregate.py]
\newcommand{\NumWTwoNTwentyTwoSSixtyFourMmhmNorbNnzMax}{1{,}175{,}666}
% w2_N22_s64 / mmhm-norb: nnz(R)+nnz(V) after INIT (mean of seeds) [configs/E2.json, aggregate.py]
\newcommand{\NumWTwoNTwentyTwoSSixtyFourMmhmNorbNnzStart}{1{,}105{,}146}
% w2_N22_s64 / mmhm-norb: p95 latency (us) [configs/E2.json, aggregate.py]
\newcommand{\NumWTwoNTwentyTwoSSixtyFourMmhmNorbPNinetyFive}{1{,}270}
% w2_N22_s64 / mmhm-norb: p99 latency (us) [configs/E2.json, aggregate.py]
\newcommand{\NumWTwoNTwentyTwoSSixtyFourMmhmNorbPNinetyNine}{1{,}640}
% w2_N22_s64 / mmhm-norb: 25th percentile latency (us) [configs/E2.json, aggregate.py]
\newcommand{\NumWTwoNTwentyTwoSSixtyFourMmhmNorbQOne}{707}
% w2_N22_s64 / mmhm-norb: 75th percentile latency (us) [configs/E2.json, aggregate.py]
\newcommand{\NumWTwoNTwentyTwoSSixtyFourMmhmNorbQThree}{955}
% w2_N22_s64 / mmhm-norb: amortized latency of full-id / this method [configs/E2.json, aggregate.py]
\newcommand{\NumWTwoNTwentyTwoSSixtyFourMmhmNorbSpeedupFullId}{5.46}
% w2_N22_s64 / mmhm-norb: median latency of full-id / this method [configs/E2.json, aggregate.py]
\newcommand{\NumWTwoNTwentyTwoSSixtyFourMmhmNorbSpeedupFullIdMed}{5.60}
% w2_N22_s64 / mmhm-norb: updates timed per run (T or prefix T') [configs/E2.json, aggregate.py]
\newcommand{\NumWTwoNTwentyTwoSSixtyFourMmhmNorbT}{2{,}000}
% w2_N22_s64 / mmhm: p95 latency (us) [configs/E2.json, aggregate.py]
\newcommand{\NumWTwoNTwentyTwoSSixtyFourMmhmPNinetyFive}{1{,}400}
% w2_N22_s64 / mmhm: p99 latency (us) [configs/E2.json, aggregate.py]
\newcommand{\NumWTwoNTwentyTwoSSixtyFourMmhmPNinetyNine}{1{,}930}
% w2_N22_s64 / mmhm: 25th percentile latency (us) [configs/E2.json, aggregate.py]
\newcommand{\NumWTwoNTwentyTwoSSixtyFourMmhmQOne}{772}
% w2_N22_s64 / mmhm: 75th percentile latency (us) [configs/E2.json, aggregate.py]
\newcommand{\NumWTwoNTwentyTwoSSixtyFourMmhmQThree}{1{,}040}
% w2_N22_s64 / mmhm: rebuilds per run (mean of seeds), default policy alpha=1, period off [configs/E2.json, aggregate.py]
\newcommand{\NumWTwoNTwentyTwoSSixtyFourMmhmRebuilds}{0}
% w2_N22_s64 / mmhm: amortized latency of full-id / this method [configs/E2.json, aggregate.py]
\newcommand{\NumWTwoNTwentyTwoSSixtyFourMmhmSpeedupFullId}{4.80}
% w2_N22_s64 / mmhm: median latency of full-id / this method [configs/E2.json, aggregate.py]
\newcommand{\NumWTwoNTwentyTwoSSixtyFourMmhmSpeedupFullIdMed}{5.16}
% w2_N22_s64 / mmhm: updates timed per run (T or prefix T') [configs/E2.json, aggregate.py]
\newcommand{\NumWTwoNTwentyTwoSSixtyFourMmhmT}{2{,}000}
% w2_N22_s64: mean atomic edits per update [configs/E2.json, aggregate.py]
\newcommand{\NumWTwoNTwentyTwoSSixtyFourOps}{1{,}020}
% w2_N22_s64: |K_0| (mean of seeds) [configs/E2.json, aggregate.py]
\newcommand{\NumWTwoNTwentyTwoSSixtyFourSize}{292{,}861}
% w2_N22_s32 / full-id: amortized latency = total time incl. rebuilds / updates (us) [configs/E2.json, aggregate.py]
\newcommand{\NumWTwoNTwentyTwoSThirtyTwoFullIdAmort}{4{,}100}
% w2_N22_s32 / full-id: INIT build time (ms), median of seeds [configs/E2.json, aggregate.py]
\newcommand{\NumWTwoNTwentyTwoSThirtyTwoFullIdInitMs}{164}
% w2_N22_s32 / full-id: max latency (us) [configs/E2.json, aggregate.py]
\newcommand{\NumWTwoNTwentyTwoSThirtyTwoFullIdMax}{23{,}500}
% w2_N22_s32 / full-id: median per-update latency (us), 3 seeds pooled [configs/E2.json, aggregate.py]
\newcommand{\NumWTwoNTwentyTwoSThirtyTwoFullIdMed}{4{,}000}
% w2_N22_s32 / full-id: p95 latency (us) [configs/E2.json, aggregate.py]
\newcommand{\NumWTwoNTwentyTwoSThirtyTwoFullIdPNinetyFive}{5{,}240}
% w2_N22_s32 / full-id: p99 latency (us) [configs/E2.json, aggregate.py]
\newcommand{\NumWTwoNTwentyTwoSThirtyTwoFullIdPNinetyNine}{6{,}280}
% w2_N22_s32 / full-id: 25th percentile latency (us) [configs/E2.json, aggregate.py]
\newcommand{\NumWTwoNTwentyTwoSThirtyTwoFullIdQOne}{3{,}740}
% w2_N22_s32 / full-id: 75th percentile latency (us) [configs/E2.json, aggregate.py]
\newcommand{\NumWTwoNTwentyTwoSThirtyTwoFullIdQThree}{4{,}290}
% w2_N22_s32 / full-id: updates timed per run (T or prefix T') [configs/E2.json, aggregate.py]
\newcommand{\NumWTwoNTwentyTwoSThirtyTwoFullIdT}{2{,}000}
% w2_N22_s32 / full-morse: amortized latency = total time incl. rebuilds / updates (us) [configs/E2.json, aggregate.py]
\newcommand{\NumWTwoNTwentyTwoSThirtyTwoFullMorseAmort}{12{,}600}
% w2_N22_s32 / full-morse: INIT build time (ms), median of seeds [configs/E2.json, aggregate.py]
\newcommand{\NumWTwoNTwentyTwoSThirtyTwoFullMorseInitMs}{189}
% w2_N22_s32 / full-morse: max latency (us) [configs/E2.json, aggregate.py]
\newcommand{\NumWTwoNTwentyTwoSThirtyTwoFullMorseMax}{34{,}100}
% w2_N22_s32 / full-morse: median per-update latency (us), 3 seeds pooled [configs/E2.json, aggregate.py]
\newcommand{\NumWTwoNTwentyTwoSThirtyTwoFullMorseMed}{12{,}300}
% w2_N22_s32: amortized latency full-morse / full-id [configs/E2.json, aggregate.py]
\newcommand{\NumWTwoNTwentyTwoSThirtyTwoFullMorseOverFullId}{3.08}
% w2_N22_s32 / full-morse: p95 latency (us) [configs/E2.json, aggregate.py]
\newcommand{\NumWTwoNTwentyTwoSThirtyTwoFullMorsePNinetyFive}{17{,}000}
% w2_N22_s32 / full-morse: p99 latency (us) [configs/E2.json, aggregate.py]
\newcommand{\NumWTwoNTwentyTwoSThirtyTwoFullMorsePNinetyNine}{21{,}600}
% w2_N22_s32 / full-morse: 25th percentile latency (us) [configs/E2.json, aggregate.py]
\newcommand{\NumWTwoNTwentyTwoSThirtyTwoFullMorseQOne}{11{,}400}
% w2_N22_s32 / full-morse: 75th percentile latency (us) [configs/E2.json, aggregate.py]
\newcommand{\NumWTwoNTwentyTwoSThirtyTwoFullMorseQThree}{13{,}100}
% w2_N22_s32 / full-morse: amortized latency of full-id / this method [configs/E2.json, aggregate.py]
\newcommand{\NumWTwoNTwentyTwoSThirtyTwoFullMorseSpeedupFullId}{0.324}
% w2_N22_s32 / full-morse: median latency of full-id / this method [configs/E2.json, aggregate.py]
\newcommand{\NumWTwoNTwentyTwoSThirtyTwoFullMorseSpeedupFullIdMed}{0.325}
% w2_N22_s32 / full-morse: updates timed per run (T or prefix T') [configs/E2.json, aggregate.py]
\newcommand{\NumWTwoNTwentyTwoSThirtyTwoFullMorseT}{2{,}000}
% w2_N22_s32: mean atomic edits per update [configs/E2.json, aggregate.py]
\newcommand{\NumWTwoNTwentyTwoSThirtyTwoOps}{512}
% w2_N22_s32: |K_0| (mean of seeds) [configs/E2.json, aggregate.py]
\newcommand{\NumWTwoNTwentyTwoSThirtyTwoSize}{289{,}277}
% w2_N22_s2 / full-morse: amortized latency = total time incl. rebuilds / updates (us) [configs/E2.json, aggregate.py]
\newcommand{\NumWTwoNTwentyTwoSTwoFullMorseAmort}{10{,}800}
% w2_N22_s2 / full-morse: INIT build time (ms), median of seeds [configs/E2.json, aggregate.py]
\newcommand{\NumWTwoNTwentyTwoSTwoFullMorseInitMs}{189}
% w2_N22_s2 / full-morse: max latency (us) [configs/E2.json, aggregate.py]
\newcommand{\NumWTwoNTwentyTwoSTwoFullMorseMax}{47{,}300}
% w2_N22_s2 / full-morse: median per-update latency (us), 3 seeds pooled [configs/E2.json, aggregate.py]
\newcommand{\NumWTwoNTwentyTwoSTwoFullMorseMed}{10{,}500}
% w2_N22_s2 / full-morse: p95 latency (us) [configs/E2.json, aggregate.py]
\newcommand{\NumWTwoNTwentyTwoSTwoFullMorsePNinetyFive}{14{,}300}
% w2_N22_s2 / full-morse: p99 latency (us) [configs/E2.json, aggregate.py]
\newcommand{\NumWTwoNTwentyTwoSTwoFullMorsePNinetyNine}{18{,}100}
% w2_N22_s2 / full-morse: 25th percentile latency (us) [configs/E2.json, aggregate.py]
\newcommand{\NumWTwoNTwentyTwoSTwoFullMorseQOne}{9{,}760}
% w2_N22_s2 / full-morse: 75th percentile latency (us) [configs/E2.json, aggregate.py]
\newcommand{\NumWTwoNTwentyTwoSTwoFullMorseQThree}{11{,}100}
% w2_N22_s2 / full-morse: updates timed per run (T or prefix T') [configs/E2.json, aggregate.py]
\newcommand{\NumWTwoNTwentyTwoSTwoFullMorseT}{2{,}000}
}{}
\IfFileExists{generated/numbers_main.tex}{% Generated by experiments/main_tables.py -- do not edit.
% Speedup: ratio of amortised latencies over the same updates, against the
% faster (amortised) of full-id and full-sorted on that instance; ProbNinetyNineRatio:
% p99 of the DRB variant divided by the p99 of that reference.
% Variants: drb-id, mmhm-norb, drb-id-rb, mmhm (WThreeMmhm...: mmhm only).
\newcommand{\NumMainWTwoLargestSpeedupMin}{88.6}
\newcommand{\NumMainWTwoLargestSpeedupMax}{89.9}
\newcommand{\NumMainWTwoSmallestSpeedupMin}{2.29}
\newcommand{\NumMainWTwoSmallestSpeedupMax}{2.89}
\newcommand{\NumMainWThreeMmhmSpeedupMin}{1.45}
\newcommand{\NumMainWThreeMmhmSpeedupMax}{5.99}
\newcommand{\NumMainWFourSpeedupMin}{3.25}
\newcommand{\NumMainWFourSpeedupMax}{109}
\newcommand{\NumMainWFourLargeSpeedupMin}{12.2}
\newcommand{\NumMainWFourLargeSpeedupMax}{109}
\newcommand{\NumMainWFourSmallSpeedupMin}{3.25}
\newcommand{\NumMainWFourSmallSpeedupMax}{6.09}
\newcommand{\NumMainWOneSpeedupMin}{0.592}
\newcommand{\NumMainWOneSpeedupMax}{1.98}
\newcommand{\NumMainWOneProbNinetyNineRatioMin}{2.02}
\newcommand{\NumMainWOneProbNinetyNineRatioMax}{6.81}
}{}
\DeclareRobustCommand{\NumMissing}[1]{%
  \PackageWarning{mmhm}{Number \string\Num#1\space undefined (generated/numbers.tex)}%
  \textbf{??}}
\newcommand{\NumFallback}[1]{%
  \expandafter\providecommand\csname Num#1\endcsname{\NumMissing{#1}}}
% Protocol parameters that are fixed by the experiment configs.
\NumFallback{Updates}          % T: updates per timed run
\NumFallback{Prefix}           % T': prefix length for static baselines at the largest sizes
\NumFallback{GudhiSizeMax}     % largest |K| at which the gudhi baseline is timed
\NumFallback{SizeMax}          % largest |K| in a timed run
% Correctness (E0).
\NumFallback{VerifiedStates}   % (method, run, state) triples compared with an oracle, all exact methods
\NumFallback{ExactMismatches}  % mismatches of exact methods (all oracles)
% W4 truncation effect.
\NumFallback{WFourBetaDRatioMin} % min over W4 instances of beta_D / n at INIT
\NumFallback{WFourBetaDRatioMax} % max over W4 instances of beta_D / n at INIT
\makeatother

